% concatenated from BreakingUniversality_LowerOrderTerms_CuspForms_SMALL2025_v10.tex
\documentclass[11pt, reqno]{amsart}
\usepackage{mathrsfs}
\usepackage{lmodern}
\usepackage[lmargin=1in,rmargin=1in,tmargin=1in,bmargin=1in,marginparwidth=5 in, marginparsep= .2 in]{geometry}
\usepackage[dvipsnames]{xcolor}
\usepackage[utf8]{inputenc}
\usepackage{bbold}
\usepackage[colorlinks=true,
            linkcolor=blue,
            citecolor=blue,
            urlcolor=cyan]{hyperref}
\usepackage{comment}
\usepackage{amsmath,amsfonts,blindtext, amssymb,amsbsy,amscd,amsthm,bbm,enumerate,enumitem,epsfig,esint,graphicx,float,marginnote,mathrsfs,mathtools,multicol}

\usepackage{tikz}



\newcommand{\bA}{\mathbb{A}} \newcommand{\bB}{\mathbb{B}}\newcommand{\bC}{\mathbb{C}}
\newcommand{\D}{\mathbb{D}}
\newcommand{\bD}{\mathbb{D}}\newcommand{\bE}{\mathbb{E}}\newcommand{\bF}{\mathbb{F}}\newcommand{\bG}{\mathbb{G}}\newcommand{\bH}{\mathbb{H}}\newcommand{\bI}{\mathbb{I}}\newcommand{\bJ}{\mathbb{J}}\newcommand{\bK}{\mathbb{K}}\newcommand{\bL}{\mathbb{L}}\newcommand{\bM}{\mathbb{M}}\newcommand{\bN}{\mathbb{N}}\newcommand{\bO}{\mathbb{O}}\newcommand{\bP}{\mathbb{P}}\newcommand{\bQ}{\mathbb{Q}}\newcommand{\bR}{\mathbb{R}}\newcommand{\bS}{\mathbb{S}}\newcommand{\bT}{\mathbb{T}}\newcommand{\bU}{\mathbb{U}}\newcommand{\bV}{\mathbb{V}}\newcommand{\bW}{\mathbb{W}}\newcommand{\bX}{\mathbb{X}}\newcommand{\bY}{\mathbb{Y}}\newcommand{\bZ}{\mathbb{Z}}\newcommand{\cA}{\mathcal{A}}\newcommand{\cB}{\mathcal{B}}\newcommand{\cC}{\mathcal{C}}\newcommand{\cD}{\mathcal{D}}\newcommand{\cE}{\mathcal{E}}\newcommand{\cF}{\mathcal{F}}\newcommand{\cG}{\mathcal{G}}\newcommand{\cH}{\mathcal{H}}\newcommand{\cJ}{\mathcal{J}}\newcommand{\cI}{\mathcal{I}}\newcommand{\cK}{\mathcal{K}}\newcommand{\cL}{\mathcal{L}}\newcommand{\cN}{\mathcal{N}}\newcommand{\cM}{\mathcal{M}}\newcommand{\cO}{\mathcal{O}}\newcommand{\cV}{\mathcal{V}}\newcommand{\cX}{\mathcal{X}}\newcommand{\cY}{\mathcal{Y}}\newcommand{\cZ}{\mathcal{Z}}\newcommand{\uC}{\underline{\mathbb{C}}}\newcommand{\RH}{\mathrm{RH}}\newcommand{\todo}[1]{\fbox{\(\color{Mahogany} \textbf{ Todo }\): #1}}\newcommand{\cbox}[1]{\fbox{\(\color{Mahogany} \textbf{ Comment: }\) #1}}\newcommand{\rup}[1]{\lceil #1 \rceil}\newcommand{\rdown}[1]{\lfloor #1 \rfloor}\newcommand{\Perv}{\mathrm{Perv}}\newcommand{\Hom}{\mathrm{Hom}}\newcommand{\cHom}{\mathop{\mathcal{H}\! \mathit{om}}}\newcommand{\RcHom}{\mathop{R\mathcal{H}\! \mathit{om}}}\newcommand{\RHom}{\mathrm{RHom}}\newcommand{\IC}{\mathrm{IC}}\newcommand{\cExt}{{\mathop{\mathcal{E}\! \mathit{xt}}}}\newcommand{\op}{\mathrm{op}}\newcommand{\Supp}{\mathrm{Supp}}\newcommand{\Exc}{\mathrm{Exc}}\renewcommand{\div}{\mathrm{div}}\newcommand{\fram}{\mathfrak{m}}\newcommand{\modd}{\ {\rm mod}\ } \newcommand{\RSA}{{\rm RSA}}\newcommand{\GL}{\mathrm{GL}}\newcommand{\SL}{\mathrm{SL}}

\makeatletter
\renewenvironment{proof}[1][\proofname]{\par
  \vspace{-.7\topsep}% remove the space after the theorem
  \pushQED{\qed}%
  \normalfont
  \topsep0pt \partopsep0pt % no space before
  \trivlist
  \item[\hskip\labelsep
        \itshape
    #1\@addpunct{.}]\ignorespaces
}{%
  \popQED\endtrivlist\@endpefalse
  \addvspace{6pt plus 6pt} % some space after
}
\makeatother


\newtheorem{theorem}{Theorem}[section]
\newtheorem{thm}{Theorem}
\newtheorem{prop}[theorem]{Proposition}
\newtheorem{cor}[theorem]{Corollary}
\newtheorem{lemma}[theorem]{Lemma}
\newtheorem{definition}{Definition}
\newtheorem{example}[theorem]{Example}
\newtheorem{problem}[theorem]{Problem}
\newtheorem{exercise}{Exercise}
\newtheorem*{thought}{Thought}
\newtheorem{bonus}{Bonus Exercise}
\newtheorem{corollary}{Corollary}
\newtheorem{conjecture}[theorem]{Conjecture}

\theoremstyle{definition}
\newtheorem{remark}[theorem]{Remark}

\theoremstyle{remark}
\newtheorem*{solution}{Solution}
\newtheorem*{claim}{Claim}
\newtheorem*{hint}{Hint}
%\newtheorem*{question}{Question}
%\newtheorem*{note}{Note}
% define useful symbols
\newcommand{\nc}{\newcommand}
\nc{\dmo}{\DeclareMathOperator}


% ~~~~~~~~~` FOR TOC ~~
\setcounter{tocdepth}{3} % to get subsubsections in toc

\let\oldtocsection=\tocsection

\let\oldtocsubsection=\tocsubsection

\let\oldtocsubsubsection=\tocsubsubsection

\renewcommand{\tocsection}[2]{\hspace{0em}\oldtocsection{#1}{#2}}
\renewcommand{\tocsubsection}[2]{\hspace{2em}\oldtocsubsection{#1}{#2}}
\renewcommand{\tocsubsubsection}[2]{\hspace{4.5em}\oldtocsubsubsection{#1}{#2}}
\usepackage{xpatch}
\makeatletter   
\xpatchcmd{\@tocline}
{\hfil\hbox to\@pnumwidth{\@tocpagenum{#7}}\par}
{\ifnum#1<0\hfill\else\dotfill\fi\hbox to\@pnumwidth{\@tocpagenum{#7}}\par}
{}{}
\makeatother

\makeatletter
\renewcommand*\@maketitle{%
  \normalfont\normalsize
  \@adminfootnotes
  \@mkboth{\@nx\shortauthors}{\@nx\shorttitle}%
  \global\topskip42\p@\relax % 5.5pc   "   "   "     "     "
  \@settitle
  \ifx\@empty\authors \else \@setauthors \fi
  \ifx\@empty\@date \else {\vskip 1em \vtop{\centering\large\@date\@@par}}\fi
  \ifx\@empty\@dedicatory
  \else
    \baselineskip18\p@
    \vtop{\centering{\footnotesize\itshape\@dedicatory\@@par}%
      \global\dimen@i\prevdepth}\prevdepth\dimen@i
  \fi
  \@setabstract
  \normalsize
  \if@titlepage
    \newpage
  \else
    \dimen@34\p@ \advance\dimen@-\baselineskip
    \vskip\dimen@\relax
  \fi
}
\makeatother

%~~~~~~~~~~~~~~~~~

\makeatletter
\renewcommand*\@maketitle{%
  \normalfont\normalsize
  \@adminfootnotes
  \@mkboth{\@nx\shortauthors}{\@nx\shorttitle}%
  \global\topskip42\p@\relax % 5.5pc   "   "   "     "     "
  \@settitle
  \ifx\@empty\authors \else \@setauthors \fi
  \ifx\@empty\@date \else {\vskip 1em \vtop{\centering\large\@date\@@par}}\fi
  \ifx\@empty\@dedicatory
  \else
    \baselineskip18\p@
    \vtop{\centering{\footnotesize\itshape\@dedicatory\@@par}%
      \global\dimen@i\prevdepth}\prevdepth\dimen@i
  \fi
  \@setabstract
  \normalsize
  \if@titlepage
    \newpage
  \else
    \dimen@34\p@ \advance\dimen@-\baselineskip
    \vskip\dimen@\relax
  \fi
}
\makeatother



% blackboard bold
\nc{\C}{\mathbb{C}}
\nc{\F}{\mathbb{F}}
\nc{\N}{\mathbb{N}}
\nc{\Q}{\mathbb{Q}}
\nc{\R}{\mathbb{R}}
\nc{\T}{\mathbb{T}}
\nc{\Z}{\mathbb{Z}}
\newcommand{\eps}{\varepsilon}

\DeclareMathOperator{\cl}{cl}
\DeclareMathOperator{\supp}{supp}
\DeclareMathOperator{\dist}{dist}
\DeclareMathOperator{\opint}{int}
\DeclareMathOperator{\osc}{osc}
\DeclareMathOperator{\opvol}{vol}
\DeclareMathOperator{\intr}{int}
\DeclareMathAlphabet{\mathcal}{OMS}{cmsy}{m}{n}
\numberwithin{equation}{section}

\newcommand{\tbf}{\textbf}
\newcommand{\mbf}{\mathbf}
\newcommand{\mbb}{\mathbb}
\newcommand{\mcal}{\mathcal}
\newcommand{\mfrak}{\mathfrak}
\newcommand{\tsup}{\textsuperscript}
\newcommand{\tsub}{\textsubscript}
\renewcommand{\hat}{\widehat}
\renewcommand{\tilde}{\widetilde}
\renewcommand{\check}{\widecheck}
\renewcommand{\bar}{\overline}

\newcommand{\intd}{\,d}
\newcommand{\del}{\partial}
\newcommand{\grad}{\nabla}
\newcommand{\lap}{\Delta}


\newcommand{\abs}[1]{\left\lvert #1 \right\rvert}
\newcommand{\norm}[1]{\lVert #1 \rVert}
\newcommand{\inprod}[1]{\left\langle #1 \right\rangle}
\newcommand{\upindex}[1]{^{(#1)}}

%bib
\usepackage[style=alphabetic,backend=biber, maxbibnames=99]{biblatex}
\addbibresource{refs.bib}

%title page
\title[Breaking Universality in the Lower Order Terms in the 1-level and 2-level Density of Holomorphic Cusp Newforms]{Breaking Universality in the Lower Order Terms in the 1-level and 2-level Density of Holomorphic Cusp Newforms}

\author{Lawrence Dillon}
\address{Department of Mathematics, University of Washington, Seattle, WA 98105}
\email{\href{mailto:}{lvid@uw.edu}}

\author{Xiaoyao Huang}
\address{Department of Mathematics, University of Michigan, Ann Arbor, MI 48109}
\email{\href{mailto:}{xyrushac@umich.edu}}

\author{Say-Yeon Kwon}
\address{Department of Mathematics, Princeton University, Princeton, NJ 08544}
\email{\href{mailto:sk9017@princeton.edu}{sk9017@princeton.edu}}

\author{Meiling Laurence}
\address{Department of Mathematics, Yale University, New Haven, CT 06520}
\email{\href{mailto:}{meiling.laurence@yale.edu}}

\author{Steven J. Miller}
\address{Department of Mathematics \& Statistics, Williams College, Williamstown, MA 01267}
\email{\href{mailto:sjm1@williams.edu}{sjm1@williams.edu}}

\author{Vishal Muthuvel}
\address{Department of Mathematics, Columbia University, New York, NY 10027}
\email{\href{mailto:vm2696@columbia.edu}{vm2696@columbia.edu}}

\author{Luke Rowen}
\address{Department of Mathematics, Carleton College, Northfield, MN 55057 }
\email{\href{mailto:rowenl@carleton.edu}{rowenl@carleton.edu}}

\author{Pramana Saldin}
\address{Department of Mathematics, University of Wisconsin, Madison, WI 53706}
\email{\href{mailto:saldin@wisc.edu}{saldin@wisc.edu}}

\author{Steven Zanetti}
\address{Department of Mathematics, University of Michigan, Ann Arbor, MI 48109}
\email{\href{mailto:szanetti@umich.edu}{szanetti@umich.edu}}

\thanks{\textit{Date}: \today. \\ \indent Acknowledgments: This research was supported by the National Science Foundation under Grant No. DMS-2241623, as well as Columbia University, Princeton University, the University of Michigan, the University of Washington, Williams College, and Yale University.}



\begin{document}

%%%%%%%%%%%%%%%%

\begin{abstract} The Katz-Sarnak density conjecture states that, as the analytic conductor $R \to \infty$, the distribution of the normalized low-lying zeros (those near the central point $s = 1/2$) converges to the scaling limits of eigenvalues clustered near 1 of subgroups of $U(N)$. There is extensive evidence supporting this conjecture for many families, including the family of holomorphic cusp newforms. Interestingly, there are very few choices for the main term of the limiting behavior. In 2009, S. J. Miller computed lower-order terms for the 1-level density of families of elliptic curve $L$-functions and compared to cuspidal newforms of prime level; while the main terms agreed, the lower order terms depended on the arithmetic of the family. We extend his work by identifying family-dependent lower-order correction terms in the weighted 1-level and 2-level densities of holomorphic cusp newforms up to $O\left(1/\log^4 R\right)$ error, sharpening Miller's $O\left(1/\log^3 R\right)$ error. We consider cases where the level is prime or when the level is a product of two, not necessarily distinct, primes. We show 
that the rates at which the prime factors of the level tend to infinity lead to different lower-order terms, breaking the universality of the main behavior.
\end{abstract}


%\date{\today}
\maketitle
\tableofcontents
\pagestyle{plain}

\section{Introduction}
    The Katz-Sarnak conjecture states that as we take the limit as the analytic conductor goes to infinity for families of $L$-functions, the zeros lying near the central point agree with the scaling limits of eigenvalues near 1 of a classical compact group  (\cite{katz1999random}, \cite{katz1999zeros}, \cite{Hej}, \cite{Mon}, \cite{Od1}, \cite{Od2}, \cite{RS}). Supporting evidence for the conjecture has been gathered for various families, by studying moments of $L$-functions (\cite{CF2000}, \cite{CFKRS2005}, \cite{KeSn1}, \cite{KeSn2}, \cite{KeSn3}) as well as by investigating the $n$-level density for suitable test functions (\cite{DM1}, \cite{FI}, \cite{Gu}, \cite{HR}, \cite{HM}, \cite{Mil2}, \cite{OS}, \cite{RR1}, \cite{Ro}, \cite{Rub}, \cite{Yo2}).

Specifically, given an $L$-function $L(s,f)$ associated with a holomorphic cusp newform $f \in H^*_k(N) $ of level $N$ and weight $k$, we assume the Generalized Riemann Hypothesis (GRH).\footnote{We do not need the GRH to define any of the statistics described below. However, assuming the GRH allows us to relate the $n$-level density with the spacing between the zeros.} Thus, all non-trivial zeros of $L(s,f)$ are of the form $\rho_f=1/2+i\gamma_f$. We then enumerate zeros by their imaginary part $0 \leq \gamma_f^{(1)}\leq \gamma_f^{(2)}\leq \ldots$ with $\gamma_f^{(-j)}=-\gamma_f^{(j)}$ by symmetry. We define the $n$-level density by 
 \begin{equation}
    D_n(f;\Phi)\ :=\ \sum_{\substack{j_1,\ldots,j_n\\j_i \neq \pm j_k}}\phi_1\left( \gamma^{(j_1)}_f \frac{\log (R)}{2\pi}\right)\cdots \phi_n\left( \gamma^{(j_n)}_f \frac{\log (R)}{2\pi}\right),
 \end{equation}
 where\footnote{In general, we do not need the test function $\Phi$ to be a product of coordinate-wise test functions.} $\Phi(x_1, \ldots, x_n)=\phi_1(x_1)\cdot \phi_2(x_2)\cdots\phi_n(x_n)$ is a Schwartz test function and $R$ is the analytic conductor of $f$. We rescale the zeros near the central point by $\log (R)/2\pi$ (in all our families of interest, $\log( R) \sim \log(N)$). The Katz-Sarnak conjecture states that for a family $\mathcal{F}=\bigcup\mathcal{F}_N$ of $L$-functions ordered by their conductors, we have
 \begin{equation}
     \lim_{N\to \infty}\frac{\sum_{f\in \mathcal{F}_N}D_n(f;\Phi)}{|\mathcal{F}_N|}\ =\ \int \cdots\int \Phi(x_1,\ldots, x_n)W_{n, G(\mathcal{F})}(x_1,\ldots,x_n)dx_1\cdots dx_n,
 \end{equation}
where $W_{G(\mathcal{F})}$ represents the limiting distribution of a similar statistic for the eigenvalues of random matrices in some classical compact group as their rank goes to infinity. While the main terms in the expansions of the $n$-level density have been shown to agree with random matrix theory, the lower order terms break this universality and provide insight in the arithmetic properties of the family and can serve to differntiate families. For example, \cite{young2005lower} analyzed lower order terms arising from families of elliptic curves, revealing that the number of low lying zeros is highly family-dependent (see also \cite{GAO_ZHAO_2023}, \cite{Fiorilli_2015}, \cite{Fourvy_2003} for investigations of the lower order terms for other families of $L$-functions). 

Of particular importance to us,  \cite{Miller2009_LOTerms_1LevelDensity} demonstrated how the lower-order terms in the 1-level density can provide insight into the arithmetic properties of the family being studied by looking specificaly at one-parameter families of elliptic curves. Moreover, in his thesis \cite{Mil2} and some subsequent papers (\cite{Miller2005}, \cite{Miller2009_LOTerms_1LevelDensity}), S. J. Miller noticed that different one parameter families of elliptic curve $L$-functions have different lower order terms for the moments of their Satake parameters, and this has implications for the distribution of zeros near the central point, explaining some of the observed excess rank. The calculations are complicated by the fact that the lower-order terms for the $n$-level densities can be greatly impacted by the small primes. This led him to develop averaging formulas to numerically approximate all the contributions. He proved that while the main terms agree, different families had different lower-order terms. He then compared that to the lower order terms for the family of all cuspidal newforms, for simplicity of prime level N tending to infinity.

We investigate how the lower-order terms of the weighted first-level and second-level density over the space of holomorphic cusp newforms, which are respectively defined by
\begin{equation}\label{eq:2ndlvl-density-def}
     D_1(H_k^*(N))\ :=\ \lim_{N\to \infty}\frac{1}{W_R(H_k^*(N))}\sum_{f \in H_k^*(N)}w_R(f)D_1(f ; \phi),
\end{equation}
\begin{equation}\label{eq:2ndlvl-density-def-2}
     D_2(H_k^*(N))\ :=\ \lim_{N\to \infty}\frac{1}{W_R(H_k^*(N))}\sum_{f \in H_k^*(N)}w_R(f)D_2(f ;  \phi_1,\phi_2),
\end{equation}where
\begin{equation}
    D_1(f, \phi) \ :=\ \displaystyle\sum_{i}\phi\left(\frac{2\pi}{\log(R) }\gamma_f^{(i)}\right),
\end{equation}
\begin{equation}
D_2(f; \phi_1,\phi_2)\ :=\ \sum_{\substack{i,j\\i \neq \pm j}}\phi_1 \left( \frac{2\pi}{\log(R)} \gamma_f^{(i)} \right)\phi_2 \left( \frac{2\pi}{\log(R)} \gamma_f^{(j)} \right),
\end{equation} 
vary depending on how the prime factors of the level $N$ of the family approach infinity. By doing this, we break the universality of the main term behavior suggested by the Katz-Sarnak Conjecture.

We consider four different scenarios in which the level approaches infinity. \begin{itemize}
    \item The first case is $N$ going to infinity through primes. 
    \item The second case is when $N=q_1q_2$ for $q_1 \neq q_2$ two primes with $q_1$ fixed and $q_2 \to \infty.$
    \item The third case is when $N=q_1q_2$ for $q_1 \neq q_2$ two primes with $q_1 \sim N^{\delta}$ and $q_2 \sim N^{1-\delta}$ where $\delta \in (0,1/2].$
    \item Lastly, we consider the case when $N=p^2$.
\end{itemize}  Here, we weight each $f\in H_k^*(N)$ by non-negative weights $w_R(f)$ and define
\begin{align}
    W_R(H_k^*(N)):= \sum_{f\in H_k^*(N)}w_R(f).
\end{align} Moreover, both of our test functions $\phi_1(x)$ and $\phi_2(y)$ are even Schwartz functions whose Fourier transform has compact support.

We write $D_2(f,\phi_1,\phi_2)$ in terms of $D_1(f,\phi_i)$ using inclusion-exclusion.
\begin{equation}\label{eq:2lvl_break_down}
    D_2(f; \phi_1, \phi_2)\ =\ D_1(f, \phi_1)D_1(f, \phi_2)-2D_1(f, \phi_1\phi_2)+\frac{1-\varepsilon_f}{2}\phi_1(0)\phi_2(0),
\end{equation} where $\epsilon_f = \pm1$ depending on whether the functional equation associated with the completed $L$-function $\Lambda(f,s)$ is even or odd. Then, we get
\small\begin{equation}
    D_2\left(H_k^*(N)\right) \ =\ \lim_{N\to \infty } \frac{1}{W_R(H_k^*(N))}\sum_{f\in H_k^*(N)} w_R(f)\left(D_1(f, \phi_1)D_1(f, \phi_2)-2D_1(f, \phi_1\phi_2)+\frac{1-\varepsilon_f}{2}\phi_1(0)\phi_2(0)\right).
\end{equation}\normalsize
Whenever there is no ambiguity, we write the family $H_k^*(N)$ as $\cF$. We assume that the reader is familiar with standard properties of $L$-functions (see for example \cite{iwaniec-kwalski2004analytic}), and in particular with Satake parameters. We denote the Satake parameters of $f$ at $n$ by $\alpha_f(n)$ and $\beta_f(n)$. The Satake parameters at $n$ are related to to the corresponding Hecke eigenvalue, $\lambda(n)$. We have,
\begin{equation}
    \lambda_f(n) \ = \ \alpha_f(n) + \beta_f(n), \end{equation}\begin{equation}
        \alpha_f(n) \beta_f(n) \ =\ 1.
\end{equation}

Our main tool is the explicit formula (see \cite{ILS2} (4.11), for example), \begin{equation}\label{eq:1lvl density breakdown}
    D_1(f, \phi) \ = \ \frac{A_{k,N}(\phi)}{\log(R)} - 2\sum_{p}\sum_{m=1}^\infty \frac{\alpha_f(p)^{m}+\beta_f(p)^{m}}{p^{m/2}} \frac{\log(p)}{\log(R) } \hat\phi\left(m \frac{\log(p)}{\log(R) }\right)
\end{equation} where 
    $$A_{k,N} (\phi)\ :=\  2\hat\phi(0)\log\left(\frac{\sqrt{N}}{\pi} \right) + \int_{-\infty}^\infty \psi\left(\frac{k}{4}+ \frac{2\pi i x}{\log(R) }\right)\phi(x)dx + \int_{-\infty}^\infty \psi\left(\frac{k}{4}+ \frac{1}{2} +\frac{2\pi i x}{\log(R)}\right)\phi(x)dx . $$
We define
$$S(\phi):=-2\sum_{p}\sum_{m=1}^\infty \frac{\alpha_f(p)^{m}+\beta_f(p)^{m}}{p^{m/2}} \frac{\log(p)}{\log(R) } \hat\phi\left(m \frac{\log(p)}{\log(R) }\right)$$
so that using \eqref{eq:1lvl density breakdown}, we can write: 
\begin{align}
     D_1(f, \phi_1)D_1(f, \phi_2)&= \left(\frac{A_{k,N}(\phi_1)}{\log(R)} +S(\phi_1)\right)\left(\frac{A_{k,N}(\phi_2)}{\log(R)} +S(\phi_2)\right).
    \label{eq:rewriting (i) by explicit formula}
\end{align}
Expanding \eqref{eq:rewriting (i) by explicit formula}, we have \begin{equation}
     D_1(f, \phi_1)D_1(f, \phi_2) \ = \ \frac{A_{k,N}(\phi_1)A_{k,N}(\phi_2)}{\log^2(R)} +S(\phi_1)\left(\frac{A_{k,N}(\phi_{2})}{\log(R)}\right) +S(\phi_2)\left(\frac{A_{k,N}(\phi_{1})}{\log(R)}\right)+S(\phi_1)S(\phi_2).
\end{equation}

In addition, define \begin{equation}
    S_1(\cF,\phi)\ := \ -2\sum_{p}\sum_{m=1}^\infty\frac{1}{W_R(\cF)}\sum_{f\in\cF} w_R(f) \frac{\alpha_f(p)^m+\beta_f(p)^m}{p^{m/2}}\frac{\log(p)}{\log(R)}\hat\phi\left(m\frac{\log(p)}{\log(R)}\right)
\end{equation} and
\footnotesize \begin{align}
    S_2(\cF,\phi_1,\phi_2)\notag\ := \ 4\sum_{p_1,p_2}\sum_{\substack{m_1\in\N \\ m_2\in\N}} \frac{1}{W_R(\cF)}\sum_{f\in\cF} w_R(f) \frac{\alpha_f(p_1)^m_1+\beta_f(p_1)^m_1}{p_2^{m_2/2}}&\frac{\alpha_f(p_2)^m_2+\beta_f(p_2)^m_2}{p_2^{m_2/2}}\frac{\log(p_1)\log(p_2)}{\log^2(R)}\\&\times\hat\phi_1\left(m_1\frac{\log(p_1)}{\log(R)}\right)\hat\phi_2\left(m_2\frac{\log(p_2)}{\log(R)}\right).
\end{align}\normalsize
Then we can rewrite $D_1(\cF,\phi)$ and $D_2\left(\cF, \phi_1,\phi_2\right)$ respectively by
\begin{equation}
    D_1(\cF,\phi) \ = \ \lim_{N\to \infty}\frac{A_{k,N}(\phi)}{\log(R)} + S_1(\cF,\phi),
\end{equation}
\begin{align}
  D_2\left(\cF, \phi_1,\phi_2\right) \ =& \  \lim_{N\to \infty}\frac{A_{k,N}(\phi_1)A_{k,N}(\phi_2)}{\log^2(R)} +\left(\frac{A_{k,N}(\phi_{2})}{\log(R)}\right)S_1(\cF,\phi_1) \\\ &+\left(\frac{A_{k,N}(\phi_{1})}{\log(R)}\right)S_1(\cF,\phi_2)+S_2(\cF,\phi_1,\phi_2) \notag- 2\left(\frac{A_{k,N}(\phi_1\phi_2)}{\log(R)} + S_1(\cF,\phi_1\phi_2)\right)  \\ \ &+\phi_1(0)\phi_2(0)\frac{\sum_{f\in\cF}w_R(f)(1-\epsilon_f)}{2W_R(\cF)}. 
\end{align}
Therefore, by computing $A_{k,N}(\phi), ~ S_1(\cF,\phi),$ and $S_2(\cF,\phi_1,\phi_2)$ up to $O\left(\frac{1}{\log^4(R)}\right)$ error for arbitrary $\phi,\phi_1,$ and $\phi_2$ even Schwartz functions with compactly supported fourier transforms, we can calculate    $D_1(\cF,\phi)$ and $D_2\left(\cF,\phi_1,\phi_2\right)$ up to $O\left(\frac{1}{\log^4(R) }\right)$ error. 

%%%%%%%%%%%%%%%% MOTIVATION %%%%%%%%%%%%%%%%%%%%%

By doing so, we greatly extend Miller's work \cite{Miller2009_LOTerms_1LevelDensity}, allowing the level $N$ to be a product of two, not necessarily distinct, primes. The assumption of prime level greatly simplifies the inclusion-exclusion analysis and resulting formulas, as well as provides explicit formulas for the sign of the functional equation, which allows us to split by sign. When $N$ is not prime, we perform more delicate analysis with main ingredients coming from ILS \cite{ILS2} when $N$ is square-free, while we use tools from \cite{BarrettEtAl2016arXiv} for general $N$.

While all the families we study share the same main term, they have different lower-order terms, depending on the factorization of the level. We can compute up to errors of size $O((1/\log(N))^4)$. In all the cases we study, so long as the largest factor of $N$ is at least a fixed power smaller than $N$, the lower-order terms agree to this degree of precision with those from the case when the level is prime. On the other hand, the lower-order terms differ when the smallest prime is at most a given size. Therefore, whether or not there is a space of oldforms of comparable size to the newforms turns out to be the reason for the breaking of the universality of behavior. 

%%%%%%%%%%%%% OUR RESULT %%%%%%%%%%%%% 

Through computing the lower order terms for each of the four scenarios of how $N$ goes to infinity, we prove that for $N = q_1q_2,$ the case in which $q_1$ is fixed has different lower order terms compared to 
other cases where both factors go to infinity at $N^{c}$ rate for some $c \in (0,1),$ which agree with the $N$ prime case. In other words, as long as both factors of the level $N$ go to infinity fast enough, the factorization of the level does not affect the lower order terms of first nor second level density up to $O\left(1/{\log^4(R)}\right)$ error. We conjecture that the same phenomenon holds in general. 

\begin{conjecture}
 Given $N$, as long as the prime factors of $N$ go to infinity fast enough relative to the error we are computing up to, the lower order terms of the $n$ level density are the same as the $N$ prime case. Interesting new factors emerge when one of the factors goes to infinity slower than the reciprocal of the error term, say at the rate of $\log(N)$, $\log^2(N)$, or $\log^3(N)$. 
 \end{conjecture}

In Section \ref{chp: A_k,N compute}, we compute $A_{k,N}(\phi)$ up to $O\left(\frac{1}{\log^4(R) }\right)$ error. In Section \ref{chp: S1 compute chpt}, we compute explicitly $S_{A'}(\cF)$ and $S_{A}(\cF)$ where $$S_1(\cF,\phi) = S_A(\cF) + S_{A'}(\cF) + O\left(\frac{1}{\log^4(R)}\right),$$ extending \cite{Miller2009_LOTerms_1LevelDensity}'s result to improve the error term. See Theorem \ref{theorem: S_1} for definitions of $S_{A'}(\cF) $ and $S_{A}(\cF)$. In Section \ref{chp: S2 compute chp}, we compute $S_{B''}(\cF) $, $S_{B'}(\cF), S_{B_f}(\cF), $ and $S_{B_\infty}(\cF)$ to write $$S_2(\cF,\phi_1,\phi_2) =S_{B''}(\cF) +S_{B'}(\cF)+S_{B_f}(\cF)+S_{B_\infty}(\cF)+ O\left(\frac{1}{\log^4(R)}\right).$$ See Theorem \ref{thm: S_2.Formula} for definitions of $S_{B''}(\cF) $, $S_{B'}(\cF), S_{B_f}(\cF), $ and $S_{B_\infty}(\cF)$. Then, we have an explicit formula for $D_1(\cF)$ in terms of $A_{k,N}(\phi)$ and $S_1(\cF,\phi)$ up to $O\left(1/\log^2(R)\right)$ error and an explicit formula for $D_2(\cF)$ in terms of $A_{k,N}(\phi), ~ S_1(\cF,\phi),$ and $S_2(\cF,\phi_1,\phi_2)$ up to $O\left({1}/{\log^4(R)}\right)$ error. Once we have a formula for both $D_1(\cF)$ and $D_2(\cF)$, we start investigating how the level approaches infinity affects the lower-order terms of $D_1(\cF)$ and $D_2(\cF)$. In Section  \ref{chp: sum of weights and family specific terms}, we compute explicitly the terms $A'_{r,\cF}(p), \ A_{r,\cF}(p), B''_{r_1,r_2,\cF}(p_1,p_2), \ B'_{r_1,r_2,\cF}(p_1,p_2), \ $ and $B_{r_1,r_2,\cF}(p_1,p_2)$ which are defined repectively in \eqref{def:A'_r},\eqref{def:A_r}, \eqref{not:B'' i,j}, \eqref{not:B' i,j}, \eqref{not:B i,j}. If the family we are averaging over is clear, we omit specifying it in a subscript. These quantities determine the value of $S_1(\cF,\phi_1),$ $S_1(\cF,\phi_2),$ and $S_2(\cF,\phi_1,\phi_2)$ for each of the four families we are considering. To this end, we use the Peterson trace formula and \cite{BarrettEtAl2016arXiv}'s generalized formula. Lastly, in Section \ref{chp: lower order terms}, we substitute the results from Section \ref{chp: sum of weights and family specific terms} into the formula of $S_1(\cF,\phi)$ and $S_2(\cF,\phi_1,\phi_2)$ we obtained in Sections \ref{chp: S1 compute chpt} and \ref{chp: S2 compute chp} in order to evaluate the lower order terms. 

Computing the 2-level density, as opposed to calculating just the 1-level density, comes with various new obstructions. From doing inclusion-exclusion, the error term of $D_2(\cF)$ is dictated by the sharpness of the error term in the 1-level density; if we continue to use inclusion-exclusion for a general $n$-level density, the error term will continue to be dictated by the 1-level density. This is because terms $S_1(\cF,\phi)$ occurs with just one factor of $1/\log(R)$ while $S_2(\cF,\phi_1,\phi_2)$ occurs with  $1/\log^2(R)$. Therefore, we need to obtain more refined lower-order terms for 1-level density. Furthermore, as the level increases, the number of terms that need to be considered greatly increases; computing $S_2(\cF,\phi_1,\phi_2)$ requires far more terms than $S_1(\cF,\phi)$. In addition, looking at our work with the harmonic weights, new behaviors appear in the moments; there, careful case work is required depending on the factorization of the level. Notably, the case where $N=p^2$ is nearly identical to $N=q_1q_2$ with both factors going to infinity, while new challenges and behaviors emerge when one of the factors is fixed. 

\begin{remark}
One helpful analogy to our results due to Miller \cite{Miller2009_LOTerms_1LevelDensity} is demonstrated by the proof of the Central Limit Theorem (CLT). Suppose
$X_1,\ldots, X_N$ are 
\lq nice' independent, identically distributed random variables with mean $\mu$ and variance
$\sigma^2$. Then, the running average $(X_1 + \cdots +X_N-N\mu)/\sigma \sqrt{N} $ converges to the standard normal  $\mathcal{N}(0,1)$ as $N\to \infty$.
The universality suggested by the CLT is that, modulo normalization, the main terms are independent of the initial distribution. 
However, the higher moments of the
distribution determine the rate of convergence to $\mathcal{N}(0,1)$ (See the Berry-Esseen Theorem \cite{Berry1941}, \cite{Esseen1942}). We observe a similar phenomenon with the 2-level density. Universally, as suggested by Katz and Sarnak, we know what the main terms should be in the limit as $R \to \infty$. However, the higher moments of the Fourier coefficients determine the lower-order terms of the 2-level density.
\end{remark}
\begin{remark}
     Erik Samuelsson \cite{SamuelssonPC} of the University of Gothenburg is working on a similar problem, but for a fixed level ($N = 1$) with the weight $k$ going to infinity.
\end{remark}

%%%%%%%%%%%%%%%%%%%%%%%%%%%%%%%%%%%%%%%%%%%%%%%
%%%%%%%%%%%%%%%%%%%%%%%%%%%%%%%%%%%%%%%%%%%%%%
\vspace{1em}
\section{Preliminaries}
We define the $N^{\text{th}}$ congruence subgroup of $SL_2(\mathbb{Z})$ by
\begin{equation}\label{eq:cong-subgroup}
\Gamma_0(N) \ := \ \left\{ \begin{pmatrix} 
a & b \\ c & d 
\end{pmatrix} \in \mathrm{SL}_2(\mathbb{Z}) \ \Bigg| \ c \equiv 0 \pmod{N} \right\}.
\end{equation}
We are interested in how  $N^{\text{th}}$ congruence subgroup group acts on certain holomorphic functions $f: \mathbb{H}\to \C$, where $\mathbb{H}$ denotes the upper half complex plane. Specifically, we are interested in functions $f$ that admit a symmetry when acted on by $\Gamma_0(N)$ by the following relation:
\begin{equation}
f\left( \frac{az + b}{cz + d} \right)\ = \ (cz + d)^{k} f(z).
\end{equation}
If these functions satisfy the extra stipulation that they vanish at their cusps, we call these functions $f$ a holomorphic cusp form of weight $k$ and level $N$.
Such functions admit a Fourier expansion
\begin{equation}\label{eq:Fourier}
f(z)\ =\ \sum_{n=1}^{\infty}a_f(n)e^{2\pi inz},
\end{equation}
where the coefficients $a_f(n)$ are normalized so that $a_f(1)=1.$\footnote{Since we are looking at cusp forms, $a_f(0)=0$; for general forms, this is not necessarily true.}
We denote the space of all holomorphic cusp forms of weight $k$ and level $N$ by $S_k(N)$. This space is a finite-dimensional Hilbert space whose inner product is given by
\begin{equation}\label{eq:Petersson}
\langle f,g\rangle \ := \ \frac{1}{\nu(N)}\int_{\Gamma_0(N) \backslash\mathbb{H}}f(z)\overline{g(z)}\,y^k\,\frac{dx\,dy}{y^2},
\end{equation}
where $\nu(N):=[\mathrm{SL}_2(\mathbb{Z}):\Gamma_0(N)]$. This inner product is called the Petersson inner product. 

In 1970, Atkin and Lehner \cite{AL} built a theory on newforms of $S_k(N)$. For every form in $S_k(N)$, we can induce a form in $S_k(M)$ where $N|M$ and $M \neq N$. The induced forms are referred to as ``oldforms,'' while the rest, which form an orthogonal space to the space spanned by oldforms, are called ``newforms.'' While $S_k(N)$ contains several oldforms, we filter them out and focus on the remaining set $H^*_k(N)$ of exclusively newforms. These newforms satisfy many desirable properties like being eigenfunctions of the Hecke operators $T_n$ (Hecke eigenforms for short) for all $n \in \N.$. The Hecke operators $T_n$ are given by
\begin{equation}\label{eq:Hecke}
T_nf(z) \ := \ n^{k-1}\sum_{\substack{ad=n \\(a,N)=1}}\sum_{b=0}^{d-1}d^{-k}f\left(\frac{az+b}{d}\right) .
\end{equation} 
Moreover, we can find an orthogonal basis $\mathcal{B}_k(N)$ of the newforms. See \cite{ILS2} for construction of this basis.
Since $S_k(N)$ is finite-dimensional, we know $|\mathcal{B}_k(N)|$ is finite; in fact, from \cite[(2.73)]{ILS},
\begin{equation}
\left|\mathcal{B}_k(N)\right|\sim\frac{k-1}{12}\varphi(N)+O\left((kN)^{5/6}\right),
\end{equation} 
where $\varphi(N)$ is Euler's totient function. Given Hecke eigenform $f\in\mathcal{B}_k(N)$, we refer to its eigenvalue under $T_n$ as the $n^{\text{th}}$ Hecke eigenvalue of $f$, $\lambda_f(n)$. Thanks to Deligne \cite{Deligne1974}, we know that $\lambda_f(p)\in [-2,2]$. Moreover, Hecke eigenvalues of $f$ are closely related to its Fourier coefficients, satisfying the following relation 
\begin{equation}\label{eq:Fourier-Hecke}
    a_f(n) \ = \ \lambda_f(n)n^{(k-1)/2}.
\end{equation}
 Furthermore, the Hecke eigenvalues of $f$ possess useful multiplicative properties,
\begin{align}\label{eq:mult-prop}
\lambda_f(m)\lambda_f(n)\ =\sum_{\substack{d\mid(m,n)\\ (d,N)=1}}\lambda_f\left(\frac{mn}{d^2}\right),\end{align}
which means that if $(m,n)=1$, then
\begin{align}\lambda_f(mn)=\lambda_f(m)\lambda_f(n).
\end{align}
Using this multiplicative property, we define the $L$-function associated to $f$ as:
\begin{equation}
L(s,f) \ := \ \sum_{n=1}^{\infty}\frac{\lambda_f(n)}{n^s}, \quad \Re(s)>1.
\end{equation}

For a cusp form $f$, we define the normalized Fourier coefficients by 
$$\Psi_f(n) \ := \ \left( \frac{\Gamma(k-1)}{(4\pi)^{k-1}}\right)^{1/2}||f||^{-1}\lambda_f(n)$$
with $||f||^2=\langle f,f\rangle.$ We wish to consider the sum
$$\Delta_{k,N}(m,n)=\sum_{f \in \mathcal{B}_k(N)}\bar{\Psi_f(m)}\Psi_f(n)$$
where $\mathcal{B}_k(N)$ is an orthonormal basis of $S_k(N).$ This sum is computed using the Petersson trace formula \cite{Petersson1932}\cite{iwaniec-kwalski2004analytic}. 
\begin{prop}\label{eq:peter} We have
    \begin{equation}
        \Delta_{k,N}(m,n)\ =\ \delta(m,n)+2\pi i^k+\sum_{c \equiv 0 \text{ mod}\: N}c^{-1}S(m,n;c)J_{k-1}\left( \frac{4\pi \sqrt{mn}}{c}\right)
         \end{equation}
    where $\delta(m,n)=1$ if $m \ = \ n$ and is 0 otherwise, $J_{k-1}\left( x\right)$ denotes the Bessel function, and
    $S(m,n;c)$ is the classical Kloosterman sum.
\end{prop}
Effective bounding yields the following estimation from \cite{ILS2}.
\begin{prop}\label{prop:peter aprox. ILS2.2}
For $m,n \geq 1$,
        \begin{equation}
        \Delta_{k,N}(m,n)\ =\ \delta(m,n)+O\left(\frac{\tau(N)}{k^{5/6}N}\frac{(m,n,N)\tau_3((m,n))}{((m,N)+(n,N))^{1/2}} \left(\frac{mn}{(mn)^{1/2}+kN} \right)^{1/2}\log 2mn\right),
        \end{equation}
    where the implied constant is absolute and $\tau_3(\ell)=\{\#(a,b,c)\; | \; abc=\ell \}$
\end{prop}
Moreover, we define
\begin{align}\label{def:Z_func}
Z(s, f)\ & := \ \sum_{n=1}^{\infty} \frac{\lambda_f\left(n^2\right)}{n^s}\ =\ \frac{\zeta_N(s) L(s, f \otimes f)}{\zeta(s)},
\\\label{def:Z_N_func}
Z_N(s,f)\ & :=\ \sum_{n\mid N^\infty} \frac{\lambda_f(n^2)}{n^s}.
\end{align}
We wish to evaluate the arithmetically weighted sum 
\begin{align}\label{def:Deltastar}
    \Delta^*_{k,N}(m,n)\ =\ \sum_{f \in H^*_k(M)}\frac{Z_N(1,f)}{Z(1,f)}\lambda_f(m)\lambda_f(n).
\end{align}
Iwaniec, Luo, and Sarank derived the following expression in the case of $N$ squarefree.
\begin{prop} \label{prop:ILS_Lem_2.7}
    Let $N$ be squarefree, $(m,N)=1$, and $(n,N^2)|N$. Then
    \begin{equation}\label{eq:harm avg}
        \Delta^*_{k,N}(m,n)\ =\ \frac{k-1}{12}\sum_{ML=N}\frac{\mu(L)M}{\nu((n,L))}\sum_{\ell |L^{\infty}}\ell^{-1}\Delta_{k,M}(m\ell^2,n).
    \end{equation}
\end{prop}
In addition, Corollary 2.10 in \cite{ILS2} provides the following estimate: 
\begin{prop}\label{eq:harm avg aprox}
 Let $N$ be squarefree, $(m,N)=1$, and $(n,N^2)|N$. Then,
    \begin{equation} 
        \Delta^*_{k,N}(m,n)\ =\ \frac{k-1}{12}\varphi(N)\delta(m,n)+O\left(k^{1/6}(mn)^{1/4}(n,N)^{-1/2}\tau^2(N)\tau_3((m,n))\log 2mnN \right). \label{aproxILS}
    \end{equation}
\end{prop}
  As Proposition \ref{eq:harm avg aprox} requires $N$ to be squarefree, for general $N$, we use proposition 4.1 in \cite{BarrettEtAl2016arXiv} which removes the squarefree restriction on $N$.
    \begin{prop} \label{prop:barrett}
    Let $(m,N)=1$ and $(n,N)=1$. Then
 \begin{equation} 
        \Delta^*_{k,N}(m,n)\ =\ \frac{k-1}{12}\sum_{ML=N}\prod_{p^2|M}\left( \frac{p^2}{p^2-1}\right)^{-1}\sum_{\substack{\ell |L^{\infty}\\(\ell,M)=1}}\ell^{-1}\Delta_{k,M}(m\ell^2,n).
    \end{equation}
\end{prop}
%%%%%%%%%%%%%%%%%%%%%%%%%%%%%%%%%%%%%%%%%%%%%%%
%%%%%%%%%%%%%%%%%%%%%%%%%%%%%%%%%%%%%%%%%%%%%%


\section{Computing \texorpdfstring{$\mathbf{A_{k,N}(\phi)}$}{A{k,N}(phi)}} \label{chp: A_k,N compute}
We start the computation of the 1st and 2nd level density by computing ${A_{k,N}(\phi)}$ up to $O\left(\frac{1}{\log^4(R)}\right)$ error. 
\begin{theorem}\label{gammaf}  Let $\phi$ be an even Schwartz function whose Fourier transform has compact support and let $\psi$ be the digamma function. We have the following estimate.
    \begin{align}
A_{k,N} (\phi)& \ = \ \hat\phi(0)\log(N) +\hat\phi(0)  \left(\psi\left(\frac{k}{4}\right)+\psi\left(\frac{k}{4}+ \frac{1}{2}\right)-2\log(\pi)\right) \notag \\ &~~~~- \frac{2\pi^2}{\log^2(R)}\left(\int_{-\infty}^\infty \phi(x)x^2 dx\right)\left(\psi''\left(\frac{k}{4}\right)+\psi''\left(\frac{k}{4}+ \frac{1}{2}\right)\right) +O\left(\frac{1}{\log^4(R)}\right).
\end{align} 
\end{theorem}
\textit{Proof of Theorem \ref{gammaf}.} We start by proving a lemma.
\begin{lemma}  Let $\phi$ be an even Schwartz function whose Fourier transform has compact support and let $\psi$ be the digamma function. We have the following estimate.
    \begin{equation}
        \int_{-\infty}^\infty \phi(x)\psi\left(\frac{k}{4}+ \frac{2\pi i x}{\log(R)}\right) dx \ =\ \psi\left(\frac{k}{4}\right)\hat\phi(0) -\frac{\psi^{''}(\frac{k}{4})2\pi^2}{\log^2(R)}\int_{-\infty}^\infty \phi(x)x^2 dx + O\left(\frac{1}{\log^4(R)}\right)
    \end{equation}
\end{lemma}
\begin{proof}
Using (8.363.3) of \cite{GradshteynRyzhik1965} gives us that the leading term is $\psi(\frac{k}{4})\hat\phi(0)$. Therefore, we subtract the leading term and compute the error term more accurately. 
Using the fact that $\phi$ is even
\begin{align}
    &\int_{-\infty}^\infty \phi(x)\psi\left(\frac{k}{4}+ \frac{2\pi i x}{\log(R)}\right) dx - \psi\left(\frac{k}{4}\right)\hat\phi(0)\notag\\ & \ = \
    \frac{1}{2}\int_{-\infty}^\infty \phi(x)\left[\psi\left(\frac{k}{4}+ \frac{2\pi i x}{\log(R)}\right) + \psi\left(\frac{k}{4}- \frac{2\pi i x}{\log(R)}\right) - 2\psi\left(\frac{k}{4}\right)\right]dx. \label{eq:goal 1 what we want di gamma expand}
\end{align}

Suppose $f:\Omega\to \C,$ where $\Omega$ is an open subset of $\C$. Suppose further that $z_0 \in \Omega$. Taylor expanding about the $z_0 \in \Omega$ yields
\begin{equation}\label{eq:symm Taylor expansion up to x^4 term}
    f(z_0+z) + f(z_0-z) - 2f(z_0) \ =\  f''(z_0)z^2 + O(z^4)
\end{equation} for all $z$ such that $|z|<R$ where $R$ is the radius of convergence. Moreover, we know from Cauchy's Inequality that $R \geq \text{dist}(z_0,\del\Omega)$. 
We apply the formula in \eqref{eq:symm Taylor expansion up to x^4 term} with $\Omega = \{z\in\C : \text{Re}(z) > 0\}$, $f(z) = \psi(z)$, $z_0 = \frac{k}{4}$, and $z = 2\pi ix/\log(R)$. We remark that $\text{dist}(\frac{k}{4},\del(\{z\in\C : \text{Re}(z) > 0\})) = \frac{k}{4}$ and hence $R\geq \frac{k}{4}$. Then, we get
\begin{equation}\label{eq:digamma taylor expand}
    \psi\left(\frac{k}{4} + z\right) +  \psi\left(\frac{k}{4} - z\right) - 2 \psi\left(\frac{k}{4}\right) \ =\  \psi^{''}\left(\frac{k}{4}\right) z^2 + O(z^4)
\end{equation} for all $x$ such that $|x|<k\log(R)/(8\pi)$.

Now we integrate \eqref{eq:digamma taylor expand} against the test function $\phi$ for $|x| < k\log(R)/(8\pi)$ and \eqref{eq:goal 1 what we want di gamma expand}. Then
\begin{align}\label{eq: digamma substution}
&\int_{-k\log(R)/(8\pi)}^{k\log(R)/(8\pi)}\phi(x)\psi\left(\frac{k}{4}+ \frac{2\pi i x}{\log(R)}\right) dx - \psi\left(\frac{k}{4}\right)\hat\phi(0) \notag \\ & \ = \ \frac{1}{2}\int_{-k\log(R)/(8\pi)}^{k\log(R)/(8\pi)} \phi(x)\left[-\psi''\left(\frac{k}{4}\right)\left(\frac{2\pi x}{\log(R)}\right)^2 + O\left(\frac{2\pi x}{\log^4(R)}\right)\right]dx \notag \\
& \ = \ \frac{1}{2}\int_{-k\log(R)/(8\pi)}^{k\log(R)/(8\pi)} \phi(x)\left[-\psi''\left(\frac{k}{4}\right)\left(\frac{2\pi x}{\log(R)}\right)^2\right]dx + O\left(\frac{1}{\log^4(R)}\right)\left(\int_{-k\log(R)/(8\pi)}^{k\log(R)/(8\pi)} \phi(x)x^4dx\right) \notag \\ 
& \ = \ \frac{1}{2}\int_{-k\log(R)/(8\pi)}^{k\log(R)/(8\pi)} \phi(x)\left[-\psi''\left(\frac{k}{4}\right)\left(\frac{2\pi x}{\log(R)}\right)^2\right]dx + O\left(\frac{1}{\log^4(R)}\right).
\end{align}

We now show the tail part of the integral is negligible. Consider the series expansion of $\psi(z)$ from \cite{abramowitz1972handbook}, which holds for all $z \notin \Z^{\leq 0}$:
\begin{equation}
\psi(z) \ = \ -\gamma + \sum_{n=0}^\infty \frac{z-1}{(n+1)(n+z)}
\end{equation}
where $\gamma$ is the Euler–Mascheroni constant. For $z = \frac{k}{4}+2\pi i x/ \log(R),$ we then have that for all $n \geq 1,$:
\begin{align}\label{ineq:helping convergence of tail}
  \left|  \frac{\left(\frac{k}{4}+2\pi i x/ \log(R)\right)-1}{(n+1)(n+\left(\frac{k}{4}+2\pi i x/ \log(R)\right))}\right| & \ = \ \left|\frac{8 i \pi x + (-4 + k) \log(R)}{(1 + n)\left(8 i \pi x + (k + 4n) \log(R)\right)}\right| \notag \\
  & \ \leq \ \frac{8\pi x + |k-4| \log(R)}{(1 + n)\left((k + 4n) \log(R)\right)} \notag \\
  & \ \leq \ \frac{8\pi x+|k-4|\log(R)}{n^2}.
\end{align} 
We also know that because $\phi(x)$ is Schwartz, $\int_{k\log(R)/(8\pi)}^\infty \phi(x) P(x)dx = O\left(\frac{1}{\log^B R}\right)$ for any polynomial $P(x)$. 
Therefore, using this fact along with \eqref{ineq:helping convergence of tail}, we find 
\begin{align}\label{eq: remainder part 1/R} &\left|\int_{k\log(R)/(8\pi)}^\infty  \phi(x)\left(-\gamma + \sum_{n=0}^\infty \frac{z-1}{(n+1)(n+z)}\right) dx \right|\notag\\
   & \ \leq \ (-\gamma +1)\left|\int_{k\log(R)/(8\pi)}^\infty  \phi(x)dx \right|+  \sum_{n=1}^\infty\left|\int_{k\log(R)/(8\pi)}^\infty\phi(x)\left(\frac{8\pi x+|k-4|\log(R)}{n^2}\right)dx \right| \notag\\
   & \ = \ \left(-\gamma +1+ \left(\sum_{n=1}^\infty \frac{|k-4|\log(R)}{n^2}\right)\right)\left|\int_{k\log(R)/(8\pi)}^\infty  \phi(x)dx\right| +  8\pi \left(\sum_{n=1}^\infty \frac{1}{n^2}\right)\left|\int_{k\log(R)/(8\pi)}^\infty  x\phi(x)dx\right|
   \notag \\
   & \ = \ O\left(\frac{1}{\log^4(R)}\right).
\end{align} Similarly, we know that $$\left|\int_{k\log(R)/(8\pi)}^\infty  \phi(x)\left(-\gamma + \sum_{n=0}^\infty \frac{z-1}{(n+1)(n+z)}\right) dx \right| \ = \ O\left(\frac{1}{\log^4(R)}\right).$$ 
In addition, we know that \begin{equation} \label{eq: other piece of remainder}
    \int_{k\log(R)/(8\pi)}^{\infty} \phi(x)\left[-\psi''\left(\frac{k}{4}\right)\left(\frac{2\pi x}{\log(R)}\right)^2\right]dx + \int_{-\infty}^{-k\log(R)/(8\pi)} \phi(x)\left[-\psi''\left(\frac{k}{4}\right)\left(\frac{2\pi x}{\log(R)}\right)^2\right]dx
\end{equation} is  $O\left({1}/{\log^4(R)}\right)$ because $\phi$ is Schwartz.
Therefore, combining \eqref{eq: digamma substution},  \eqref{eq: remainder part 1/R}, and \eqref{eq: other piece of remainder}, and noticing that \begin{align}
    \int_{-\infty}^{\infty} \phi(x)x^2dx = -\hat\phi''(0)
\end{align} we have that
\begin{equation}
\int_{-\infty}^\infty \phi(x)\psi\left(\frac{k}{4}+ \frac{2\pi i x}{\log(R)}\right) dx - \psi\left(\frac{k}{4}\right)\hat\phi(0) \  =  \ \frac{2\pi^2\psi''\left(\frac{k}{4}\right)}{\log^2(R)}\hat\phi''(0) + O\left(\frac{1}{\log^4(R)}\right)    
\end{equation} as desired.
\end{proof}
By the same argument, but for $z_0 = \frac{k}{4}+ \frac{1}{2}$, we get the following lemma. 
\begin{lemma} Let $\phi$ be an even Schwartz function whose Fourier transform has compact support and let $\psi$ be the digamma function. Then
      \begin{equation}
        \int_{-\infty}^\infty \phi(x)\psi\left(\frac{k}{4}+ \frac{1}{2}+ \frac{2\pi i x}{\log(R)}\right) dx \ =\ \psi\left(\frac{k}{4} + \frac{1}{2}\right)\hat\phi(0) +\frac{2\pi^2\psi^{''}(\frac{k}{4} + \frac{1}{2})}{\log^2(R)}\hat\phi''(0) + O\left(\frac{1}{\log^4(R)}\right). 
        \end{equation}
    \end{lemma}
Using the two above lemmas, we have Theorem \ref{gammaf}. \qed

\vspace{1em} 

%%%%%%%%%%%%%%%%%%%%%%%%%%%%%%%%%%%%%%%%%%%%%%%
%%%%%%%%%%%%%%%%%%%%%%%%%%%%%%%%%%%%%%%%%%%%%%%

\vspace{1em}

\section{\texorpdfstring{Computing $\mathbf{S_1(\cF,\phi)}$ up to $O\left(\frac{1}{\log^4(R)}\right)$ error}{Computing S1 up to 1/log4 error}}\label{chp: S1 compute chpt}
Similar to Miller \cite{Miller2009_LOTerms_1LevelDensity}, we define 
\begin{align} \label{def:A'_r}
    A'_{r,\cF}(p) \ &:=\ \frac{1}{W_R(\cF)}\sum_{\substack{f \in \cF \\ p|N}}w_R(f)\lambda_f(p)^r\\ \label{def:A_r}
    A_{r,\cF}(p)\  &:= \ \frac{1}{W_R(\cF)}\sum_{\substack{f \in \cF \\ p \nmid N}}w_R(f)\lambda_f(p)^r
\end{align}
where $W_R(\cF):=\sum_{f \in \cF}w_R(f)$. This brings us to the following theorem.
\begin{theorem}\label{theorem: S_1}
    Let $$\mathcal{S}_1( \mathcal{F},\phi)\ :=\ -2\sum_p\sum_{m=1}^{\infty}\frac{1}{W_{R}(\mathcal{F})}\sum_{f \in \mathcal{F}}w_R(f)\frac{\alpha_f(p)^m+\beta_f(p)^m}{p^{m/2}}\frac{\log(p)}{\log(R)}\hat{\phi}\left(m \frac{\log(p)}{\log(R)} \right),$$
    where $\log(R)$ is the average log conductor, then 
\begin{align}
\mathcal{S}_1( \mathcal{F},\phi)\ =& \ \ S_{A'}(\cF) + S_{A}(\cF) \end{align} where \small\begin{align}
  S_{A'}(\cF)\ :=\ & -2\sum_p\sum_{m=1}^{\infty}\frac{A'_{m, \mathcal{F}}(p)}{p^{m/2}}\frac{\log(p)}{\log(R)}\hat{\phi}\left(m \frac{\log(p)}{\log(R)} \right)  \\
  S_{A}(\cF)\ :=\ &-2\hat{\phi}\left(0 \right)\sum_p\frac{2A_{0, \mathcal{F}}(p)\log(p)}{p(p+1)\log(R)}+2\sum_p
\frac{2A_{0, \mathcal{F}}(p)\log(p)}{p\log(R)}\hat{\phi}\left(2 \frac{\log(p)}{\log(R)} \right) \notag \\ 
&-2\sum_p\frac{A_{1, \mathcal{F}}(p)\log(p)}{p^{1/2}\log(R)}\hat{\phi}\left(\frac{\log(p)}{\log(R)} \right)+2\hat{\phi}\left(0 \right)\sum_p\frac{A_{1,\mathcal{F}}(p)(3p+1)\log(p)}{p^{1/2}(p+1)^2\log(R)} \notag \\
&-2\sum_p\frac{A_{2,\mathcal{F}}(p)\log(p)}{p\log(R)}\hat{\phi}\left(2 \frac{\log(p)}{\log(R)} \right)+2\hat{\phi}\left(0 \right) \sum_p\frac{A_{2, \mathcal{F}}(p)(p^2+3p+1)\log(p)}{p(p+1)^3 \log(R)} \notag \\
&+\hat{\phi}''(0)\sum_p\frac{A_{0, \mathcal{F}}(p)(32p^2+24p+8)\log^3(p)}{p(p+1)^3\log(R)}-\hat{\phi}''(0)\sum_p\frac{A_{1, \mathcal{F}}(p)(27p^3-17p^2+5p+1)\log^3(p)}{p^{1/2}(p+1)^4}
\notag\\
&-\hat{\phi}''(0)\sum_p\frac{A_{2, \mathcal{F}}(p)(64p^4-4p^3+44p^2+20p+4)\log^3(p)}{p(p+1)^5\log^3(R)}\\&-2\hat{\phi}\left(0 \right)\sum_{p}\sum_{r=3}^{\infty}\frac{A_{r, \mathcal{F}}(p)p^{r/2}(p-1)\log(p)}{(p+1)^{r+1}\log(R)} \notag
\\
&+\hat{\phi}''\left(0 \right)\sum_p\sum_{r=3}^{\infty}\frac{A_{r, \mathcal{F}}(p-1)(r^2(p-1)^2-12rp-8p)p^{r/2}\log^3(p)}{(p+1)^{r+3}\log(R)}+O\left(\frac{1}{\log^4(R)} \right).
\end{align}\normalsize
\end{theorem}  
\text{Proof.} Identical to Miller's analysis in \cite{Miller2009_LOTerms_1LevelDensity}, we break the explicit formula from \cite{ILS2} into the case when $p|N$ and $p\nmid N$ and use properties of the Satake parameters to obtain
\begin{align}
\mathcal{S}_1({\mathcal{F}},\phi)\ =&\ -2\sum_p\sum_{m=1}^{\infty}\frac{1}{W_R(\mathcal{F})}\sum_{\substack{f \in \mathcal{F} \notag \\p|N}}w_R(f)\frac{\lambda_f(p)^m}{p^{m/2}}\frac{\log(p)}{\log(R)}\hat{\phi}\left(m \frac{\log(p)}{\log(R)} \right) \notag \\
&-2\sum_p\frac{1}{W_R(\mathcal{F})}\sum_{\substack{f \in \mathcal{F} \notag \\p \nmid N}}w_R(f)\frac{\lambda_f(p)}{p^{1/2}}\frac{\log(p)}{\log(R)}\hat{\phi}\left(\frac{\log(p)}{\log(R)} \right) \notag \\
&-2\sum_p\frac{1}{W_R(\mathcal{F})}\sum_{\substack{f \in \mathcal{F}\\p \nmid N}}w_R(f)\frac{\lambda_f(p)^2-2}{p}\frac{\log(p)}{\log(R)}\hat{\phi}\left(2\frac{\log(p)}{\log(R)} \right) \notag
\\
&-2\sum_p\sum_{m=3}^{\infty}\frac{1}{W_R(\mathcal{F})}\sum_{\substack{f \in \mathcal{F}\\p \nmid N}}w_R(f)\frac{\alpha_f(p)^m+\beta_f(p)^m}{p^{m/2}}\frac{\log(p)}{\log(R)}\hat{\phi}\left(m \frac{\log(p)}{\log(R)} \right)\label{eq:yes3}.
\end{align}
The first three sums are already in the desired form, as they can easily be expressed as weighted averages of Hecke eigenvalues. We now look to evaluate the last sum. Through Taylor expansion, we see that 
\begin{equation}
   \hat{\phi}(mx)-\hat{\phi}(x)\ =\ \frac{m^2-1}{2}\hat{\phi}''(0)x^2+O\left((m^4-1)x^4 \right). 
\end{equation}
Thus we have 
\begin{equation}\label{eq:taylorexpansionphi}
    \hat{\phi}\left(m\frac{\log(p)}{\log(R)}\right)-\hat{\phi}\left(\frac{\log(p)}{\log(R)}\right)\ =\ \frac{m^2-1}{2}\hat{\phi}''(0)\left( \frac{\log(p)}{\log(R)}\right)^2+O\left((m^4-1)\left(\frac{\log(p)}{\log(R)} \right)^4 \right).
\end{equation}
Substituting (\ref{eq:taylorexpansionphi}) into the last sum in (\ref{eq:yes3}), we get
\begin{align}\label{eq:S1 first}
  \sum_p\sum_{m=3}^{\infty}\frac{1}{W_R(\mathcal{F})}\sum_{\substack{f \in \mathcal{F}\\p \nmid N}}w_R(f)\frac{\alpha_f(p)^{m}+\beta_f(p)^{m}}{p^{m/2}} \frac{\log(p)}{\log(R)}\left(\hat{\phi} \left(\right. \right. &\left. \left. \frac{\log(p)}{\log(R)} \right)+\frac{m^2-1}{2}\hat{\phi}''(0)\left( \frac{\log(p)}{\log(R)}\right)^2 \right. \notag \\ &+\left. O\left((m^4-1)\left(\frac{\log(p)}{\log(R)} \right)^4 \right) \right). 
\end{align}
\noindent Define \begin{align}\label{def: M_(c,k)}
    M_{c,k}(p)\ :=\ \sum_{m=c}^{\infty}m^k\frac{\alpha_f(p)^{m}+\beta_f(p)^{m}}{p^{m/2}}.
\end{align} We now write equation \eqref{eq:S1 first}
as \begin{align}
    \sum_p\frac{1}{W_R(\mathcal{F})}\sum_{\substack{f \in \mathcal{F}\\p \nmid N}}w_R(f)M_{3,0}(p) &\frac{\log(p)}{\log(R)}\hat{\phi} \left(\frac{\log(p)}{\log(R)} \right)\\
    &+\frac{\hat{\phi}''(0)}{2}\sum_p\frac{1}{W_R(\mathcal{F})}\sum_{\substack{f \in \mathcal{F}\\p \nmid N}}w_R(M_{3,2}(p)-M_{3,0}(p))\frac{\log^3(p)}{\log^3 R} \notag \\
    &+O\left(\sum_p\frac{1}{W_R(\mathcal{F})}\sum_{\substack{f \in \mathcal{F}\\p \nmid N}}w_R(M_{3,4}(p)-M_{3,0}(p))\frac{\log^5 (p)}{\log^4(R)}  \right).
\end{align}
Substituting $\hat{\phi}\left(\frac{\log(p)}{\log(R)} \right)\ =\ \hat{\phi}(0)+\frac{\hat{\phi}''(0)}{2}\frac{\log^2(p)}{\log^2(R)}+O\left(\frac{1}{\log^4(R)}\right)$ gives
\begin{align} \label{eq: S1 fourth}
    \hat{\phi}(0)\sum_p\frac{1}{W_R(\mathcal{F})}\sum_{\substack{f \in \mathcal{F}\\p \nmid N}} w_R(f)M_{3,0}(p)\frac{\log(p)}{\log(R)}+\frac{\hat{\phi}''(0)}{2}\sum_p\frac{1}{W_R(\mathcal{F})}\sum_{\substack{f \in \mathcal{F}\\p \nmid N}} w_R(f)M_{3,2}(p)\frac{\log^3(p)}{\log^3(R)}& \notag \\
    +O\left(\sum_p\frac{1}{W_R(\mathcal{F})}\sum_{\substack{f \in \mathcal{F}\\p \nmid N}}w_R(f)(M_{3,4}(p)-M_{3,0}(p))\frac{\log^5 (p)}{\log^4(R)}  \right)&
\end{align}
In Miller's proof of Theorem 1.1, he shows
$$M_{3,0}(p)\ =\ \frac{2}{p(p+1)}-\frac{p^{1/2}(3p+1)}{p(p+1)^2}\lambda_f(p)-\frac{(p^2+3p+1)}{p(p+1)^3}\lambda_f(p)^2+\sum_{m=3}^{\infty}\frac{p^{m/2}(p-1)\lambda_f(p)^m}{(p+1)^{m+1}}.$$
We now compute $M_{3,2}(p).$ See Appendix \ref{appendix: Proof of Lemma M3,2} for the proof.
\begin{lemma}\label{Lemma, S_1}
    For prime $p$, we have 
    \begin{align} \label{eq: M_3,2 lemma compute}
        M_{3,2}(p)\ =&\ \frac{32p^2+24p+8}{p(p+1)^3} - \frac{27p^3-17p^2+5p+1}{\sqrt{p}(p+1)^4}\lambda_f(p)-\frac{64p^4-4p^3+44p^2+20p+4}{p(p+1)^5} \lambda_f(p)^2
        \notag\\&+ \sum_{r=3}^{\infty}\frac{(p-1)(r^2(p-1)^2-12rp-8p)p^{r/2}\lambda_f(p)^r}{(p+1)^{r+3}}.
    \end{align} 
\end{lemma}
Substituting equation \eqref{eq: M_3,2 lemma compute} into equation \eqref{eq: S1 fourth} yields the main term in \eqref{theorem: S_1}.
Now, we look at the error term
\begin{align}
    \sum_p\sum_{m=3}^{\infty}\frac{1}{W_R(\mathcal{F})}\sum_{\substack{f \in \mathcal{F}\\p \nmid N}}\frac{\alpha_f(p)^{m}+\beta_f(p)^{m}}{p^{m/2}} \frac{\log(p)}{\log(R)}\left(O\left((m^4-1)\left(\frac{\log(p)}{\log(R)} \right)^4 \right) \right).\label{eq:MOh}
\end{align}
As $\mid \alpha_f(p)^m+\beta_f(p)^m\mid \ \leq \ 2$ for all $m \ \in \ \N$, we have
\begin{align}
   \eqref{eq:MOh}\ &=\ O \left( \sum_p\sum_{m=3}^{\infty}\frac{1}{W_R(\mathcal{F})}\sum_{\substack{f \in \mathcal{F}\\p \nmid N}}\frac{\alpha_f(p)^{m}+\beta_f(p)^{m}}{p^{m/2}} \frac{\log(p)}{\log(R)}\left((m^4-1)\left(\frac{\log(p)}{\log(R)} \right)^4 \right) \right) \notag \\
   \ &=\ O \left( \sum_p \left( \frac{\log(p)}{\log(R)} \right)^5 \sum_{m=3}^{\infty} \sum_{\substack{f \in \mathcal{F}\\p \nmid N}} \frac{2 w_R(f)}{W_R(\mathcal{F})}\frac{1}{p^{m/2}} \left(m^4-1\right) \right) \notag \\
   \ &=\ O \left( \sum_p \left( \frac{\log(p)}{\log(R)} \right)^5 \sum_{m=3}^{\infty} \frac{m^4-1}{p^{m/2}} \right).
\end{align}

Thus we have 
\begin{align}
    \sum_p \left( \frac{\log(p)}{\log(R)} \right)^5 \sum_{m=3}^{\infty} \frac{m^4-1}{p^{m/2}} \lesssim \frac{1}{\log^4(R)}\sum_p\frac{\log^5 (p)}{p^{3/2}}
    &\lesssim  \frac{1}{\log^5(R)}.\\
\end{align}
As the sum $\sum_p\frac{\log^5 (p)}{p^{3/2}}$ converges by the comparison test, our error term is sufficient.
\qed
%%%%%%%%%%%%%%%%%%%%%%%%%%%%%%%%%%%%%%%%%%%%%%%
%%%%%%%%%%%%%%%%%%%%%%%%%%%%%%%%%%%%%%%%%%%%%%%

\vspace{1em}

\newpage

\section{\texorpdfstring{Computing ${S_2(\cF,\phi_1,\phi_2)}$ up to ${O\left(\frac{1}{\log^4(R)}\right)}$ error}{Computing S2 up to 1/log4 R error}} \label{chp: S2 compute chp}

We are interested in finding a formula for $S_2(\cF,\phi_1,\phi_2)\:=\frac{1}{W_R(\cF)}\sum_{f\in \cF}w_R(f)S(\phi_1)S(\phi_2)$ where
\begin{align} \label{eq: S_1S_2 formula}
    S(\phi_1)S(\phi_2) \ = \  &\left(\sum_{p_1}\sum_{m_1=1}^\infty \frac{\alpha_f(p_1)^{m}+\beta_f(p_1)^{m_1}}{p_1^{m_1/2}} \frac{\log(p_1)}{\log(R)} \hat\phi_1\left(m_1 \frac{\log(p_1)}{\log(R)}\right)\right)\notag\\&\cdot \left(\sum_{p_2}\sum_{m_2=1}^\infty \frac{\alpha_f(p_2)^{m_2}+\beta_f(p_2)^{m_2}}{p^{m_2/2}} \frac{\log(p_2)}{\log(R)} \hat\phi_2\left(m_2 \frac{\log(p_2)}{\log(R)}\right)\right)
\end{align} up to $O\left(\frac{1}{\log^4(R)}\right)$ error term.

For fixed $f \in \cF,$ we break this sum given by \eqref{eq: S_1S_2 formula} into four cases depending on whether each $p_1$ and $p_2$ divide $N_f$ or not. We freely change the order of summation because convergence is guaranteed by $\hat\phi_1$ and $\hat\phi_2$ being Schwartz. Using the fact that if $p|N_f$, then $\alpha_f(p)^m+\beta_f(p)^m\ = \ \lambda_f(p)^m$, we obtain: 
{\footnotesize\begin{align} \label{eq:S_2 formula}
    &S_2(\cF,\phi_1,\phi_2)\ = \ \notag\\ &~~~~~ \frac{1}{W_R(\cF)}\sum_{p_1,p_2}\sum_{\substack{m_1\in N\\m_2\in N}}\sum_{\substack{f\in \cF\\p_1|N_f\\p_2|N_f }}w_R(f)\frac{\lambda_f(p_1)^{m_1}}{p_1^{m_1/2}}  \frac{\lambda_f(p_2)^{m_2}}{p_2^{m_2/2}} \frac{\log(p_1)\log(p_2)}{\log^2(R)}\hat\phi_1 \left(m_1 \frac{\log(p_1)}{\log(R)}\right)\hat \phi_2 \left(m_2 \frac{\log(p_2)}{\log(R)}\right)\notag \\
    &+ \frac{1}{W_R(\cF)}\sum_{p_1,p_2}\sum_{\substack{m_1\in N\\m_2\in N}}\sum_{\substack{f\in \cF\\p_1|N_f\\p_2\nmid N_f }}w_R(f)\frac{\lambda_f(p_1)^{m_1}}{p_1^{m_1/2}}\frac{\alpha_f(p_2)^{m_2}+\beta_f(p_2)^{m_2}}{p_2^{m_2/2}} \frac{\log(p_1)\log(p_2)}{\log^2(R)}\hat\phi_1 \left(m_1 \frac{\log(p_1)}{\log(R)}\right)\hat \phi_2 \left(m_2 \frac{\log(p_2)}{\log(R)}\right)  \notag\\
    &+ \frac{1}{W_R(\cF)}\sum_{p_1,p_2}\sum_{\substack{m_1\in N\\m_2\in N}}\sum_{\substack{f\in \cF\\p_1\nmid N_f\\p_2|N_f }}w_R(f)\frac{\alpha_f(p_1)^{m_1}+\beta_f(p_1)^{m_1}}{p_1^{m_1/2}} \frac{\lambda_f(p_2)^{m_2}}{p_2^{m_2/2}}\frac{\log(p_1)\log(p_2)}{\log^2(R)}\hat\phi_1 \left(m_1 \frac{\log(p_1)}{\log(R)}\right)\hat \phi_2 \left(m_2 \frac{\log(p_2))}{\log(R)}\right)\notag \\
    &+\frac{1}{W_R(\cF)}\sum_{p_1,p_2}\sum_{\substack{m_1\in N\\m_2\in N}}\sum_{\substack{f\in \cF\\p_1\nmid N_f\\p_2\nmid N_f }}w_R(f)\frac{\alpha_f(p_1)^{m_1}+\beta_f(p_1)^{m_1}}{p_1^{m_1/2}} \frac{\alpha_f(p_2)^{m_2}+\beta_f(p_2)^{m_2}}{p_2^{m_2/2}}\notag \\ &\qquad \qquad\qquad\qquad\qquad\qquad\qquad\cdot\frac{\log(p_1)\log (p_2)}{\log^2(R)}\hat\phi_1 \left(m_1 \frac{\log(p_1)}{\log(R)}\right)\hat \phi_2 \left(m_2 \frac{\log(p_2)}{\log(R)}\right),
\end{align}}\normalsize where $W_R(\cF)\:= \ \sum_{f\in \cF} w_R(f).$
Further, we introduce the notations:
\begin{align}\label{not:B'' i,j}
    B''_{r_1,r_2,\cF}(p_1,p_2)&\ = \ \frac{1}{W_R(\cF)}\sum_{f \in \cF,~p_1| N_f,~p_2 | N_f} w_R(f) \lambda_f(p_1)^{r_1}\lambda_f(p_2)^{r_2} \\ \label{not:B' i,j}
     B'_{r_1,r_2,\cF}(p_1,p_2)&\ = \ \frac{1}{W_R(\cF)}\sum_{f \in \cF,~p_1| N_f,~p_2 \nmid N_f} w_R(f) \lambda_f(p_1)^{r_1}\lambda_f(p_2)^{r_2}\\\label{not:B i,j}
     B_{r_1,r_2,\cF}(p_1,p_2)&\ = \ \frac{1}{W_R(\cF)}\sum_{f \in \cF,~p_1\nmid N_f,~p_2 \nmid N_f} w_R(f) \lambda_f(p_1)^{r_1}\lambda_f(p_2)^{r_2}.
\end{align}\normalsize
Whenever there is no confusion, we drop the indication of the family $\cF$, and write  $B''_{r_1,r_2}(p_1,p_2)$, $B'_{r_1,r_2}(p_1,p_2)$, or $B_{r_1,r_2}(p_1,p_2)$. We now compute each of the sums explicitly, using the notations above.

\begin{theorem}\label{thm: S_2.Formula}
    Define $P_0(x)\ = \ \frac{2}{x(x+1)}$, $P_1(x)=-\frac{\sqrt{x}(3 x+1) }{x(x+1)^2}$, $P_2(x)=-\frac{\left(x^2+3 x+1\right) }{x(x+1)^3}$
$P_{m\geq3}(x)=\frac{x^{m / 2}(x-1) }{(x+1)^{m+1}}$. Let $A=\{(1,2),(2,1) \}.$ We have that 
 \begin{align*}
        S_2(\cF,\phi_1,\phi_2)&\ = \  S_{B''}(\cF) + S_{B'}(\cF)+S_{B_f}(\cF)+ S_{B_\infty}(\cF)+O\left(\frac{1}{\log^4(R)}\right),\end{align*} where
 \footnotesize\begin{align}
    S_{B''}(\cF) \ = \  &\sum_{p_1,p_2}\sum_{m_1,m_2=1}^{\infty}  B''_{m_1,m_2}(p_1,p_2) \frac{\log(p_1)\log(p_2)}{p_1^{m_1/2}p_2^{m_2/2}\log^2(R)}\hat\phi_1 \left(m_1 \frac{\log(p_1)}{\log(R)}\right)\hat \phi_2 \left(m_2 \frac{\log(p_2)}{\log(R)}\right),\\
    S_{B'}(\cF)   \ =\  &  \sum_{p_1,p_2}\sum_{m_1=1}^{\infty} B'_{m_1,1}(p_1,p_2) \frac{\log(p_1)\log(p_2)}{p_1^{m_1/2}\sqrt{p_2}\log^2(R)} \left(\sum_{(i,j) \in A}\hat\phi_i\left(\frac{m_1\log(p_1)}{\log(R)}\right)\hat\phi_j\left(\frac{\log(p_2)}{\log(R)}\right)\right)\notag\\&+\sum_{p_1,p_2}\sum_{m_1=1}^{\infty} (B'_{m_1,2}(p_1,p_2)-2B'_{m_1,0}(p_1,p_2)) \frac{\log(p_1)\log(p_2)}{p_1^{m_1/2}p_2\log^2(R)} \left(\sum_{(i,j) \in A}\hat\phi_i\left(m_1\frac{\log(p_1)}{\log(R)}\right)\hat\phi_j\left(2\frac{\log(p_2)}{\log(R)}\right)\right)  \notag\\
    &+\sum_{p_1,p_2}\sum_{m_1=1,m_2=0}^{\infty}B'_{m_1,m_2}(p_1,p_2)\frac{P_{m_2}(p_2)}{p_1^{m_1/2}}\frac{\log(p_1)\log (p_2)}{\log^2(R)}\left(\sum_{(i,j) \in A}\hat\phi_i\left(\frac{m_1\log(p_1)}{\log(R)}\right)\hat\phi_j\left(0\right)\right)
    \end{align}
    \begin{align}
    S_{B_f}(\cF) \ = \  &\sum_{p_1,p_2} B_{1,1}(p_1,p_2)\frac{\log(p_1)\log(p_2)}{\sqrt{p_1}\sqrt{p_2}\log^2(R)} \hat\phi_1\left(\frac{\log(p_1)}{\log(R)}\right)\hat\phi_2\left(\frac{\log(p_2)}{\log(R)}\right)\notag\\
    &+\sum_{p_1,p_2} \left(B_{1,2}(p_1,p_2)-2B_{1,0}(p_1)\right)\frac{\log(p_1)\log(p_2)}{\sqrt{p_1}p_2\log^2(R)} \hat\phi_1\left(\frac{\log(p_1)}{\log(R)}\right)\hat\phi_2\left(\frac{2\log(p_2)}{\log(R)}\right)\notag\\
    &+\sum_{p_1,p_2} \left(B_{2,1}(p_1,p_2)-2B_{0,1}(p_2)\right)\frac{\log(p_1)\log(p_2)}{p_1\sqrt{p_2}\log^2(R)} \hat\phi_1\left(\frac{2\log(p_1)}{\log(R)}\right)\hat\phi_2\left(\frac{\log(p_2)}{\log(R)}\right)\notag\\
    &+\sum_{p_1,p_2} \left(B_{2,2}(p_1,p_2)-2B_{2,0}(p_1)-2B_{0,2}(p_2)+4\right)\frac{\log(p_1)\log(p_2)}{p_1p_2\log^2(R)}\hat\phi_1\left(\frac{2\log(p_1)}{\log(R)}\right)\hat\phi_2\left(\frac{2\log(p_2)}{\log(R)}\right),\\
        S_{B_\infty}(\cF) \ = \  &\hat{\phi}_1(0)\sum_{p_1,p_2}\sum_{m_1=0}^{\infty}\frac{\log(p_1) \log(p_2)}{\log^2(R)}\hat{\phi}_2\left( \frac{2\log(p_2)}{\log(R)}\right)(B_{m_1,2}(p_1,p_2)-2B_{m_1,0}(p_1,p_2))\frac{P_{m_1}(p_1)}{p_2}\notag\\&+\hat{\phi}_1(0)\sum_{p_1,p_1}\sum_{m_1=0}^{\infty}\hat{\phi}_2\left(\frac{\log(p_2)}{\log(R)} \right)\frac{\log(p_1) \log(p_2)}{\log^2(R)}B_{m_1,1}(p_1,p_2)\frac{P_{m_1}(p_1)}{\sqrt{p_2}}\notag\\
&+\hat{\phi}_2(0)\sum_{p_1,p_1}\sum_{m_1=0}^{\infty}\hat{\phi}_1\left(\frac{\log(p_1)}{\log(R)} \right)\frac{\log(p_1) \log(p_2)}{\log^2(R)}B_{1,m_1}(p_1,p_2)\frac{P_{m_1}(p_2)}{\sqrt{p_1}}\notag\\
&+\hat{\phi}_2(0)\sum_{p_1,p_2}\sum_{m_1=0}^{\infty}\frac{\log(p_1) \log(p_2)}{\log^2(R)}\hat{\phi}_1\left( \frac{2\log(p_1)}{\log(R)}\right)(B_{2,m_1}(p_1,p_2)-2B_{0,m_1}(p_1,p_2))\frac{P_{m_1}(p_2)}{p_1} \notag\\   &+\hat{\phi}_1(0)\hat{\phi}_2(0)\sum_{p_1,p_2}\sum_{m_1,m_2=0}^\infty\frac{\log(p_1) \log(p_2)}{\log^2(R)}B_{m_1,m_2}(p_1,p_2)P_{m_1}(p_1)P_{m_2}(p_2),
    \end{align}
 \normalsize    
 with the convention that $B_{0,r}(q,p)=B_{r,0}(p,q) = B'_{0,r}(q,p)=A_r(p)$, $B'_{r,0}(p,q) = A'_r(p),$ and $B''_{0,0}(p,q)=B'_{0,0}(p,q) =B_{0,0}(p,q) =A'_{0}(p)=A_0(p) = 1 $ for every $p$ and $q$ prime.
\end{theorem}
To prove the theorem, we show that the first line of \eqref{eq:S_2 formula} is $S_{B''}(\cF),$ the sum of the second and third line of \eqref{eq:S_2 formula} equals $S_{B'}(\cF)$, and the fourth line of \eqref{eq:S_2 formula} equals $S_{B_f}(\cF) + S_{B_\infty}(\cF)$. Below, we break the proof into three parts, where each part shows one of the equalities mentioned above.
\subsection{\texorpdfstring{First line of \eqref{eq:S_2 formula}: $p_1 | N_f$, $p_2 | N_f$}{Computation of the first line of S2}} This case is simple. Using the notation established above, the first line of \eqref{eq:S_2 formula} is equal to
\begin{align}
   &\frac{1}{W_R(\cF)}\sum_{p_1,p_2}\sum_{\substack{m_1\in N\\m_2\in N}}\sum_{\substack{f\in \cF\\p_1|N_f\\p_2|N_f }}w_R(f)\frac{\lambda_f(p_1)^{m_1}}{p_1^{m_1/2}}  \frac{\lambda_f(p_2)^{m_2}}{p_2^{m_2/2}} \frac{\log(p_1)\log (p_2)}{\log^2(R)}\hat\phi_1 \left(m_1 \frac{\log(p_1)}{\log(R)}\right)\hat \phi_2 \left(m_2 \frac{\log(p_2)}{\log(R)}\right)  \notag \\
    &\ = \  \sum_{p_1,p_2}\sum_{\substack{m_1\in N\\m_2\in N}}  B''_{m_1,m_2}(p_1,p_2) \frac{\log(p_1)\log(p_2)}{p_1^{m_1/2}p_2^{m_2/2}\log^2(R)}\hat\phi_1 \left(m_1 \frac{\log(p_1)}{\log(R)}\right)\hat \phi_2 \left(m_2 \frac{\log(p_2)}{\log(R)}\right) \ = \  S_{B''}(\cF).
\end{align}

\subsection{\texorpdfstring{Sum of Second and Third Lines of \eqref{eq:S_2 formula}:  $p_1 | N_f$, $p_2 \nmid N_f$ and $p_2 | N_f$, $p_1 \nmid N_f$}{Sum of second and third lines of S2}}
We begin by noting that the sums for $p_1 | N_f, p_2 \nmid N_f$ and $p_2 | N_f, p_1 \nmid N_f$ are essential symmetric with the exception of the test function.  Therefore,  for $A=\{(1,2), (2,1)\},$ the second and third line of \eqref{eq:S_2 formula} sum to
 \footnotesize{\begin{align} \label{eq:case 2 sum} 
   \sum_{(i,j) \in A} \left[\frac{1}{W_R(\cF)}\sum_{p_1,p_2}\sum_{\substack{m_1\in N\\m_2\in N}}\sum_{\substack{f\in \cF\\p_1|N_f\\p_2\nmid N_f }}w_R(f)\frac{\lambda_f(p_1)^{m_1}}{p_1^{m_1/2}}\frac{\alpha_f(p_2)^{m_2}+\beta_f(p_2)^{m_2}}{p_2^{m_2/2}} \frac{\log(p_1)\log (p_2)}{\log^2(R)}\hat\phi_i \left(m_1 \frac{\log(p_1)}{\log(R)}\right)\hat \phi_j \left(m_2 \frac{\log(p_2)}{\log(R)}\right)\right].
\end{align}}
\normalsize
For $(i,j)\in A$ fixed, we denote by $(\bigstar)$ the expression inside of the brackets in \eqref{eq:case 2 sum}. By symmetry, it suffices to compute ($\bigstar$) for $(i,j) = (1,2)$.

We start by breaking the $(\bigstar)$ into three parts, depending on whether $m_2 =1, m_2=2, $ or $m_2 \geq 3$.
\small{\begin{align}
    (\bigstar)\,= \, \nonumber\,&\frac{1}{W_R(\cF)}\sum_{p_1,p_2}\sum_{m_1\in \N}\sum_{\substack{f\in \cF\\p_1|N_f\\p_2\nmid N_f }}w_R(f)\frac{\lambda_f(p_1)^{m_1}}{p_1^{m_1/2}}\frac{\lambda_f(p_2)}{\sqrt{p_2}} \frac{\log(p_1)\log (p_2)}{\log^2(R)}\hat\phi_1 \left(m_1 \frac{\log(p_1)}{\log(R)}\right)\hat \phi_2 \left(\frac{\log(p_2)}{\log(R)}\right) \\\nonumber&+  \frac{1}{W_R(\cF)}\sum_{p_1,p_2}\sum_{m_1\in \N}\sum_{\substack{f\in \cF\\p_1|N_f\\p_2\nmid N_f }}w_R(f)\frac{\lambda_f(p_1)^{m_1}}{p_1^{m_1/2}}\frac{(\lambda_f(p_2)^2-2)}{p_2} \frac{\log(p_1)\log (p_2)}{\log^2(R)}\hat\phi_1 \left(m_1 \frac{\log(p_1)}{\log(R)}\right)\hat \phi_2 \left(\frac{2\log(p_2)}{\log(R)}\right)\\\label{eq: case 2 sum m2 geq 3}
    &+\frac{1}{W_R(\cF)}\sum_{p_1,p_2}\sum_{\substack{m_1\in \N\\m_2\geq 3}}\sum_{\substack{f\in \cF\\p_1|N_f\\p_2\nmid N_f }}w_R(f)\frac{\lambda_f(p_1)^{m_1}}{p_1^{m_1/2}}\frac{\alpha_f(p_2)^{m_2}+\beta_f(p_2)^{m_2}}{p_2^{m_2/2}} \frac{\log(p_1)\log (p_2)}{\log^2(R)} \notag \\ &\qquad \qquad \qquad \qquad \qquad \qquad \qquad \qquad \qquad \qquad \cdot \hat\phi_1 \left(m_1 \frac{\log(p_1)}{\log(R)}\right)\hat \phi_2 \left(\frac{m_2\log(p_2)}{\log(R)}\right).
\end{align}} \normalsize

We write $(\bigstar)=(\bigstar')+(\bigstar'')+(\bigstar''')$ where, for example, $(\bigstar')$ is the first line in \eqref{eq: case 2 sum m2 geq 3}. Using the definition of $B'_{r_1,r_2}(p_1,p_2)$ in \eqref{not:B' i,j} and $A_{r}(p)$ in \eqref{def:A_r}, as well as the convention that $A'_{r}(p) = B'_{r,0}(p,q)$ (for any prime $q\nmid N$), we can easily deal with the case when $m_2=1$ and $m_2=2$. We have, 
\begin{equation}\label{eq for them case 2, m1 =1}
    (\bigstar') \ =\  \sum_{p_1,p_2}\sum_{m_1\in \N} B'_{m_1,1}(p_1,p_2) \frac{\log(p_1)\log(p_2)}{p_1^{m_1/2}\sqrt{p_2}\log^2(R)} \hat\phi_1\left(\frac{m_1\log(p_1)}{\log(R)}\right)\hat\phi_2\left(\frac{\log(p_2)}{\log(R)}\right)
\end{equation} 
and
\begin{align}\label{eq for them case 2, m1 =2}
    (\bigstar'') &\ = \  \sum_{p_1,p_2}\sum_{m_1\in \N} \left(B'_{m_1,2}(p_1,p_2)-2B_{m_1,0}(p_1,p_2)\right)\frac{\log(p_1)\log(p_2)}{p_1^{m_1/2}p_2\log^2(R)} \hat\phi_1\left(m_1\frac{\log(p_1)}{\log(R)}\right)\hat\phi_2\left(2\frac{\log(p_2)}{\log(R)}\right).
\end{align}
We now aim to simplify $(\bigstar''')$, which is the case when $m_2 \geq 3$. The purpose of the following lemma is to remove the dependency of $m_2$ in the argument of $\hat\phi_2$ by arguing that we can replace $\hat\phi_2\left(m_2\frac{\log(p_2)}{\log(R)}\right)$ with $\hat\phi_2(0)$ at the cost of $O\left(1/\log^4(R)\right)$.
\begin{lemma} \label{lemma for case 2} Suppose $\phi_1$ and $\phi_2$ are two even Schwartz test functions with Fourier transform supported in $[-\sigma,\sigma]$. Then,
\small\begin{align}
    (\bigstar''')\ &=\ \frac{1}{W_R(\cF)}\sum_{\substack{p_1\\ p_2 < R^{\sigma}}}\sum_{\substack{m_1\in \N\\m_2\geq 3}}\sum_{\substack{f\in \cF\\p_1|N_f\\p_2\nmid N_f }}w_R(f)\frac{\lambda_f(p_1)^{m_1}}{p_1^{m_1/2}}\frac{\alpha_f(p_2)^{m_2}+\beta_f(p_2)^{m_2}}{p_2^{m_2/2}} \frac{\log(p_1)\log (p_2)}{\log^2(R)}\hat\phi_1 \left(m_1 \frac{\log(p_1)}{\log(R)}\right)\hat \phi_2 (0)\notag \\ & \ \qquad+ O\left(\frac{1}{\log^4(R)}\right).
\end{align}\normalsize
\end{lemma}
\begin{proof}
By assumption, $\hat\phi_2$ is even and therefore $\hat\phi'_2(0)=0$. Thus, using the Taylor expansion around the origin, we notice that \begin{equation}\label{eq.lemma, case 2:taylor expansion}
    \hat\phi_2 \left(m_2 \frac{\log(p_2)}{\log(R)}\right) - \hat\phi_2(0)\ =\ O\left(m_2^2 \frac{\log(p_2)^2}{\log(R)^2}\right). 
\end{equation}
Moreover, we know that for every $m \in \N$  and prime $p,$ \begin{equation}\label{eq.lemma, case 2: alpha beta bound}
    |\alpha_f(p)^m + \beta_f(p)^m| \ \leq \ 2 .
\end{equation}     
Thus, because for $p|N_f$, $\lambda_f(p)^{m} = \alpha_f(p)^m + \beta_f(p)^m,$ we have that if $p|N_f$,
\begin{equation}
    |\lambda_f(p)^{m}|\ \leq \ 2.
\end{equation}
Therefore, we get the following estimate:
\footnotesize{\begin{align}\label{case2, lemma first eq}
    &\frac{1}{W_R(\cF)}\sum_{\substack{p_1\\ p_2 < R^{\sigma}}}\sum_{\substack{m_1\in \N\\m_2\geq 3}}\sum_{\substack{f\in \cF\\p_1|N_f\\p_2\nmid N_f }}w_R(f)\frac{\lambda_f(p_1)^{m_1}}{p_1^{m_1/2}}\frac{\alpha_f(p_2)^{m_2}+\beta_f(p_2)^{m_2}}{p_2^{m_2/2}} \frac{\log(p_1)\log (p_2)}{\log^2(R)} \notag \\ & \qquad \qquad \qquad \qquad \qquad \qquad \qquad  \qquad \qquad \qquad \cdot \hat\phi_1 \left(m_1 \frac{\log(p_1)}{\log(R)}\right)\left(\hat \phi_2 \left(m_2\frac{\log(p_2)}{\log(R)}\right)- \phi_2 (0)\right) \notag \\
    &\lesssim \frac{1}{\log^4(R)}\left(\sum_{p_1} \sum_{m_1 \in \N}\frac{\log(p_1)}{p_1^{m_1/2}} \hat\phi_1\left(m_1\frac{\log(p_1)}{\log(R)}\right)\right) \left(\sum_{p_2}\log(p_2)^3 \sum_{m_2 \geq 3}\frac{m_2^2}{p_2^{m_2/2}}\right) \\
   &
   \lesssim \frac{1}{\log^4(R)}\left(\sum_{p_1} \sum_{m_1 \in \N}\frac{\log(p_1)}{p_1^{m_1/2}} \hat\phi_1\left(m_1\frac{\log(p_1)}{\log(R)}\right)\right) \left(\sum_{p_2}\frac{\log(p_2)^3 }{p_2^{3/2}}\right),
\end{align}}\normalsize
where the last line follows using the fact that $m_2^2 \lesssim \left(\frac{3}{2}\right)^{m_2/2}$ and the geometric series formula.
In fact, the proof is complete because $\left(\sum_{p_1} \sum_{m_1 \in \N}\frac{\log(p_1)}{p_1^{m_1/2}} \hat\phi_1\left(m_1\frac{\log(p_1)}{\log(R)}\right)\right)$ and $\left(\sum_{p_2}\frac{\log(p_2)^3 }{p_2^{3/2}}\right)$ are convergent sums.
\end{proof}
With Lemma \ref{lemma for case 2} and using the notation of $M_{3,0}(p)$ from \eqref{def: M_(c,k)}, we have
\small{\begin{align*}
    (\bigstar''') \,=\, \frac{1}{W_R(\cF)}\sum_{p_1,p_2}\sum_{m_1\in \N}\sum_{\substack{f\in \cF\\p_1|N_f\\p_2\nmid N_f }}w_R(f)\frac{\lambda_f(p_1)^{m_1}}{p_1^{m_1/2}}M_{3,0}(p_2)\frac{\log(p_1)\log (p_2)}{\log^2(R)}\hat\phi_1 \left(m_1 \frac{\log(p_1)}{\log(R)}\right)\hat \phi_2 (0) + O\left(\frac{1}{\log^4(R)}\right).
\end{align*}}\normalsize We know from \cite{Miller2009_LOTerms_1LevelDensity}
 that \begin{align*}
     M_{3,0}(p)\ = \ \frac{2}{p(p+1)}-\frac{p^{1/2}(3p+1)}{p(p+1)^2}\lambda_f(p)-\frac{(p^2+3p+1)}{p(p+1)^3}\lambda_f(p)^2+\sum_{m=3}^{\infty}\frac{p^{m/2}(p-1)\lambda_f(p)^m}{(p+1)^{m+1}}.
 \end{align*}
Moreover, by observing that $M_{3,0}(p_2)=\sum_{m_2=0}^\infty P_m(p_2)$, we have that
\begin{align}
\label{eq for case 2 m1 geq 2}
(\bigstar''')\ = \ \sum_{p_1,p_2}\sum_{m_1=1,m_2=0}^{\infty}B'_{m_1,m_2}(p_1,p_2)\frac{P_{m_2}(p_2)}{p_1^{m_1/2}}\frac{\log(p_1)\log (p_2)}{\log^2(R)}\hat\phi_1\left(\frac{m_1\log(p_1)}{\log(R)}\right)\hat\phi_2\left(0\right).
\end{align} 
Now substituting \eqref{eq for them case 2, m1 =1}, \eqref{eq for them case 2, m1 =2}, and \eqref{eq for case 2 m1 geq 2} for $(\bigstar'),(\bigstar'')$, and $(\bigstar''')$ respectively, we get that the second line of \eqref{eq:S_2 formula} equals $S_{B'}(\cF)$.
 
\subsection{\texorpdfstring{Fourth line of \eqref{eq:S_2 formula}: $p_1 \nmid N_f$, $p_2 \nmid N_f$}{Fourth line of S2}}
Lastly, we show that up to $O(\log^{-4} R)$ error, the fourth line of \eqref{eq:S_2 formula} equals $S_{B_f}(\cF) +S_{B_\infty}(\cF) $  .

For fixed primes $p_1,p_2$, let 
\footnotesize\begin{equation}\label{eq: breaking into C(i,j)'s 5.9}
    C(m_1,m_2)\ := \ \sum_{\substack{f\in \cF\\p_1\nmid N_f\\p_2\nmid N_f }}w_R(f)\frac{\alpha_f(p_1)^{m_1}+\beta_f(p_1)^{m_1}}{p_1^{m_1/2}} \frac{\alpha_f(p_2)^{m_2}+\beta_f(p_2)^{m_2}}{p_2^{m_2/2}}\frac{\log(p_1)\log (p_2)}{\log^2(R)}\hat\phi_1 \left(m_1 \frac{\log(p_1)}{\log(R)}\right)\hat \phi_2 \left(m_2 \frac{\log(p_2)}{\log(R)}\right).
\end{equation}\normalsize
The fourth line of \eqref{eq:S_2 formula} is equal to
\footnotesize \begin{align}\nonumber
  &  \frac{1}{W_R(\cF)}\sum_{p_1,p_2}\sum_{\substack{m_1\in N\\m_2\in N}}\sum_{\substack{f\in \cF\\p_1\nmid N_f\\p_2\nmid N_f }}w_R(f)\frac{\alpha_f(p_1)^{m_1}+\beta_f(p_1)^{m_1}}{p_1^{m_1/2}} \frac{\alpha_f(p_2)^{m_2}+\beta_f(p_2)^{m_2}}{p_2^{m_2/2}}\frac{\log(p_1)\log (p_2)}{\log^2(R)}\hat\phi_1 \left(m_1 \frac{\log(p_1)}{\log(R)}\right)\hat \phi_2 \left(m_2 \frac{\log(p_2)}{\log(R)}\right)\notag\\
&\ = \ \frac{1}{W_R(\cF)}\sum_{p_1,p_2}\big(C(1,1)+C(1,2)+C(2,1)+C(2,2)\big) \label{eq:rewrite fourth line finite}
\\ & ~~~~~\ +\frac{1}{W_R(\cF)}\sum_{p_1,p_2}\sum_{m_1,m_2\geq 3} \big(C(m_1,1)+C(m_1,2)+C(1,m_2)+C(2,m_2)+C(m_1,m_2)\big).\label{eq:rewrite fourth line infinite}
\end{align}
\normalsize
We show that the finite part, line \eqref{eq:rewrite fourth line finite}, equals $S_{B_f}(\cF)$ while the infinite part, line \eqref{eq:rewrite fourth line infinite}, equals $S_{B_\infty}(\cF).$

The former immediately follows from the multiplicativity of Hecke eigenvalues. We know that for all prime $p$, $\alpha_f(p) + \beta_f(p) = \lambda_f(p)$ and $\alpha_f(p)^2 + \beta_f(p)^2 = \lambda_f(p)^2-2$. For the latter, we need a lemma. See Appendix \ref{appendix: Proof of Lemma both m1,m2 geq 3} for the proof of it.
\begin{lemma}\label{lemma1: both m_1,m_2 geq 3} Suppose $\phi_1$ and $\phi_2$ are even Schwartz functions with $\hat\phi_1$ and $\hat\phi_2$ having support in $[-\sigma,\sigma]$.
We then have the following estimate: 
\small{\begin{align}\label{lem:phi0.estimate.statment}
    &\frac{1}{W_R(\cF)}\sum_{\substack{p_1,p_2}}\sum_{m_1,m_2\geq 3} C(m_1,m_2)\notag \\
    & \ = \ \frac{1}{W_R(\cF)}\sum_{p_1,p_2<R^\sigma}\sum_{m_1,m_2\geq 3}\sum_{\substack{f\in \cF\\p_1\nmid N_f\\p_2\nmid N_f }}w_R(f)\frac{\sum_{j=1}^2\left(\alpha_f(p_j)^{m_j}+\beta_f(p_j)^{m_j}\right)}{p_1^{m_1/2}p_2^{m_2/2}} \frac{\log(p_1)\log (p_2)}{\log^2(R)}\hat\phi_1(0) \hat \phi_2(0) \notag \\
    & \qquad\qquad + O\left(\frac{1}{\log^4(R)}\right).
\end{align}}
\end{lemma}\normalsize
\begin{lemma}\label{lemma 5.4.2: only one m_i geq 3}
Suppose $\phi_1$ and $\phi_2$ are even Schwartz functions with $\hat\phi_1$ and $\hat\phi_2$ having compact support in $[-\sigma,\sigma]$.
We can estimate every term of \eqref{eq:rewrite fourth line infinite} up to the desired error by replacing $\hat\phi_i\left(\frac{m_i\log(p_i)}{\log(R)}\right)$ by $\hat\phi_i(0)$ and summing over primes $p_i < R^\sigma$. For example, for the first term, we have 
\small\begin{align}
   &\frac{1}{W_R(\cF)}\sum_{\substack{p_1,p_2}}\sum_{m_1\geq 3} C(m_1,1)\notag \\ & \ = \  \frac{1}{W_R(\cF)}\sum_{p_1,p_2<R^\sigma}\sum_{m_1\geq 3}\sum_{\substack{f\in \cF\\p_1\nmid N_f\\p_2\nmid N_f }}w_R(f)\frac{\left(\alpha_f(p_1)^{m_1}+\beta_f(p_1)^{m_1}\right)\lambda_f(p_2)}{p_1^{m_1/2}p_2^{1/2}} \frac{\log(p_1)\log (p_2)}{\log^2(R)}\hat\phi_1(0) \hat \phi_2\left(\frac{\log(p_2)}{\log(R)}\right) \notag \\
    & \qquad\qquad + O\left(\frac{1}{\log^4(R)}\right).
\end{align}\normalsize
\end{lemma}
The proof of Lemma \ref{lemma 5.4.2: only one m_i geq 3} follows from the proof of Lemma \ref{lemma1: both m_1,m_2 geq 3}. \qed

Now using lemmas \ref{lemma1: both m_1,m_2 geq 3} and \ref{lemma 5.4.2: only one m_i geq 3}, as well as using formulas for $M_{3,0}(p)$ from \cite{Miller2009_LOTerms_1LevelDensity}, we get \begin{align*}
&\frac{1}{W_R(\cF)}\sum_{p_1,p_2}\sum_{m_2\geq 3} C(1,m_2) \notag \\
&\ = \  \frac{1}{W_R(\cF)} \sum_{p_1,p_2} \sum_{\substack{f\in \cF\\p_1\nmid N_f\\p_2\nmid N_f }}w_R(f) \frac{\lambda_f(p_1)}{\sqrt{p_1}}M_{3,0}(p_2)\frac{\log(p_1)\log(p_2)}{\log^2(R)} \hat\phi_1\left(\frac{\log(p_1)}{\log(R)}\right)\hat\phi_2(0) + O\left(\frac{1}{\log^4(R)}\right) \\
&\ = \  \hat{\phi}_1(0)\sum_{p_1,p_1}\sum_{m_1=0}^{\infty}\hat{\phi}_2\left(\frac{\log(p_2)}{\log(R)} \right)\frac{\log(p_1) \log(p_2)}{\log^2(R)}B_{m_1,1}(p_1,p_2)\frac{P_{m_1}(p_1)}{\sqrt{p_2}} + O\left(\frac{1}{\log^4(R)}\right).
\end{align*}\normalsize
Performing a similar substitution, we have 
\footnotesize\begin{eqnarray}
& &\sum_{p_1,p_2}\sum_{m_2\geq 3} \frac{C(2,m_2)}{W_R(\cF)}\nonumber\\  & &= \  \hat{\phi}_1(0)\sum_{p_1,p_2<R^\sigma}\sum_{m_1=0}^{\infty}\frac{\log(p_1) \log(p_2)}{\log^2(R)}\hat{\phi}_2\left( \frac{2\log(p_2)}{\log(R)}\right)(B_{m_1,2}(p_1,p_2)-2B_{m_1,0}(p_1,p_2))\frac{P_{m_1}(p_1)}{p_2}
+ O\left(\frac{1}{\log^4(R)}\right), \notag
\\
& &\sum_{p_1,p_2}\sum_{m_1\geq 3} \frac{C(m_1,1)}{W_R(\cF)}\nonumber\\  &  &=  \  \hat{\phi}_2(0)\sum_{p_1,p_1<R^\sigma}\sum_{m_1=0}^{\infty}\hat{\phi}_1\left(\frac{\log(p_1)}{\log(R)} \right)\frac{\log(p_1) \log(p_2)}{\log^2(R)}B_{m_1,1}(p_1,p_2)\frac{P_{m_1}(p_2)}{\sqrt{p_2}}+ O\left(\frac{1}{\log^4(R)}\right),\nonumber\\
& &\sum_{p_1,p_2}\sum_{m_1\geq 3} \frac{C(m_1,2)}{W_R(\cF)} \nonumber\\ & & =\   \hat{\phi}_2(0)\sum_{p_1,p_2<R^\sigma}\sum_{m_1=0}^{\infty}\frac{\log(p_1) \log(p_2)}{\log^2(R)}\hat{\phi}_1\left( \frac{2\log(p_1)}{\log(R)}\right)(B_{m_1,2}(p_1,p_2)-2B_{m_1,0}(p_1,p_2))\frac{P_{m_1}(p_2)}{p_2}
+ O\left(\frac{1}{\log^4(R)}\right),\nonumber\\
& &\sum_{p_1,p_2}\sum_{m_1,m_2\geq 3} \frac{C(m_1,m_2)}{W_R(\cF)} \nonumber\\
& & = \  \hat{\phi}_1(0)\hat{\phi}_2(0)\sum_{p_1,p_2<R^\sigma}\sum_{m_1,m_2=0}^\infty\frac{\log(p_1) \log(p_2)}{\log^2(R)}B_{m_1,m_2}(p_1,p_2)P_{m_1}(p_1)P_{m_2}(p_2) + O\left(\frac{1}{\log^4(R)}\right).
\end{eqnarray}\normalsize
With this, we can see that $S_{B\infty}(\cF)$ is equivalent to \eqref{eq:rewrite fourth line infinite}, finishing the proof of Theorem \ref{thm: S_2.Formula}.


%%%%%%%%%%%%%%%%%%%%%%%%%%%%%%%%%%%%%%%%%%%%%%%
%%%%%%%%%%%%%%%%%%%%%%%%%%%%%%%%%%%%%%%%%%%%%%%
\vspace{1em}

\section{Formulas for Family Specific Terms} \label{chp: sum of weights and family specific terms}

We use the harmonic weights\footnote{The harmonic weights are nearly constant across the family: by \cite{I1}, \cite{HL},\cite{Miller2009_LOTerms_1LevelDensity} we have $w_R(f)\ =\ N^{\epsilon}$ uniformly for $f$ in the family. If we allow ineffective constants, the factor $N^{\varepsilon}$ can be replaced by $\log(N)$ for sufficiently large $N$.There are other ways of normalizations. \cite{Miller2009_LOTerms_1LevelDensity} and our paper averages with 
$w_R(f)\ =\ \frac{\zeta_N(2)}{Z(1,f)} = \frac{\zeta(2)}{L(1,\mathrm{sym}^2 f)}$. Rouymi \cite{Rouymi2011TraceNonAnnulation} uses the Petersson-harmonic weight $w_R'(f)=\frac{\Gamma(k-1)}{(4\pi)^{k-1}}\langle f,f\rangle_N^{-1}$.
Rankin-Selberg shows $\langle f,f\rangle_N \asymp_{k,N} L(1,\mathrm{sym}^2 f)$, so $w_R(f)\propto_{k,N} w_R'(f)$.
In addition, using this weight, Rouymi derived formulas for  level $N = p^{\alpha}$ for fixed prime $p$ and $\alpha \to \infty$. We also note that different weights exhibit different behaviors as seen in \cite{KnightlyReno2019} and \cite{dillon2025centered}.} to facilitate computing explicitly the asymptotic. Recall we are working with the weights 
$$w_R(f)\ =\ \frac{Z_N(1,f)}{Z(1,f)}$$ where $Z_N(s,f)$ and $Z(s,f)$ are defined in \eqref{def:Z_N_func} and \eqref{def:Z_func} respectively. We consider $\cF = H_k^*(N)$ the family of holomorphic cusp newforms with weight $k$ and level $N$. We consider four different scenarios in which the level approaches infinity. 
\begin{itemize}
    \item $N$ is prime. \\
    \item $N\ =\ q_1q_2$ for two primes $q_1,q_2$, with $q_1 \neq q_2$ and $q_1$ fixed. \\
    \item $N\ =\ q_1q_2$ for two primes with $q_1 \sim N^{\delta}$ and $q_2 \sim N^{1-\delta}$ where $\delta \in (0,1/2].$ \\
    \item $N\ =\ p^2$ where $p$ is prime
\end{itemize}
The goal of this section is to compute explicitly the terms $A'_r(p), A_r(p), B''_{r_1,r_2}(p_1,p_2),B'_{r_1,r_2}(p_1,p_2),$ and $B_{r_1,r_2}(p_1,p_2)$, which determine the value of $S_1(\cF,\phi_1,\phi_2)$ and $S_2(\cF,\phi_1,\phi_2)$. We will find we only see significant differences in the moment when $q_1$ is fixed and $q_2 \to \infty.$

Recall that using the definition of the above terms involve the sum of the weights over the family $W_R(\cF)\ := \ \sum_{f\in \cF} w_R(f).$ Therefore, we start by computing $W_R(\cF)$ for each scenario.

\subsection{Sum of the Weights}
We compute the sum of the weights for each scenario below. In brief, we do so by first appealing the Peterson trace formula and resulting bounds in \cite{ILS2} in the $N$ square-free case and to proposition 4.1 in \cite{BarrettEtAl2016arXiv} for the $N = p^2$ case. Below, Lemma \ref{lem: sum of weights squarefree} computes $W_R(\cF)$ for the first three square-free cases while Lemma \ref{lem: sum of weights for N=p2} computes it for $N\ =\ p^2$ case.
\begin{lemma}\label{lem: sum of weights squarefree}
    Let $N\ =\ q_1q_2$ be square-free with $q_1,q_2$ being  distinct primes with $$q_1\sim C_1N^{\delta},q_2\sim C_2N^{1-\delta}.$$ Here, $\delta \in [0,1/2]$ while $C_1,C_2$ can be both prime constants if $\delta\neq0$, or $C_1=1$ or prime  and $C_2$ prime if $\delta \ =\ 0$. Then,
  \begin{align}
     W_R(\cF):= \sum_{f \in \mathcal{F}}w_r(f)=\frac{k-1}{12}\varphi(N)+O(E(N, \delta)),
  \end{align}
    where
    \[
    E(N,\delta)\ =\  
    \begin{cases}
    \frac{\log(N)}{N^{\frac{1-\delta}{2}}} & \text{if } 0< \delta \leq \frac{1}{2} \text{ or } \delta\ =\ 0, C_1=1 \\[1.5ex]
    \frac{\log^2(N)}{N^{1/2}} & \text{if } \delta \ =\  0 \text{ and } C_1 \text{ prime}.
    \end{cases}
    \]\label{lem:sum of weights}
\end{lemma}
\textit{Proof.} We begin by noticing that $W_R(\cF) = \Delta_{k,N}^*(1,1)$, which by Proposition \ref{prop:ILS_Lem_2.7} (\cite{ILS2} Lemma 2.7) becomes
\begin{align}
W_R(\mathcal{F})\ &=\ \Delta^{*}_{k,N}(1,1)\ =\ \frac{k-1}{12}\sum_{LM=N}\mu(L)M\sum_{\ell|L^{\infty}}\ell^{-1}\Delta_{k,M}(\ell^2,1).
\end{align}
Notice that if $C_1\ =\ 1$, then $N$ is a prime and we get that $$W_R(\cF)\ =\ \frac{k-1}{12}\left[N\Delta_{k,N}(1,1)-\sum_{\ell|N^{\infty}}\ell^{-1}\Delta_{k,q_1}(\ell^2,1)\right].$$ On the other hand, if $C_1 \neq 1$, $N$ is product of 2 distinct primes; so, we have: 
$$W_R(\cF)\ =\ \frac{k-1}{12}\left[N\Delta_{k,N}(1,1)-q_1\sum_{\ell|q_2^{\infty}}\ell^{-1}\Delta_{k,q_1}(\ell^2,1)-q_2\sum_{\ell|q_1^{\infty}}\ell^{-1}\Delta_{k,q_2}(\ell^2,1)+\sum_{\ell|N^{\infty}}\ell^{-1}\Delta_{k,1}(\ell^2,1) \right].$$
Using Proposition \ref{prop:peter aprox. ILS2.2} (\cite{ILS2} Cor 2.2), we have that
$$\Delta_{k,N}(1,1)\ =\ 1+O_k\left( \frac{1}{N^{3/2}} \right) ~~\text{ and } ~~\Delta_{k,M}(\ell^2,1)\ =\ \delta(\ell^2,1)+O_k\left( \frac{\ell\log(2\ell^2)}{M(\ell+kM)^{1/2}} \right),$$ where the implied constant is absolute in $k.$
We first evaluate the two sums that appear in both the $N$ prime and $N$ being the product of two distinct primes cases, then move to the sums specific to the two prime case.\\
\\
\textbf{Case 1:} $N\Delta_{k,N}(1,1)$\\
We see that since $$\Delta_{k,N}(1,1)\ =\ 1+O_k\left( \frac{1}{N^{3/2}} \right),$$
we have $$N\Delta_{k,N}(1,1)\ =\ N+O_k\left(\frac{1}{N^{1/2}}\right)$$
for the first case.\\
\\
\textbf{Case 2:} $\sum_{\ell|N^{\infty}}\ell^{-1}\Delta_{k,q_1}(\ell^2,1)$\\
We have that 
\begin{align}
\sum_{\ell|N^{\infty}}\ell^{-1}\Delta_{k,1}(\ell^2,1) &= \sum_{r_1,r_2=0}^{\infty}q_1^{-r_1}q_2^{-r_2}\Delta_{k,1}(q_1^{2r_1}q_2^{2r_2},1) \notag \\
&=\sum_{r_1,r_2=0}^{\infty}\left[q_1^{-r_1}q_2^{-r_2}\delta(q_1^{2r_1}q_2^{2r_2},1)+O_k\left( \frac{\log(2q_1^{2r_1}q_2^{2r_2})}{(q_1^{r_1}q_2^{r_2}+k)^{1/2}} \right)\right].
\end{align}
As $q_1^{-r_1}q_2^{-r_2}\delta(q_1^{2r_1}q_2^{2r_2},1)=1$ if and only if $r_1=r_2=0$ and is 0 otherwise $$\sum_{r_1,r_2=0}^{\infty}q_1^{-r_1}q_2^{-r_2}\delta(q_1^{2r_1}q_2^{2r_2},1)=1.$$
Since the implied constant is absolute and $k$ is constant, we can combine the sum of the big O terms into one error, and obtain
$$1+O_k \left(\sum_{r_1,r_2=0}^{\infty}  \frac{\log(2q_1^{2r_1}q_2^{2r_2})}{(q_1^{r_1}q_2^{r_2}+k)^{1/2}}\right).$$
We now see
\begin{align*}
     \sum_{r_1,r_2=1}^\infty \frac{\log(2q_1^{r_1}q_2^{r_2})}{\sqrt{q_1^{r_1}q_2^{r_2}+k}} \ &\lesssim_k \frac{\log(q_1)}{\sqrt{q_2}}\left( \sum_{r_1=1}^\infty \frac{r_1}{\sqrt{q_1^{r_1}}}\right)+\frac{\log(q_2)}{\sqrt{q_1}}\left(\sum_{r_2=1}^\infty\frac{r_2}{\sqrt{q_2^{r_2}}}\right)\lesssim_k \frac{\log(N)}{\sqrt{N}}.
\end{align*}
Thus we have $$1+O_k\left(\frac{\log(N)}{N^{1/2}} \right).$$\\
We have now exhausted the cases for $N$ prime, giving us an error term of $O_k\left( \frac{\log(N)}{N^{1/2}}\right).$ We now operate under the assumption that if $\delta=0$, then $C_1$ is prime.\\
\\
\textbf{Case 3:} $q_1\sum_{\ell|q_2^{\infty}}\ell^{-1}\Delta_{k,q_1}(\ell^2,1)$\\
Similarly, 
\begin{align*}
q_1\sum_{\ell|q_2^{\infty}}\ell^{-1}\Delta_{k,q_1}(\ell^2,1)&= \ q_1\sum_{m=0}^{\infty}q_2^{-m}\Delta_{k,q_1}(q_2^{2m},1)\\
    &= \ q_1\sum_{m=0}^{\infty}q_2^{-m}\left(\delta(q_2^{2m},1)+O_k\left( \frac{q_2^{m}\log(2q_2^{2m})}{q_1(q_2^{m}+kq_1)^{1/2}} \right) \right) \\
    &= \ q_1+O_k\left(\sum_{m=0}^{\infty} \frac{\log(2q_2^{2m})}{(q_2^{m}+kq_1)^{1/2}} \right)
\end{align*}
for the same reasons as in Case 2. We now turn our attention to the error term
\begin{align*}
    \sum_{m=0}^{\infty} \frac{\log(2q_2^{2m})}{(q_2^{m}+kq_1)^{1/2}} \lesssim_k \log(q_2)\sum_{m=0}^{\infty} \frac{m}{q_2^{m/2}}
    \lesssim_k\frac{\log(q_2)}{q_2^{1/2}}\lesssim_k\frac{\log   N}{N^{(1-\delta)/2}}.
\end{align*}
Thus we have
$$q_1\sum_{\ell|q_2^{\infty}}\ell^{-1}\Delta_{k,q_1}(\ell^2,1)\ =\ q_1+O\left(\frac{\log(N)}{N^{(1-\delta)/2}}\right).$$\\ 
\\
\textbf{Case 4:} 
$q_2\sum_{\ell|q_1^{\infty}}\ell^{-1}\Delta_{k,q_2}(\ell^2,1)$\\ 
Latly, we have
\begin{align*}
    q_2\sum_{\ell|q_1^{\infty}}\ell^{-1}\Delta_{k,q_2}(\ell^2,1)&=\ q_2\sum_{m=0}^{\infty}q_1^{-m}\Delta_{k,q_2}(q_1^{2m},1)\\
    &=\ q_2\sum_{m=0}^{\infty}q_1^{-m}\left(\delta(q_1^{2m},1)+O_k\left( \frac{q_1^{m}\log(2q_1^{2m})}{q_2(q_1^{m}+kq_2)^{1/2}} \right) \right) \\
    &=\ q_2+O_k\left(\sum_{m=0}^{\infty} \frac{\log(2q_1^{2m})}{(q_1^{m}+kq_2)^{1/2}} \right).
\end{align*}
We note that this error term requires additional care as when $\delta=0$, the method used in Case 3 achieves an $O(1)$ error. Thus we write
$$\sum_{m=0}^{\infty} \frac{\log(2q_1^{2m})}{(q_1^{m}+kq_2)^{1/2}}\lesssim_k\frac{\log(q_1)}{q_2^{1/2}}\left[\sum_{m\leq \frac{\log(kq_2)}{\log(q_1)}} \frac{m}{(\frac{q_1^{m}}{kq_2}+1)^{1/2}}+\sum_{m> \frac{\log(kq_2)}{\log(q_1)}} \frac{m}{(\frac{q_1^{m}}{kq_2}+1)^{1/2}} \right].$$
For the first sum, since $q_1^{m}/kq_2 \leq 1$, we have
\begin{align*}
    \sum_{m\leq \frac{\log(kq_2)}{\log(q_1)}} \frac{m}{(\frac{q_1^{m}}{kq_2}+1)^{1/2}} &\lesssim_k \sum_{m\leq \frac{\log(kq_2)}{\log(q_1)}} m\lesssim_k \frac{\log^2 (kq_2)}{\log^2 (q_1)}.
\end{align*}
The second sum requires more care. Define $\alpha= \lceil\log(kq_2)/\log(q_1)\rceil$. We can rewrite the sum as
\begin{align}
    \sum_{m \geq\alpha} \frac{m}{(\frac{q_1^{m}}{kq_2}+1)^{1/2}} &=\sum_{m=0}^\infty(m+\alpha)\frac{1}{(\frac{q_1^{m+\alpha}}{kq_2}+1)^{1/2}} \lesssim\ \sum_{m=0}^\infty(m+\alpha)\frac{1}{(q_1^{m})^{1/2}}\lesssim \frac{\log(kq_2)}{\log(q_1)}.
\end{align}
Combining these together yields 
\small\begin{align}
  \frac{\log(q_1)}{q_2^{1/2}}\left[\sum_{m\leq \frac{\log(kq_2)}{\log(q_1)}} \frac{m}{(\frac{q_1^{m}}{kq_2}+1)^{1/2}}+\sum_{m> \frac{\log(kq_2)}{\log(q_1)}} \frac{m}{(\frac{q_1^{m}}{kq_2}+1)^{1/2}} \right] &\lesssim  \frac{\log(q_1)}{q_2^{1/2}}\left[\frac{\log^2( kq_2)}{\log^2 (q_1)}+\frac{\log(kq_2)}{\log(q_1)} \right]\notag \\
  &\lesssim \frac{\log^2(C_2N^{1-\delta})}{(C_2N^{1-\delta})^{1/2}\log^2 (C_1N^{\delta})}+\frac{\log (C_2N^{1-\delta})}{(C_2N^{1-\delta})^{1/2}}.
\end{align}\normalsize
Here we note that if $\delta=0$, left term dominates, otherwise the right term dominates, thus 
$$\sum_{m=0}^{\infty} \frac{\log(2q_1^{2m})}{(q_1^{m}+kq_2)^{1/2}}\ =\ O(E(N,\delta))$$
  \[
    E(N,\delta)\ =\ 
    \begin{cases}
    \frac{\log(N)}{N^{\frac{1-\delta}{2}}} & \text{if } 0< \delta \leq \frac{1}{2} \\[1.5ex]
    \frac{\log^2(N)}{N^{1/2}} & \text{if } \delta = 0.
    \end{cases}
    \]
Combining this with the other error terms and observing that $\varphi(N)=\sum_{d|N}\mu(d)N/d$ yields the lemma. \qed
\\
While the Peterson trace formula applies for square-free level, we wish to investigate beyond this condition.  Thus we appeal to proposition 4.1 in \cite{BarrettEtAl2016arXiv} to obtain the following lemma.
\begin{lemma} \label{lem: sum of weights for N=p2}
    Let $N\ =\ p^2$ for some prime $p$, then
    \begin{align}
        W_R(\mathcal{F}) \ = \ \sum_{f\in\mathcal{H}_k^*(N)}w_R(f) \ = \ \frac{k-1}{12}(p^2-p-1) \ + \ O\left(\frac{1}{N^{1/4}}\right) \ = \ \frac{k-1}{12}(\varphi(N)-1) \ + \ O\left(\frac{1}{N^{1/4}}\right).
    \end{align}
\end{lemma}
    We have \begin{align} \label{weights_N=p^2}
     W_R(\mathcal{F})= \ \Delta^*_{k,N}(1,1).
\end{align}
By applying Proposition \ref{prop:barrett}, we have
\begin{align*}
     W_R(\mathcal{F})\ &= \ \frac{k-1}{12}\left(\mu(1)p^2\left(\frac{p^2-1}{p^2}\right)\Delta_{k,~p^2}(1,1) \ + \ \mu(p)p\Delta_{k,~p}(1,1) \ + \ \mu(p^2)\sum_{\substack{\ell \mid (p^2)^\infty \\ (\ell, 1)=1}}\ell^{-1}\Delta_{k,~1}(\ell^2,1) \right) \\
    &= \ \frac{k-1}{12}((p^2-1)\Delta_{k,~p^2}(1,1) \ - \ p\Delta_{k,~p}(1,1)).
\end{align*}
We now use \cite{ILS2} Proposition 2.2 to compute $\Delta_{k,~p}(1,1)$ and $\Delta_{k,~p^2}(1,1)$:
\begin{align*}
    \Delta_{k,~p}(1,1) \ &= \ 1 \ + \ O\left( \frac{\tau(p)}{pk^{5/6}}\left(\frac{1}{1+kp} \right)^{1/2}\log(2)\right)= 1 + O_k\left(\frac{1}{p^{3/2}}\right).
\end{align*}
Similarly, 
\begin{align*}
       \Delta_{k,~p^2}(1,1) \ &= \ 1 \ + \ O\left( \frac{\tau(p^2)}{p^2k^{5/6}}\left(\frac{1}{1+kp^2} \right)^{1/2}\log(2)\right) = 1 + O_k\left(\frac{1}{p^3}\right).
\end{align*}
Thus, we have
\begin{align*}
    W_R(\mathcal{F}) \ = \ \frac{k-1}{12}(p^2-p-1) \ + \ O_k\left( \frac{1}{N^{1/4}} \right).
\end{align*}
\qed 
\subsection{Computing Terms}
Recall that  our extended explicit formula is expressed in terms of weighted average moments of Hecke eigenvalues, namely 
$$\sum_{f \in H^{*}_k(N)}\lambda(p_1)^{r_1}\lambda(p_2)^{r_2},$$
wherein restrictions are placed on $p_1,p_2$ relative to the level $N$. We wish to compute these average weighted moments for our four cases: $N$ prime, $N=q_1q_2$ with $q_1$ fixed and distinct from $q_2$, $N=q_1q_2$ with both distinct and neither fixed, and lastly $N=p^2$. We summarize our results in the following four lemmas. 
\begin{remark}
    Note that unless otherwise stated, the value of the moments in the following lemmas is 0.
\end{remark}

\begin{lemma}\label{thm: N prime, weight term values}
    For $\cF= H_k^*(N)$ where $N$ is a prime, the following holds:
    \footnotesize \begin{align*}
A'_r(p)\ &=\ O(N^{-r/2}) \\[6pt]
A_r(p)\ &= \
\begin{cases}
  C_{r/2} + O\!\left(\frac{r 2^r p^{r/4}\log(p^rN)}{N}\right) 
    & \text{if $r$ even}, \\
  O\!\left(\frac{r 2^r p^{r/4}\log(p^rN)}{N}\right) 
    & \text{if $r$ odd}.
\end{cases} \\[10pt]
B_{r_1,r_2}''(p_1,p_2)\ &=\ O\!\left(N^{-(r_1+r_2)/2}\right) \\[6pt]
B'_{r_1,r_2}(p_1,p_2)\  &=\ O\!\left(2^{r_2}N^{-r_1/2}\right) \\[6pt]
B_{r_1,r_2}(p_1,p_2)\ &=\
\begin{cases}
  C_{(r_1+r_2)/2} 
    + O\!\left(2^{r_1+r_2}(p_1^{r_1}p_2^{r_2})^{1/4}
    \frac{\log(2p_1^{r_1}p_2^{r_2}N)\log(N)}{k^{5/6}N}\right) 
    & \text{if $p_1 = p_2$ and $(r_1-r_2 \bmod 2)=0$}, \\[6pt]
  C_{r_1/2}C_{r_2/2} 
    + O\!\left(2^{r_1+r_2}(p_1^{r_1}p_2^{r_2})^{1/4}
    \frac{\log(2p_1^{r_1}p_2^{r_2}N)\log(N)}{k^{5/6}N}\right) 
    & \text{if $p_1 \neq p_2$ and $r_1,r_2$ even}, \\[6pt]
  O\!\left(2^{r_1+r_2}(p_1^{r_1}p_2^{r_2})^{1/4}
    \frac{\log(2p_1^{r_1}p_2^{r_2}N)\log(N)}{k^{5/6}N}\right) 
    & \text{otherwise}.
\end{cases}
\end{align*}\normalsize
\end{lemma}
\noindent\textit{Proof of Lemma \ref{thm: N prime, weight term values}.} First, we show asymptotics for $A_r(p)$ and $A'_r(p).$  We know that if $p| N,$ then $p\ =\ N.$ Therefore, we first determine how $A'_r(N)$ behaves. We know that $\lambda_f(N)^2 \ =\ \frac{1}{N}. $ Therefore, we have 
\begin{equation}
    A'_r(N) \ = \ O(N^{-r/2}).
\end{equation}
Moreover, we compute $A_r(p)$ for $p\neq N.$ To do this, we use  \eqref{eq: p not dividing N, hecke eigenvalue val} and \eqref{eq: average two hecke}: \begin{align}
    A_r(p)\ &=\ \frac{1}{W_R(\cF)} \sum_{f\in\cF} w_R(f) \lambda_f(p)^{r} \notag \\
    \ & =\ \frac{1}{W_R(\cF)} \sum_{f\in\cF} w_R(f) \sum_{k=0}^{\lfloor r/2\rfloor }b_{r,r-2k}\lambda_f(p^{r-2k}) \notag \\
   \ &=\  \sum_{\ell=0}^{\lfloor r/2\rfloor }b_{r,r-2\ell}\left(\delta\left(p^{r-2\ell},1\right)+O\left(p^\frac{r-2\ell}{4} \frac{\log(2p^{r-2\ell}N)}{N}\right)\right) \notag \\
   \ &=\ \begin{cases}
  C_{r/2} + O\!\left(\frac{r 2^r p^{r/4}\log(p^rN)}{N}\right) 
    & \text{if $r$ even}, \\[2pt]
  O\!\left(\frac{r 2^r p^{r/4}\log(p^rN)}{N}\right) 
    & \text{if $r$ odd}.
\end{cases}
\end{align} where $C_{r/2}$ is the $(r/2)^{\text{th}}$ Catalan number (see sequences in A000108, A000984 in \cite{Sl}). For the last equality, we use the fact that $b_{r,r-2\ell}\ \leq \ 2^r.$We compute $B''_{r_1,r_2,H_k^*(N)}(p_1,p_2),\ B'_{r_1,r_2, H_k^*(N)}(p_1,p_2),$ and $B_{r_1,r_2,H_k^*(N)}(p_1,p_2)$ for various $r_1,r_2,p_1,$ and $p_2.$  Because $N$ is prime, notice that if prime $p$ divides $N$, then $p\ =\ N$. First, notice that \begin{equation}
    B''_{r_1,r_2,\cH_k^*(N)}(N,N)\ \lesssim \ \frac{1}{N^{(r_1+r_2)/2}}, 
\end{equation}because $N$ is prime, thus square-free, and in this case, for all prime $p|N,$ $\lambda_f(p)^2\ =\ \frac{1}{p}.$

Secondly, for $p_2 \nmid N$, we have 
\begin{align}
     B'_{r_1,r_2,H_k^*(N)}(N,p_2) \ &=\ \frac{1}{W_R(\cF)}\sum_{f\in \cF} w_R(f) \lambda_f(N)^{r_1}\lambda_f(p_2)^{r_2} \notag\\ &\lesssim \frac{1}{N^{r_1/2}}\left(\frac{1}{W_R(\cF)}\sum_{f\in \cF} w_R(f) \lambda_f(p_2)^{r_2}\right) \lesssim \frac{2^{r_2}}{N^{r_1/2}}.
\end{align} 
Recall for $(p_1,p_2,N) \ =\ 1$, by equation \ref{eq:harm avg aprox} we have 
$$ \sum_{f\in H_k^*(N)} w_R(f) \lambda(m)\lambda(n)\ =\  \frac{k-1}{12} \varphi(N) \delta(m, n)  +O\left(k^{1 / 6}(m n)^{1 / 4}(n, N)^{-1 / 2} \tau^2(N) \tau_3((m, n)) \log (2 m n N)\right).$$ 
Therefore using \eqref{lem: sum of weights squarefree}, we find
\begin{equation} \label{eq: average two hecke}
    \frac{1}{W_R(\cF)}\sum_{f\in H_k^*(N)} w_R(f) \lambda(m)\lambda(n) \ = \ \delta(m,n) + O\left((mn)^{1/4}\frac{\log(2mnN)}{k^{5/6}N}\right).
\end{equation}
In addition, note that \begin{align} \label{eq: p not dividing N, hecke eigenvalue val}   \lambda_f(p)^r\ =\ \sum_{k=0}^{\lfloor r/ 2\rfloor} b_{r, r-2 k} \lambda_f\left(p^{r-2 k}\right). 
\end{align}We can compute $b_{r,r-2k}$ recursively using the fact that $\lambda_f(m)\lambda_f(n) \ =\ \sum_{d|(m,n)}\lambda_f\left(\frac{mn}{d^2}\right)$.

Thus, we have that for $p_1 \nmid N$ and $p_2 \nmid N:$ 
\footnotesize\begin{align}
   & B_{r_1,r_2,H_k^*(N)}(p_1,p_2)\\
 \notag \  &=\ \sum_{k_1=0}^{\lfloor r_1/2\rfloor}\sum_{k_2=0}^{\lfloor r_2/2\rfloor} b_{r_1,r_1-2k_1}b_{r_2,r_2-2k_2} \left[\delta(p_1^{r_1-2k_1},p_2^{r_2-2k_2}) + O\left((p_1^{r_1}p_2^{r_2})^{1/4}\frac{\log(2p_1^{r_1}p_2^{r_2}N)}{k^{5/6}N}\right)\right] \notag \\
    &=\ \sum_{k_1=0}^{\lfloor r_1/2\rfloor}\sum_{k_2=0}^{\lfloor r_2/2\rfloor} b_{r_1,r_1-2k_1}b_{r_2,r_2-2k_2} \delta(p_1^{r_1-2k_1},p_2^{r_2-2k_2}) + O\left((p_1^{r_1}p_2^{r_2})^{1/4}\frac{\log(2p_1^{r_1}p_2^{r_2}N)}{k^{5/6}N}\right) \notag \\
    &=\ \begin{cases}
  O\!\left(2^{r_1+r_2}(p_1^{r_1}p_2^{r_2})^{1/4} 
    \frac{\log(2p_1^{r_1}p_2^{r_2}N)}{k^{5/6}N}\right) 
    & \text{if $p_1\ =\ p_2$ and $(r_1 - r_2 \bmod 2)\ =\ 1$} \\[2pt]
  %\sum_{k=0}^{\min(r_1,r_2)} b_{r_1,k} b_{r_2,k} 
    C_{(r_1+r_2)/2} 
    + O\!\left(2^{r_1+r_2}(p_1^{r_1}p_2^{r_2})^{1/4} 
    \frac{\log(2p_1^{r_1}p_2^{r_2}N)}{k^{5/6}N}\right) 
    & \text{if $p_1\  =\ p_2$ and $(r_1 - r_2 \bmod 2)\ = \ 0$} \\[2pt]
  O\!\left(2^{r_1+r_2}(p_1^{r_1}p_2^{r_2})^{1/4} 
    \frac{\log(2p_1^{r_1}p_2^{r_2}N)}{k^{5/6}N}\right) 
    & \text{if $p_1 \neq p_2$ and $(r_1 - r_2 \bmod 2)\ =\ 1$ 
      or $p_1 \neq p_2$ and $r_1,r_2$ odd} \\[2pt]
  C_{r_1/2} C_{r_2/2} 
    + O\!\left(2^{r_1+r_2}(p_1^{r_1}p_2^{r_2})^{1/4} 
    \frac{\log(2p_1^{r_1}p_2^{r_2}N)}{k^{5/6}N}\right) 
    & \text{if $p_1 \neq p_2$ and $r_1,r_2$ both even.}
\end{cases}
\end{align}\normalsize where $C_{r}$ is the $r^{\text{th}}$ Catalan number. 
The second-to-last equality follows because the sums here are finite. \qed

Next, we compute the terms $A'_{r,\cF}(p), \ A_{r,\cF}(p), B''_{r_1,r_2,\cF}(p_1,p_2), \ B'_{r_1,r_2,\cF}(p_1,p_2), \ $ and $B_{r_1,r_2,\cF}(p_1,p_2)$ for the other three cases we consider. The proof of them is similar to that of Lemma \ref{thm: N prime, weight term values} and can be found in Appendix \ref{appendix: Proof of Lemmas for computing As and Bs}.

\begin{lemma}\label{lem: 2 primes 1 fixed}
    Suppose $\cF = H_k^*(N)$ where $N = q_1\cdot q_2$ for fixed prime $q_1$ and $q_2 \to \infty.$ Then
    \footnotesize\begin{align*}
A'_{r,\cF}(p)\ &= \
\begin{cases}
  q_1^{-r/2} 
    & \text{if $p\ =\ q_1$ and $r$ even}, \\[6pt]
  O\!\left(\tfrac{\log(N)}{N}\right) 
    & \text{if $p\ =\ q_1$ and $r$ odd}, \\[6pt]
  O\!\left(N^{-r/2}\right) 
    & \text{if $p\ =\ q_2$}.
\end{cases} \\[10pt]
A_{r,\cF}(p)\ &= \
\begin{cases}
  C_{r/2} + O_k\!\left(\tfrac{2^r p^{r/4}\log(p^rN)\log(N)}{N}\right) 
    & \text{if $r$ even}, \\[6pt]
  O_k\!\left(\tfrac{2^r p^{r/4}\log(p^rN)\log(N)}{N}\right) 
    & \text{if $r$ odd}.
\end{cases}\\
B''_{r_1,r_2,\cF}(p_1,p_2)\ &= \
\begin{cases}
  q_1^{-(r_1+r_2)/2} 
    & \text{if $p_1\ =\ p_2\ =\ q_1$, $r_1\  \equiv \ r_2 \pmod{2}$}, \\[6pt]
  O\!\left(\tfrac{\log^2(N)}{N}\right) 
    & \text{if $p_1\ =\ p_2\ =\ q_1$, $r_1 \ \not\equiv \ r_2 \pmod{2}$}, \\[6pt]
  O\!\left(N^{-\lfloor \min(r_1,r_2)/2 \rfloor}\right) 
    & \text{if $p_1\ =\ q_2$ or $p_2\ =\ q_2$}.
\end{cases} \\[10pt]
B'_{r_1,r_2,\cF}(p_1,p_2)\ &= \
\begin{cases}
  \tfrac{C_{r_2/2}}{q_1^{r_1/2}} 
    + O_{k,q_1}\!\left(\tfrac{2^{r_2}p_2^{r_2/4}\log(p_2^{r_2}N)}{N}\right) 
    & \text{if $p_1\ =\ q_2$, $r_1$ even, $r_2$ even}, \\[6pt]
  O_{k,q_1}\!\left(\tfrac{2^{r_2}p_2^{r_2/4}\log(p_2^{r_2}N)}{N}\right) 
    & \text{if $p_1\ =\ q_2$, $r_1$ even, $r_2$ odd}, \\[6pt]
  O_{q_1,k}\!\left(\tfrac{2^{r_2}p_2^{r_2/4}\log(p_2^{r_2})}{N}\right) 
    & \text{if $p_1\ =\ q_1$ and $r_1$ odd}, \\[6pt]
  O\!\left(\tfrac{2^{r_2}p_2^{r_2/4}\log(p_2^{r_2}N)}{N^{\lfloor r_1/2\rfloor+1}}\right) 
    & \text{if $p_1\ =\ q_2$}.
\end{cases} \\[10pt]
B_{r_1,r_2}(p_1,p_2)\ &=\
\begin{cases}
  C_{(r_1+r_2)/2} 
    + O\!\left(2^{r_1+r_2}(p_1^{r_1}p_2^{r_2})^{1/4}
    \frac{\log(2p_1^{r_1}p_2^{r_2}N)\log(N)}{k^{5/6}N}\right) 
    & \text{if $p_1 = p_2$ and $(r_1-r_2 \bmod 2)=0$}, \\[6pt]
  C_{r_1/2}C_{r_2/2} 
    + O\!\left(2^{r_1+r_2}(p_1^{r_1}p_2^{r_2})^{1/4}
    \frac{\log(2p_1^{r_1}p_2^{r_2}N)\log(N)}{k^{5/6}N}\right) 
    & \text{if $p_1 \neq p_2$ and $r_1,r_2$ even}, \\[6pt]
  O\!\left(2^{r_1+r_2}(p_1^{r_1}p_2^{r_2})^{1/4}
    \frac{\log(2p_1^{r_1}p_2^{r_2}N)\log(N)}{k^{5/6}N}\right) 
    & \text{otherwise}.
\end{cases}
\end{align*}\normalsize
\end{lemma}
\begin{comment}
B_{r_1,r_2,\cF}(p_1,p_2)\ &= \
\begin{cases}
  C_{r_1/2}C_{r_2/2} 
    + O\!\left(\tfrac{2^{r_1+r_2}p_1^{r_1/4}p_2^{r_2/4}\log(p_1^{r_1}p_2^{r_2}N)}{N}\right) 
    & \text{if $r_1,r_2$ both even}, \\[6pt]
  O\!\left(\tfrac{2^{r_1+r_2}p_1^{r_1/4}p_2^{r_2/4}\log(p_1^{r_1}p_2^{r_2}N)}{N}\right) 
    & \text{if $r_1$ or $r_2$ even}.
\end{cases} Stella: it should be consistent with Lemma 6.4
\end{comment}

\begin{lemma}\label{thm:q1q2infmom}
Let $\cF=H^*_k(N)$ and $N=q_1q_2$ with $q_1,q_2$ prime and $q_1\sim N^{\delta},q_2\sim N^{1-\delta}$, where $\delta \in (0,1/2].$ Then,
\footnotesize
\begin{align*}
A'_r(p) \;=\;&\; O\!\left(\tfrac{1}{N^{r\delta/2}} \right) \\[8pt]
A_{r,\cF}(p) \;=\;& 
\begin{cases}
  C_{r/2} + O_k\!\left(\tfrac{2^r p^{r/4}\log(p^rN)}{N}\right) 
    & \text{if $r$ even}, \\[6pt]
  O_k\!\left(\tfrac{2^r p^{r/4}\log(p^rN)}{N}\right) 
    & \text{if $r$ odd}.
\end{cases} \\[10pt]
B''_{r_1,r_2}(p_1,p_2) \;=\;&\; O\!\left(\tfrac{1}{N^{(r_1+r_2)\delta/2}} \right) \\[8pt]
B'_{r_1,r_2}(p_1,p_2) \;=\;&\; O\!\left(\tfrac{2^{r_2}}{N^{r_1\delta/2}}\right) \\[10pt]
B_{r_1,r_2}(p_1,p_2) \;=\;& 
\begin{cases}
  C_{r_1/2}C_{r_2/2} 
    + O\!\left( 2^{r_1+r_2}p_1^{r_1/4}p_2^{r_2/4}N^{-1}
      \log\!\big(2p_1^{r_1/4}p_2^{r_2/4}N\big)\right) 
    & \text{if $r_1,r_2$ even and $p_1 \neq p_2$}, \\[6pt]
  C_{(r_1+r_2)/2} 
    + O\!\left( 2^{r_1+r_2}p_1^{r_1/4}p_2^{r_2/4}N^{-1}
      \log\!\big(2p_1^{r_1/4}p_2^{r_2/4}N\big)\right) 
    & \text{if $r_1-r_2 \equiv 0 \pmod{2}$ and $p_1 = p_2$}, \\[6pt]
  O\!\left( 2^{r_1+r_2}p_1^{r_1/4}p_2^{r_2/4}N^{-1}
    \log\!\big(2p_1^{r_1/4}p_2^{r_2/4}N\big)\right) 
    & \text{otherwise}.
\end{cases}
\end{align*}

\end{lemma}

\begin{lemma} \label{lem:psquareinfmom}
    Let $\cF=H^*_k(N)$ and $N=p^2$ with $p\to \infty.$ Then,
    \footnotesize
        \begin{align*}
        A'_{r,\cF}(q) \ &= \ 0\\
        A_{r,\cF}(q) \ &=  \begin{cases}
        C_{r/2} + O_k\left(2^rq^\frac{r}{4} \frac{\log(2q^{r})}{N}\right)& r  \ \text{even} \\
       O_k\left(2^rq^\frac{r}{4} \frac{\log(2q^{r})}{N}\right) & r  \ \text{odd}
    \end{cases}\\
        B''_{r_1,r_2,\cF}(p_1,p_2) \ &= \ 0 \\
        B'_{r_1,r_2,\cF}(p_1,p_2) \ &= \ 0 \\
        B_{r_1,r_2,\cF}(p_1,p_2)&= \begin{cases}
            C_{r_1/2}C_{r_2/2} \ + \ O_k\left(\frac{1}{N} 2^{r_1+r_2}{p_1^{r_1/4}p_2^{r_2/4}}\log\left(2p_1^{r_1}p_2^{r_2}\right)\right) & \text{if} \ p_1 \neq p_2 \ \text{and} \ r_1,~r_2  \ \text{even}\\
            O_k\left(\frac{1}{N} 2^{r_1+r_2}{p_1^{r_1/4}p_2^{r_2/4}}\log\left(2p_1^{r_1}p_2^{r_2}\right)\right) & \text{if} \ p_1 \neq p_2 \ \text{and} \ r_1,~r_2  \ \text{odd}\\
            C_{\frac{r_1+r_2}{2}} \ + \ O_k\left(\frac{1}{N} 2^{r_1+r_2}{p_1^{r_1/4}p_2^{r_2/4}}\log\left(2p_1^{r_1}p_2^{r_2}\right)\right) & \text{if} \ p_1 = p_2 \ \text{and} \ r_1-r_2\equiv0 \mod{2}\\
            O_k\left(\frac{1}{N} 2^{r_1+r_2}{p_1^{r_1/4}p_2^{r_2/4}}\log\left(2p_1^{r_1}p_2^{r_2}\right)\right) & \text{if} \ p_1 = p_2 \ \text{and} \ r_1-r_2\equiv1 \mod{2}.\\
        \end{cases}
    \end{align*}
\end{lemma}

\begin{remark}
Notice that for the three cases where $N$ prime, $N$ product of two primes with both factors going to infinity, and $N$ square, the terms $A', B'',$ and $B'$ do not admit a main term. On the other hand, for $N =q_1q_2$ where $q_1$ is fixed, $A', B'',$ and $B'$ have main terms. Moreover, in all four cases depending on $N$, even up to the error terms, the $A$ and $B$ terms are equal. This is an important remark to keep in mind because it will save us a lot of computations later.
\end{remark}

\vspace{1em}
\section{Lower Order Terms} \label{chp: lower order terms}
\begin{theorem}\label{thm:SA'}
Suppose the test function $\phi$ is an even Schwartz function with $\hat\phi$ supported in $[-\sigma,\sigma]$ for $\sigma<0.22$ \footnote{The restriction for the support arises in the proof.} If $N$ is prime or $N=q_1q_2$ where both $q_1 $ and $q_2$ goes to infinity, $S_{A'}(\cF)$ is negligible, i.e., $$S_{A'}(\cF)\ = \ O\left(\frac{1}{\log^4(R)}\right).$$
    However, when $N = q_1q_2$ for fixed $q_1,$ we have
    \begin{align}
        S_{A'}(\cF) \ =\ -2\frac{\log(q_1)}{\log(R)}\frac{\hat\phi(0)}{q_1^2-1} - \frac{\log(q_1)}{\log(R)}\frac{\hat\phi''(0)}{q_1^2-1} + O\left(\frac{1}{\log^5(R)}\right).
    \end{align}
\end{theorem}
\noindent \textit{Proof.} We compute the term 
\[
-2\sum_{p: p|N}\sum_{m=1}^{\infty}\frac{A'_{m, \mathcal{F}}(p)}{p^{m/2}}\frac{\log(p)}{\log(R)}\hat{\phi}\left(m \frac{\log(p)}{\log(R)} \right)\notag \\
\]
from $S_1(\mathcal{F}, \phi)$ up to the desired error.

\vspace{1em}

\noindent \textbf{Case 1. } When $N$ is prime, by Lemma \ref{thm: N prime, weight term values}
\[
A'_{r,\cF}(p)= O(N^{-r/2}).
\]
Note that in this case $p = N$ is the only prime factor of $N$ and hence the sum becomes
\begin{equation}
-2\sum_{m = 1}^\infty \frac{1}{N^m}\frac{\log(N)}{\log(R)}\hat\phi\left(m \frac{\log(N)}{\log(R)}\right)
\ \lesssim\ \frac{1}{N}\left(\frac{\log(N)}{\log(R)}+\frac{(\log(N))^3}{(\log(R))^3}+\cdots\right)
\ \lesssim\ \frac{\log(N)}{N\log(R)}\ \lesssim\ \frac{1}{N}.
\end{equation}
\\
\noindent \textbf{Case 2. }When $N = q_1 \cdot q_2$ with $q_1 < q_2$ and $q_1$ is fixed, by Lemma \ref{lem: 2 primes 1 fixed}
\[
A'_{r,\cF}(p)= \begin{cases}
    q_1^{-r/2} &\text{ if $p=q_1$ and $r$ even,} \\
    O\left(\frac{\log(N)}{N}\right) &\text{ if $p=q_1 $ and $r$ odd}\\
    O\left(N^{-r/2}\right) &\text{ if $p = q_2$}.
\end{cases}\\
\]
Notice that in this case, the main term only comes from when $p=q_1$ and $r$  is even. Therefore, the main term is given by
\small\begin{align}
    -2\sum_{r=1}^\infty \frac{1}{q_1^{2r}}\frac{\log(q_1)}{\log(R)}\hat\phi\left(\frac{2r\log(q_1)}{\log(R)}\right)\ =\ -2\frac{\log(q_1)}{\log(R)}\frac{\hat\phi(0)}{q_1^2(1-q_1^{-2})}- \frac{\log(q_1)}{\log(R)}\frac{\hat\phi''(0)}{q_1^2(1-q_1^{-2})} + O\left(\frac{1}{\log^5(R)}\right).
\end{align}\normalsize
The equality follows from the power series expansion of the even function $\hat\phi.$ Moreover, the error term is bounded by 
\begin{align}
     &-2\sum_{r=1}^\infty \frac{\log(N)}{N q_1^{r-\frac{1}{2}}}\frac{\log(q_1)}{\log(R)}\hat\phi\left(\frac{(2r-1)\log(q_1)}{\log(R)}\right) -2\sum_{m=1}^\infty \frac{1}{N^{m/2} q_2^{m/2}}\frac{\log(q_2)}{\log(R)}\hat\phi\left(\frac{(2r-1)\log(q_2)}{\log(R)}\right)\nonumber \\
    &\lesssim \frac{1}{N}\left(\sum_{r=1}^\infty \frac{1}{q_1^r}\right) + \sum_{m=1}^\infty \frac{1}{N^{m}}\nonumber \\ &\lesssim \frac{1}{N}.
\end{align}

Thus, we get that in this case, 
\begin{align}
    -2\sum_{p: p|N}\sum_{m=1}^{\infty}\frac{A'_{m, \mathcal{F}}(p)}{p^{m/2}}&\frac{\log(p)}{\log(R)}\hat{\phi}\left(m \frac{\log(p)}{\log(R)} \right)\nonumber\\ &=\ -2\frac{\log(q_1)}{\log(R)}\frac{\hat\phi(0)}{q_1^2(1-q_1^{-2})}- \frac{\log(q_1)}{\log(R)}\frac{\hat\phi''(0)}{q_1^2(1-q_1^{-2})} + O\left(\frac{1}{\log^5(R)}\right).
\end{align} \\ \ \\
\textbf{Case 3.} When $N = q_1\cdot q_2$ and $q_1,q_2 \to \infty$, by Lemma \ref{thm:q1q2infmom} we have
\[
A'_{r,\cF}(p) \ =\ O\left(N^{-r\delta/2} \right), \quad \text{where } 0 \leq \delta < \frac{1}{2}.
\] Then, the term becomes an error term because 
\begin{align*}
    -2\sum_{p: p|N}\sum_{m=1}^{\infty}\frac{A'_{m, \mathcal{F}}(p)}{p^{m/2}}\frac{\log(p)}{\log(R)}\hat{\phi}\left(m \frac{\log(p)}{\log(R)} \right) &\ \lesssim\ \frac{1}{\log(R)}\left(\sum_{m=1}^{\infty}\frac{\log(q_1)}{N^{m\delta/2}q_1^{1+m/2}}\hat{\phi}\left(m \frac{\log(q_1)}{\log(R)} \right) \right. \notag \\ & \qquad\qquad\qquad \left. +\sum_{m=1}^{\infty}\frac{\log(q_2)}{N^{m\delta/2}q_2^{1+m/2}}\hat{\phi}\left(m \frac{\log(q_2)}{\log(R)} \right) \right)\\
    &\ \lesssim \frac{1}{\log(R)N^{\delta}} .
\end{align*}
\qed
\vspace{1em}
\begin{theorem}\label{thm:SA} Suppose the test function $\phi$ is an even Schwartz function with $\hat\phi$ supported in $[-\sigma,\sigma]$ for $\sigma<0.22$. For all our cases, we have
       \begin{align}
         S_A(\cF) 
        &\ =\ \frac{\phi(0)}{2} + \frac{\hat\phi(0)}{\log(R)}\left(2\gamma_{PNT3} - \gamma_{A,2} + \gamma_{A,3} + \gamma_{A,6}\right) + \frac{\hat\phi''(0)}{\log^3(R)} \left(4\gamma_{A,1} + \gamma_{A,4} + \gamma_{A,5} + \gamma_{A,7}\right),
    \end{align}
    where 
    \footnotesize
    \begin{align}
        \gamma_{PNT3 } &\coloneqq 1+ \int_1^\infty \frac{E(t)}{t^2}dt \approx -1.33258\nonumber
       \\ \nonumber
       \gamma_{A,1} &\coloneqq \int_{1}^{\infty}\frac{E(t)}{t^{2}}\Bigl((\log(t))^{2}-2\log(t)\Bigr)dt \approx -10.0881\\\nonumber
        \gamma_{A,2} &\coloneqq \sum_p\frac{4\log(p)}{p(p+1)} \approx 1.5382\\\nonumber
        \gamma_{A,3} &\coloneqq \sum_p\frac{2(p^2+3p+1)\log(p)}{p(p+1)^3 }\approx 0.8852\\ \nonumber
        \gamma_{A,4} &\coloneqq \sum_p\frac{(32p^2+24p+8)\log^3(p)}{p(p+1)^3} \approx 43.6045 \\\nonumber
        \gamma_{A,5} &\coloneqq\sum_p\frac{(-64p^4+4p^3-44p^2-20p-4)\log^3(p)}{p(p+1)^5}\&\approx -72.6540\\ \nonumber
        \gamma_{A,6} &\coloneqq \sum_p\frac{2(p-1)\log(p)}{(p+1)}\sum_{r=2}^{\infty}\frac{C_rp^{r}}{(p+1)^{2r}}  \approx 0.8321\\ 
        \gamma_{A,7} &\coloneqq \sum_p\frac{(p-1)\log^3(p)}{(p+1)^3}\sum_{r=3}^{\infty}\frac{C_r p^{r}(4r^2(p-1)^2-24rp-8p)}{(p+1)^{2r}} \approx 5.8746529
        \end{align}\normalsize
    and where we let $\theta(t)=\sum_{p\leq t}\log(p)$ and define $E(t):=\theta(t)-t$ be the error.
\end{theorem}

\noindent \textit{Proof of Theorem \ref{thm:SA},}
 We note that, regardless of factorization, we have
\begin{align}
    A_{r,\cF}(p)=& 
\begin{cases}
        C_{r/2} + O_k\left(\frac{r2^rp^{r/4}\log(p)\log^2(N)}{N}\right) &\text{ if $r$ even}\\
        O_k\left(\frac{r2^rp^{r/4}\log(p)\log^2(N)}{N}\right)&\text{ if $r$ odd}.
        \end{cases}
\end{align}
Thus we have 
\begin{align}
    A_{0,\cF}(p)\ =& \ 1+O_k\left(\frac{\log(p)\log^2(N)}{N}\right)\nonumber\\
    A_{1,\cF}(p)\ =& \ O_k\left(\frac{p^{1/4}\log(p)\log^2(N)}{N}\right)\nonumber\\
    A_{2,\cF}(p)\ =& \ 1 +O_k\left(\frac{p^{1/2}\log(p)\log^2(N)}{N}\right).
\end{align}
Since 
\begin{align*}
    \frac{\log^2(N)}{N \log(R)}\sum_{p\leq R^{\sigma}} \frac{\log^2(p)}{p^2 } &\ \lesssim\ \frac{\log^2(N)}{R^{\sigma}N}\\
     \frac{\log^2(N)}{N \log(R)}\sum_{p\leq R^{\sigma}} \frac{\log^2(p)}{p^{5/4} }&\ \lesssim\ \frac{\log^2(N)}{R^{\sigma/4}N} \\
     \frac{\log^2(N)}{N \log(R)}\sum_{p\leq R^{\sigma}} \frac{\log^2(p)}{p^{3/2} } &\ \lesssim\ \frac{\log^2(N)}{R^{\sigma/2}N},
\end{align*}
the error terms for all the finite sums in Theorem \eqref{theorem: S_1} that are multiplied by a Schwartz function are sufficiently small and we have that the sums equal
\begin{align*}
&-2\hat{\phi}\left(0 \right)\sum_p\frac{2\log(p)}{p(p+1)\log(R)}+2\hat{\phi}\left(0 \right) \sum_p\frac{(p^2+3p+1)\log(p)}{p(p+1)^3 \log(R)} \\ 
&+\hat{\phi}''(0)\sum_p\frac{(32p^2+24p+8)\log^3(p)}{p(p+1)^3\log^3(R)} -\hat{\phi}''(0)\sum_p\frac{(64p^4-4p^3+44p^2+20p+4)\log^3(p)}{p(p+1)^5\log^3(R)}.
\end{align*}
We now deal with the terms
\begin{align*}
&2\sum_p
\frac{2A_{0, \mathcal{F}}(p)\log(p)}{p\log(R)}\hat{\phi}\left(2 \frac{\log(p)}{\log(R)} \right) \\
-&2\sum_p\frac{A_{1, \mathcal{F}}(p)\log(p)}{p^{1/2}\log(R)}\hat{\phi}\left(\frac{\log(p)}{\log(R)} \right)\\
-&2\sum_p\frac{A_{2,\mathcal{F}}(p)\log(p)}{p\log(R)}\hat{\phi}\left(2 \frac{\log(p)}{\log(R)} \right).
\end{align*}
In the appendix \ref{appendix Lemmas for main term}, we estimate many sums up to our required error, which we now employ to compute these expressions. In many cases, the error terms are left with convergent but non-elementary integrals that must be estimated numerically. Since the main terms for $A_2(\cF)(p)$ and $A_0(\cF)(p)$ agree and $A_1(\cF)(p)$ is an error, using Lemma \ref{lem:MSwork}, we have 
\footnotesize
\begin{align*}
    2\sum_p\frac{2A_{0,\mathcal{F}}(p)\log(p)}{p\log(R)}&\hat{\phi}\left(2 \frac{\log(p)}{\log(R)} \right)-2\sum_p\frac{A_{2,\mathcal{F}}(p)\log(p)}{p\log(R)}\hat{\phi}\left(2 \frac{\log(p)}{\log(R)} \right)\\
    &= \ \frac{\phi(0)}{2}
+\frac{2\hat\phi(0)}{\log(R)}\Bigl(1+\int_{1}^{\infty}\frac{E(t)}{t^{2}}dt\Bigr)
+\frac{4\hat\phi''(0)}{(\log(R))^{3}}
\int_{1}^{\infty}\frac{E(t)}{t^{2}}\Bigl((\log(t))^{2}-2\log(t)\Bigr)dt
+O\Bigl(\frac{1}{\log^4(R)}\Bigr)
\end{align*}
\normalsize
as the main term for these two sums. We now show the error terms in these cases are sufficiently small. Using compact support, we have our error terms as
\begin{align}
&\frac{\log(R)}{N}\sum_{p \leq R^{\sigma}}
\frac{\log^2(p)}{p}\ \lesssim\ \frac{\log^3(R)}{N}\nonumber\\
&\frac{\log(R)}{N}\sum_{p \leq R^{\sigma}}\frac{\log^2(p)}{p^{1/4}}\ \lesssim\ \frac{R^{3\sigma/4}\log(N)}{N}\nonumber\\
&\frac{\log(R)}{N}\sum_{p \leq R^{\sigma}}\frac{\log^2(p)}{p^{1/2}}\ \lesssim\ \frac{R^{\sigma/2}\log(N)}{N},\
\end{align}
which for restricted support are all sufficient. We now evaluate the last two sums. Using \eqref{lem: Computation of small m} and \eqref{lem: Computation of tail m}, we have 
\begin{equation}
    \sum_{p}\sum_{r=3}^{\infty}\frac{A_{r, \mathcal{F}}(p)p^{r/2}(p-1)\log(p)}{(p+1)^{r+1}\log(R)}\ =\ \sum_p\frac{(p-1)\log(p)}{(p+1)\log(R)}\sum_{r=2}^{\infty}\frac{C_rp^{r}}{(p+1)^{2r}}+O\left(\frac{1}{\log^4(R)} \right).
\end{equation}
We now aim to evaluate 
$$\sum_p\sum_{r=3}^{\infty}\frac{A_{r, \mathcal{F}}(p)(p-1)(r^2(p-1)^2-12rp-8p)p^{r/2}\log^3(p)}{(p+1)^{r+3}\log^3(R)}.$$
We first show the tail is negligible. We see 
\footnotesize
\begin{align}
    \sum_p\sum_{r=1+2\log(R)}^{\infty}\frac{A_{r, \mathcal{F}}(p)((p-1)(r^2(p-1)^2-12rp-8p)p^{r/2}\log^3(p)}{(p+1)^{r+3}\log^3(R)} \ \lesssim \ \frac{1}{\log^3}\sum_p\log^3(p)\sum_{r=1+2\log(R)}^{\infty} r^2\left(\frac{2p^{1/2}}{p+1} \right)^r.
\end{align}
\normalsize
Further,
\begin{align}
     \frac{1}{\log^3(R)}\sum_p\log^3(p)\sum_{r=1+2\log(R)}^{\infty} r^2\left(\frac{2p^{1/2}}{p+1} \right)^r &\lesssim\ \frac{1}{\log^3(R)}\sum_p\log^3(p)\left(\frac{p^{1/2}}{p+1} \right)^{2\log(R)+1}\sum_{r=1}^{\infty} r^2\left(\frac{2p^{1/2}}{p+1} \right)^r\nonumber\\ 
     &\lesssim\ \frac{1}{\log^3(R)}\sum_p\log^3(p)\left(\frac{p^{1/2}}{p+1} \right)^{2\log(R)+1}\sum_{r=0}^{\infty} r^2\left(\frac{2p^{1/2}}{p+1} \right)^r\nonumber\\
      &\lesssim\ \frac{1}{\log^3(R)}\sum_p\log^3(p)\left(\frac{2p^{1/2}}{p+1} \right)^{2\log(R)+1}\nonumber\\
      &\lesssim\ \frac{1}{R^{-.11}\log^3(R)}.
\end{align}
Thus the tail is negligible. We now show the error term that arises from $A_r(\cF)$ in the truncated sum is also negligible. Using the same asymptotic for the rational function in the sum, we have that the error is 
\begin{align*}
     &\frac{1}{N\log^3(R)}\sum_{p\leq R^{\sigma}}\log^4(p)\sum_{r=3}^{2\log(R)}r^3\left(\frac{2p^{3/4}}{(p+1)}\right)^r \\
     &\ \lesssim\ \frac{1}{N\log^3(R)}\left[\sum_{p< 2025}\log^4(p)\sum_{r=3}^{2\log(R)}r^3\left(\frac{2p^{3/4}}{(p+1)}\right)^r \sum_{2025<p< R^{\sigma}}\log^4(p)\sum_{r=3}^{2\log(R)}r^3\left(\frac{2p^{3/4}}{(p+1)}\right)^r \right]\\
     &\ \lesssim\ \frac{1}{N}\left[2025\cdot \left(\frac{2\cdot3^{3/4}}{4}\right)^{2\log(R)} + \sum_{p\in[2027,R^{\sigma}]}\left(\frac{2p^{3/4}}{p+1} \right) \right]  \\
     &\ \lesssim\ \frac{R^{\max\{.11,3\sigma/4\}}}{N},
\end{align*}\normalsize
which is negligible for our support. Thus, we have 
    \begin{equation}
    \sum_p\frac{(p-1)\log^3(p)}{p+1}\sum_{r=2}^{\log(R)}\frac{C_r p^{r}(1 + 2r)\bigl(p^{2}(2r - 1) - p (10 + 4 r) + (2r - 1)\bigr)}{(p+1)^{2r}\log^3(R)}
\end{equation}
as our main term. This sum can be extended to infinity at the cost of $R^{-0.11}$, which is negligible. Thus, we have shown
\small\begin{align}
    S_A(\cF)\ =\ &\frac{\phi(0)}{2}
    +\frac{2\hat\phi(0)}{\log(R)}\Bigl(1+\int_{1}^{\infty}\frac{E(t)}{t^{2}}dt\Bigr)+\frac{4\hat\phi''(0)}{(\log(R))^{3}}
\int_{1}^{\infty}\frac{E(t)}{t^{2}}\Bigl((\log(t))^{2}-2\log(t)\Bigr)dt\nonumber\\
        &-2\hat{\phi}\left(0 \right)\sum_p\frac{2\log(p)}{p(p+1)\log(R)}+2\hat{\phi}\left(0 \right) \sum_p\frac{(p^2+3p+1)\log(p)}{p(p+1)^3 \log(R)}+\hat{\phi}''(0)\sum_p\frac{(32p^2+24p+8)\log^3(p)}{p(p+1)^3\log^3(R)}\nonumber\\
        &-\hat{\phi}''(0)\sum_p\frac{(64p^4-4p^3+44p^2+20p+4)\log^3(p)}{p(p+1)^5\log^3(R)}+2\hat\phi(0)\sum_p\frac{(p-1)\log(p)}{(p+1)\log(R)}\sum_{r=2}^{\infty}\frac{C_rp^{r}}{(p+1)^{2r}}\nonumber\\
        &+\hat\phi''(0)\sum_p\frac{(p-1)\log^3(p)}{(p+1)\log^3(R)}\sum_{r=3}^{\infty}\frac{C_r (4r^2(p-1)^2-24rp-8p)p^{r}}{(p+1)^{2r}}+ O\left(\frac{1}{\log^4(R)} \right).
\end{align}\normalsize Substituting our constants in yields the theorem. \qed

\begin{theorem}\label{thm:S_B''} Suppose the test functions $\phi_1,\phi_2$ are even Schwartz functions with $\hat\phi_1, \hat\phi_2$ supported in $[-\sigma,\sigma]$ for $\sigma<0.22$.
    If $N$ is prime, $N=q_1q_2$ where both $q_1 $ and $q_2$ goes to infinity, or $N=p^2,$ then $S_{B''}(\cF)$ is negligible, i.e., $$S_{B''}(\cF)\ =\ O\left(\frac{1}{\log^4(R)}\right).$$
    However, when $N = q_1q_2$ for fixed $q_1,$ we find that
    \begin{align}
S_{B''}(\cF)\ =\ \frac{\hat\phi_1(0)\hat\phi_2(0)\log^2(q_1)}{\log^2(R)q_1^4(1-q_1^{-2})^2}  + O\left(\frac{1}{\log^4(R)}\right).
    \end{align}
\end{theorem}

\begin{theorem}\label{thm:S_B'}Suppose the test functions $\phi_1,\phi_2$ are even Schwartz function with $\hat\phi_1, \hat\phi_2$ supported in $[-\sigma,\sigma]$ for $\sigma<0.22$.
  If $N$ is prime, $N=q_1q_2$ where both $q_1 $ and $q_2$ goes to infinity, or $N=p^2,$ then $S_{B''}(\cF)$ is negligible, i.e. $$S_{B'}(\cF) = O\left(\frac{1}{\log^4(R)}\right).$$
    However, when $N = q_1q_2$ for fixed $q_1,$ we have
\begin{align}
S_{B'}(\cF)
&= \frac{2\log(q_1)\hat\phi_1(0)\hat\phi_2(0)}{\log^2(R)(q_1^2-1)}\left(\frac{\log(q_1)}{q_1}-\frac{3\log(q_1)q_1}{(q_1+1)^3} + \gamma_{B',1}\right)\notag\\&- 
\frac{\log(q_1)\left(\hat\phi_1(0)\phi_2(0) + \phi_1(0)\hat\phi_2(0)\right)}{4\log^2(R)(q_1^2-1)} -\frac{2\log(q_1)\hat\phi_1(0)\hat\phi_2(0)}{\log^3(R)(q_1^2-1)}\gamma_{PNT3} + O\left(\frac{1}{\log^4(R)}\right),
\end{align}
where \begin{align}
   & \gamma_{B',1} := \int_0^\infty (t+E(t))\frac{6t-3}{(t+1)^4}dt\approx 0.7425,
    &&\gamma_{PNT3} = 1+\int_1^\infty \frac{E(t)}{t^2} \approx -1.332.
\end{align}\end{theorem}

For $S_{B_\infty }(\cF)$,
we start by noticing that the expressions for $A_r(p)$ and $B_{r_1,r_2}(p_1,p_2)$ are the same up to error terms for the families with prime level or families with level being two distinct primes. The error term only differs by an exponent on $1/N$, making the computations more or less the same. 

 First, we compute the main terms; the computation is exactly the same for all three families. 

 \begin{theorem} \label{thm:S_Bf} Suppose the test functions $\phi_1,\phi_2$ are even Schwartz function with $\hat\phi_1, \hat\phi_2$ supported in $[-\sigma,\sigma]$ for $\sigma<0.22$. Define $\phi(x) := (\phi_1 * \phi_2)(x)$, with $\hat\phi''$ the second derivative of $\hat\phi$. Then for all of our cases
     \begin{align}
        S_{B_f}(\cF) 
    &\ =\ \int_0^\infty u\hat\phi(u)du + \frac{5\phi_1(0)\phi_2(0)}{16\log^2(R)}+\frac{\hat\phi(0)}{\log^2(R)}(1-\gamma_{PNT1}-\gamma_{PNT2})-\frac{\hat\phi''(0)}{\log^2(R)}\gamma_{PNT2}\notag\\
    & +\frac{5(\hat\phi_1(0)\phi_2(0) + \hat\phi_2(0)\phi_1(0))}{4\log^3(R)}\gamma_{PNT3} + O\left(\frac{1}{\log^4(R)}\right)\end{align}
where
\begin{align*}
    &\gamma_{PNT1}:= \int_1^{\infty}\frac{E(t)}{t^2}(1-\log(t))dt \approx 2.546,
   && \gamma_{PNT2}:= \int_1^\infty \frac{E(t)(1-2\log(t))}{t^3}dt \approx 1.633\\
   & \gamma_{PNT3}:= 1+\int_{1}^{\infty}\frac{E(t)}{t^{2}}dt \approx -1.332
\end{align*}
    \end{theorem}

    
\begin{theorem} \label{thm:S_infty}
   Suppose the test functions $\phi_1,\phi_2$ are even Schwartz functions with $\hat\phi_1, \hat\phi_2$ supported in $[-\sigma,\sigma]$ for $\sigma<0.22$, and with $\hat\phi_i''$ the second derivative of $\hat\phi_i$. Then for all of our cases,
\small\begin{align*}
    S_{B_\infty}(\cF) &= \frac{\hat\phi_1(0)\hat\phi_2(0)}{\log(R)}
\left(\frac{3}{4}-\log(2)-\gamma_3+\tfrac{\gamma_4-\gamma_5}{2}
+3\gamma_2-4\gamma_2\log 2+2\gamma_2(\gamma_4-\gamma_5)-4\gamma_2\gamma_3\right) \\ &+\frac{\hat\phi_1(0)\hat\phi_2(0)}{\log^2(R)}
\left({2\gamma_1-\frac{\gamma_6}{2}+2\gamma_7-2\gamma_8-\frac{7}{18}+\gamma_9+\gamma_{10}+\gamma_{11}}\right)+\frac{\hat\phi_1(0)\hat\phi_2''(0)+\hat\phi_1''(0)\hat\phi_2(0) }{\log^2(R)}\gamma_1
\end{align*} \normalsize
where 
\small\begin{align*}
    \gamma_{1}&:= \int_{1}^\infty \frac{E(t)(2-5\log(t))}{t^{7/2}}dt &&\approx 0.2953\\
    \gamma_2 &:= 
     \int_1^\infty \frac{E(t)(2t+1)}{(t^2+t)^2}dt &&\approx 0.1914\\
      \gamma_3 &:= \int_1^\infty \frac{E(t)(2t+1)}{(t^2+t)^2}dt &&\approx 0.1914\\
      \gamma_4&:=\int_1^\infty E(t)\frac{2t^3+8t^2+4t+1}{t^2(t+1)^4}dt &&\approx -0.000501\\
      \gamma_5 &:=\int_1^\infty E(t)\frac{4x^{3} - 5x^{2} + 4x + 1}{x^{2} \left(x + 1\right)^{4}}dt &&\approx -0.4854\\
      \gamma_6 &:= \int_0^\infty t\frac{e^t+3+e^{-t}}{(e^t+1)^3}dt&&\approx 0.210279 \\
      \gamma_7 &:= \int_1^\infty \frac{E(t)((-3x^3-12x^2-8x-2)\log(x)+x^3+4x^2+4x+1)}{t^3(t+1)^4} &&\approx 0.07780\\
      \gamma_8 &:= \int_1^\infty \frac{E(t)((-9t^4+10t^2+10t+3)\log(t)+3t^4+3t^3-2t^2-3t-1))}{t^4(t+1)^4}dt &&\approx 0.06586\\
      \gamma_9 &:=\sum_{p_1, p_2} \log \left(p_1\right) \log(p_2)\left[P_0\left(p_1\right) P_0\left(p_2\right)+\sum_{\ell=1}^{\infty} C_\ell\left(P_0\left(p_1\right) P_{2 \ell}\left(p_2\right)+P_{2 \ell}\left(p_1\right) P_0\left(p_2\right)\right)\right]&&\approx 0.5135 \\
      \gamma_{10} &:= \sum_{p_1 p_2} {\log (p_1) \log(p_2)} \sum_{\ell_1, \ell_2=1}^{\infty} C_{\ell_1} C_{\ell_2} P_{2\ell_1}(p_1) P_{2\ell_2}(p_2) &&\approx 0.4014\\
      \gamma_{11} &:=\sum_p \log^2(p)\sum_{\ell_1, \ell_2=1}^{\infty}\left[\left(C_{\ell_1+\ell_2}-C_{\ell_1} C_{\ell_2}\right) P_{2\ell_1}\left(p\right) P_{2\ell_2}(p)+C_{\ell_1+\ell_2-1} P_{2\ell_1-1}(p) P_{2 \ell_2-1}(p)\right] &&\approx 1.9648.
\end{align*}\normalsize
\end{theorem}

\textit{Proof Sketch of Theorems \ref{thm:S_B''}, \ref{thm:S_B'},\ref{thm:S_Bf}, and \ref{thm:S_infty}.}

The above theorems follow from substituting the terms $A', A, B'', B,$ and $B$ obtained from Lemmas \ref{thm: N prime, weight term values}, \ref{lem: 2 primes 1 fixed}, \ref{thm:q1q2infmom}, and \ref{lem:psquareinfmom} for each case into the formulas for $S_{B''}(\cF), S_{B'}(\cF), S_{B_f}(\cF),$ and $S_{B_\infty}$ given in Theorem \ref{thm: S_2.Formula}.

There are two parts to the proof of the theorem. The first is evaluating the main terms from merely plugging in what we've already proved. For example, one of the main terms that come up in $S_{B_\infty}$ is, after plugging in our results,
\begin{align}\label{eq:SumExample}
    \hat{\phi}_1(0)\sum_{p}\sum_{m_1=0}^{\infty}\hat{\phi}_2\left(\frac{\log(p)}{\log(R)} \right)\frac{\log^2 (p)}{\sqrt{p}\log^2(R)}\left(\frac{-(3p+1)}{p(p+1)^2)}+\sum_{\ell=2}^\infty C_\ell \frac{p^{\ell-1}(p-1)}{(p+1)^{2\ell}}\right).
\end{align}
Notice that the substitution still has a sum over $p$ of   $\hat\phi_2$ evaluated at $\log(p)/\log(R)$. The general procedure is as follows: to move the dependence on $\widehat\phi_i$ outside of the sum, we substitute its power series expansion and apply the Prime Number Theorem to compute the sum associated to the constant term. The other sums are generally error terms, since they involve larger powers of $\log(R)$ in the denominator. Often, we use techniques involving generating functions to simplify our expressions.

In the example \eqref{eq:SumExample}, we find that the expression is equal to, up to $O\left(1/\log^{4}(R)\right)$ error, $$\hat{\phi}_1(0)\left(-\frac{4\hat\phi_2(0)}{9\log^2(R)}+\frac{\hat \phi_2(0)+\hat\phi_2''(0)}{\log^2(R)}\int_1^\infty \frac{E(t)(2-5\log(t))}{t^{7/2}}dt \right).$$
As mentioned above, the remaining lemmas we use to get rid of the dependency on the test function as well as to simplify evaluating the main terms can be found in Appendix \ref{appendix Lemmas for main term}.

The second part of the proof is to argue that the error terms of terms $A', A, B'', B,$ and $B$ obtained from Lemmas \ref{thm: N prime, weight term values}, \ref{lem: 2 primes 1 fixed}, \ref{thm:q1q2infmom}, and \ref{lem:psquareinfmom} become error terms in computing $S_{B''}(\cF), S_{B'}(\cF), S_{B_f}(\cF),$ and $S_{B_\infty}$ as well. Below are the two main lemmas we use to bound the errors.

\begin{lemma} \label{lem: Computation of small m}Suppose $n$ is a fixed positive integer. Then,
   \begin{align}
       \sum_{p <R^{\sigma}}\log^n(p)\sum_{m=0}^{2\log(R)} \frac{2^m m p^{3m/4}(p-1)}{(p+1)^{m+1}}  \ \lesssim \ R^{\max(0.11,\,3\sigma/4)}.
   \end{align}
\end{lemma}
\textit{Proof of Lemma \ref{lem: Computation of small m}.} Note that \small
    \begin{align*}
         \sum_{p <R^{\sigma}}\log^n(p)\sum_{m=0}^{2\log(R)} \frac{2^m m p^{3m/4}(p-1)}{(p+1)^{m+1}}  & \ \leq\ \sum_{p <R^{\sigma}}\log^n(p)\sum_{m=0}^{2\log(R)} m\left(\frac{2p^{3/4}}{p+1}\right)^m \\
         &\ \lesssim\ \sum_{p<2025} \log^n(p)\sum_{m=0}^{2\log(R)} m\left(\frac{2p^{3/4}}{p+1}\right)^m+ \sum_{p\in[2027,R^{\sigma}]}\log^n(p)\sum_{m=0}^{2\log(R)} m\left(\frac{2p^{3/4}}{p+1}\right)^m \\
         &\ \lesssim\ 2025\cdot \left(\frac{2\cdot3^{3/4}}{4}\right)^{2\log(R)} + \sum_{p\in[2027,R^{\sigma}]}\left(\frac{2p^{3/4}}{p+1}\right)\\
         &\ \lesssim\ R^{0.12} +\sum_{p\in[2027,R^{\sigma}]} p^{-1/4} \\
         & \ \lesssim  \ R^{0.12} + R^{3\sigma/4}.
    \end{align*}
    \normalsize
     In the second to last line above, we can bound $2025\cdot \left(\frac{2\cdot3^{3/4}}{4}\right)^{2\log(R)} \lesssim R^{0.12}$ because \\$\log((2\cdot3^{3/4}/{4})^2) = \log(3\sqrt{3}/4) \approx 0.1136 < 0.12$. \qed

\begin{lemma} \label{lem: Computation of tail m}Suppose $n$ is any fixed positive integer.
\begin{align}
    \sum_{p} \log^n(p) \sum_{m=2\log(R)}^{\infty} \frac{2^{m}p^{m/2}(p-1)}{(p+1)^{m+1}} \ \lesssim \ \frac{1}{R^{0.11}}
\end{align}
\end{lemma}
\begin{proof} Observe that:
    \begin{align*}
         \sum_{p} \log^n(p) \sum_{m=2\log(R)}^{\infty} \frac{2^{m}p^{m/2}(p-1)}{(p+1)^{m+1}} &\ \leq\ \sum_{p} \log^n(p) \sum_{m=2\log(R)}\left(\frac{2\sqrt{p}}{p+1}\right)^m \ \lesssim\ \sum_{p}\log^n(p) \left(\frac{2\sqrt{p}}{p+1}\right)^{2\log(R)} \\
         &\ \lesssim\ \left(2025 \cdot \left(\frac{2\sqrt{2}}{3}\right)^{2\log(R)} +\sum_{p \geq 2027} \frac{\log^n(p)}{p^{2\log(R)/3}}\right) \ \lesssim\ {R^{-0.11}}.
    \end{align*}
   Again, it is crucial that $\frac{2\sqrt{2}}{3}<1$. The restriction for the support ($\sigma < 0.22$) comes from the fact that $\log(({2\sqrt{2}}/{3})^2) = \log(8/9) \approx -0.1178$ and that $-2\log(({2\sqrt{2}}/{3})^2) \geq0.22$.
\end{proof}
Every error term from those of $A', A, B'', B,$ and $B$ obtained from Lemmas \ref{thm: N prime, weight term values}, \ref{lem: 2 primes 1 fixed}, \ref{thm:q1q2infmom}, and \ref{lem:psquareinfmom} is in big O notation where the implied constant is absolute in $k.$ Therefore, we can bring in all the sum into the big O notation. Moreover, every error term has some infinite sum over $m_1$ or $m_2$ over $\N.$ Whenever we have this sum in $m_i$, we break the into two parts, into cases when $m_i  < 2\log(R)$ and when $m_i \geq 2\log(R).$ In the first case, when $m_i  < 2\log(R)$, we bound the error using the terms given from Lemmas \ref{thm: N prime, weight term values}, \ref{lem: 2 primes 1 fixed}, \ref{thm:q1q2infmom}, and \ref{lem:psquareinfmom} then apply Lemma \ref{lem: Computation of small m}. On the other hand, in the case where $m_i \geq 2\log(R),$ we use the bound $|B^*_{r_1,r_2}(p_1,p_2)|\leq 2^{r_1+r_2}$ for all $* = (''), ('), (),$ $r_1,r_2 \geq 1,$ and $p_1,p_2$ primes then apply Lemma \ref{lem: Computation of tail m}. 

For example, when we are dealing with the term \begin{equation}\hat{\phi}_1(0)\sum_{p_1,p_2<R^\sigma}\sum_{m_1=0}^{\infty}\hat{\phi}_2\left(\frac{\log(p_2)}{\log(R)} \right)\frac{\log(p_1) \log(p_2)}{\log^2(R)}B_{m_1,1}(p_1,p_2)\frac{P_{m_1}(p_1)}{\sqrt{p_2}}
\end{equation}
appearing in the computation of $S_{B_\infty}(\cF)$, the error term is
\small\begin{align}
&\hat{\phi}_1(0)\left(\sum_{p_1,p_2<R^\sigma}\sum_{m_1=0}^{2\log(R)}\hat{\phi}_2\left(\frac{\log(p_2)}{\log(R)} \right)\frac{\log(p_1) \log(p_2)}{\log^2(R)}\text{Error}(B_{m_1,1}(p_1,p_2))\frac{P_{m_1}(p_1)}{\sqrt{p_2}} + \sum_{p_1,p_2<R^\sigma}\sum_{m_1=2\log(R)}^{\infty}(\text{Same Terms})\right) .
\end{align}\normalsize
For the first sum, we use bounds from Lemmas \ref{thm: N prime, weight term values}, \ref{lem: 2 primes 1 fixed}, \ref{thm:q1q2infmom}, and \ref{lem:psquareinfmom} then apply Lemma \ref{lem: Computation of small m}:
\begin{align*}
&\hat{\phi}_1(0)\sum_{p_1,p_2}\sum_{m_1=0}^{2\log(R)}\hat{\phi}_2\left(\frac{\log(p_2)}{\log(R)} \right)\frac{\log(p_1) \log(p_2)}{\log^2(R)}\text{Error}(B_{m_1,1}(p_1,p_2))\frac{P_{m_1}(p_1)}{\sqrt{p_2}}\\
&\lesssim \frac{R^{\sigma/2}}{\log(R)}\left(\sum_{p_1}^{R^\sigma}\log(p_1)\sum_{m_1=0}^{2\log(R)} \frac{m_12^{m_1}p_1^{3m_1/4} \log(p_1N)}{N(p_1+1)^{m_1}}\right) \\
&\lesssim \frac{R^{\sigma/2}\log(R)}{N}\left(\sum_{p_1<2025}\log^2(p_1)\sum_{m_1=0}^{2\log(R)} m_1\left(\frac{2p_1^{3/4}}{p_1+1}\right)^{m_1} + \sum_{p_1=2027}^{R^{\sigma}}\log^2(p_1)\sum_{m_1=0}^{2\log(R)} m_1\left(\frac{2p_1^{3/4}}{p_1+1}\right)^{m_1}\right) \\
&\lesssim\frac{R^{\sigma/2}\log(R)}{N} \left(2025\left(\frac{2\cdot3^{3/4}}{4}\right)^{2\log(R)}\log^2(R) + \sum_{p=2027}^{R^\sigma} \frac{2p^{3/4}}{p_1+1}\right) \\&\lesssim \frac{R^{\sigma/2}\log^2(R)\cdot R^{0.11}}{N} + \frac{R^{\sigma/2}}{N}\log(R) \sum_{p=2027}^{R^\sigma} p_1^{-1/4} \lesssim \frac{R^{\sigma/2}\log^2(N)R^{0.11}}{N} + \frac{R^{\sigma/2}\log(R)}{N}N^{3\sigma/4}.
\end{align*} This error is negligable as long as $\sigma < 0.8,$ which we assumed.

For the second sum, we bound $|B_{m_1,1}(p_1,p_2)|\lesssim 2^{m_1}$ then apply Lemma \ref{lem: Computation of tail m}: \begin{align*}
&\hat{\phi}_1(0)\sum_{p_1,p_2}\sum_{m_1=0}^{2\log(R)}\hat{\phi}_2\left(\frac{\log(p_2)}{\log(R)} \right)\frac{\log(p_1) \log(p_2)}{\log^2(R)}\text{Error}(B_{m_1,1}(p_1,p_2))\frac{P_{m_1}(p_1)}{\sqrt{p_2}}  \\
&\lesssim\ \left(\frac{R^{\sigma/2}}{\log(R)}\right)\left(\sum_{p_1}\log(p)\sum_{m_1=2\log(R)}^\infty 2^{m_1}\frac{p_1^{m_1/2}}{(p_1+1)^{m_1}}\right) \\
&\lesssim\ \left(\frac{R^{\sigma/2}}{\log(R)}\right)\left(\sum_{p_1}\log(p)\sum_{m_1=2\log(R)}^\infty \left(\frac{2p_1^{1/2}}{p_1+1}\right)^{m_1}\right) \\
&\lesssim\ \left(\frac{R^{\sigma/2}}{\log(R)}\right)\sum_{p_1}\log(p)\left(\frac{2p_1^{1/2}}{p_1+1}\right)^{2\log(R)} \\
&\lesssim\ \left(\frac{R^{\sigma/2}}{\log(R)}\right)\left(2025\left(\frac{2p_1^{1/2}}{p_1+1}\right)^{2\log(R)} + \sum_{p_1\geq 2026} \frac{\log(p_1)}{p_1^{2\log(R)/3}}\right) \ \lesssim \ \frac{R^{\sigma/2}}{R^{0.11}\log^2(R)},
\end{align*}
which is negligible as long as $\sigma < 0.22,$ which we assume.
\appendix
\section{Proof of Lemma \ref{Lemma, S_1} } \label{appendix: Proof of Lemma M3,2}
\noindent \textbf{Lemma \ref{Lemma, S_1}.}
    \textit{For prime $p$, we have }
    \begin{align}
        M_{3,2}(p)\ =&\ \frac{32p^2+24p+8}{p(p+1)^3} - \frac{27p^3-17p^2+5p+1}{\sqrt{p}(p+1)^4}\lambda_f(p)-\frac{64p^4-4p^3+44p^2+20p+4}{p(p+1)^5} \lambda_f(p)^2
        \notag\\&+ \sum_{r=3}^{\infty}\frac{(p-1)(r^2(p-1)^2-12rp-8p)p^{r/2}\lambda_f(p)^r}{(p+1)^{r+3}}.
    \end{align} 

\noindent \textit{Proof of Lemma.} We recall $M_{3,2}(p)\ =\ \sum_{m=3}^{\infty}m^2\left(\frac{\alpha_f(p)}{p^{1/2}}\right)^m+\sum_{m=3}^{\infty}m^2\left(\frac{\beta_f(p)}{p^{1/2}}\right)^m$. We note that since \begin{align*}
    \frac{1}{1-x}\ &=\ \sum_{n=0}^{\infty}x^n, \\
   \frac{x}{(1-x)^2}\ &=\ \sum_{n=1}^{\infty}nx^n, \text{ and }\\ 2\frac{x^2}{(1-x)^3}\ &=\ \sum_{n=2}^{\infty }n(n-1)x^{n}\ =\ \sum_{n=2}^{\infty }n^2x^{n}-\sum_{n=2}^{\infty }nx^{n}.
\end{align*} We have
 \begin{align}
\sum_{n=2}^{\infty}n^2x^n\ &=\ 2\frac{x^2}{(1-x)^3}+\frac{x}{(1-x)^2}-x =\frac{x^2+x}{(1-x)^3}-x.
 \end{align}
Thus,
\begin{align}
 \sum_{m=3}^{\infty}m^2\left(\frac{\alpha_f(p)}{p^{1/2}}\right)^m+\sum_{m=3}^{\infty}m^2\left(\frac{\beta_f(p)}{p^{1/2}}\right)^m \ = \ &\frac{\left(\frac{\alpha_f(p)}{p^{1/2}}\right)^2+\frac{\alpha_f(p)}{p^{1/2}}}{\left(1- \frac{\alpha_f(p)}{p^{1/2}}\right)^3}+\frac{\left(\frac{\beta_f(p)}{p^{1/2}}\right)^2+\frac{\beta_f(p)}{p^{1/2}}}{\left(1- \frac{\beta_f(p)}{p^{1/2}}\right)^3} \notag \\
    &-4\left(\left(\frac{\alpha_f(p)}{p^{1/2}}\right)^2+\left(\frac{\beta_f(p)}{p^{1/2}}\right)^2\right)-\left(\frac{\alpha_f(p)}{p^{1/2}}+\frac{\beta_f(p)}{p^{1/2}}\right) .   
\end{align}
We see that by multiplying by $p^{3/2}$ we can write 
$$\frac{\left(\frac{\alpha_f(p)}{p^{1/2}}\right)^2+\frac{\alpha_f(p)}{p^{1/2}}}{\left(1- \frac{\alpha_f(p)}{p^{1/2}}\right)^3}\ =\ \frac{\alpha_f(p)^2p^{1/2}+\alpha_f(p)p}{(p^{1/2}-\alpha_f(p))^3}$$
and we obtain a similar result for the corresponding $\beta_f(p)$ term. We now aim to combine
\begin{equation}
    \frac{\alpha_f(p)^2p^{1/2}+\alpha_f(p)p}{(p^{1/2}-\beta_f(p))^3}+\frac{\beta_f(p)^2p^{1/2}+\beta_f(p)p}{(p^{1/2}-\beta_f(p))^3}.
\end{equation}
The denominator becomes 
\begin{align*}
    (p^{1/2}-\alpha_f(p))(p^{1/2}-\beta_f(p))\ &= \ p+\alpha_f(p)\beta_f(p)-(\alpha_f(p)+\beta_f(p))p^{1/2}\\
   \ &= \ p+1-\lambda_f(p)p^{1/2}.
\end{align*}
We note we are leaving the aforementioned expression inside the cube. Thus we have
\begin{align} \label{eq:S1 second}
  \frac{\lambda_f(p)^2(p^2-p)+\lambda_f(p)(p^2\sqrt{p}-\sqrt{p})-8(p^2-p)}{(p+1-\lambda_f(p)\sqrt{p})^3}-\frac{4(\lambda_f(p)^2-2)}{p}-\frac{\lambda_f(p)}{p^{1/2}}.
\end{align}
Looking at the first term, we express the denominator $(p+1-\lambda_f(p)\sqrt{p})^{3}$ as $(p+1)^{3}\left(1-\frac{\lambda_f(p)\sqrt{p}}{p+1} \right)^{3}$ and use the power series $\frac{1}{(1-x)^3}\ = \ \frac{1}{2}\sum_{n=2}^{\infty}n(n-1)x^{n-2}$ to obtain
$$\frac{1}{\left(1-\frac{\lambda_f(p)p^{1/2}}{p+1}\right)^3}\ =\ \frac{1}{2}\sum_{r=0}^{\infty}(r+2)(r+1)\left(\frac{\lambda_f(p)p^{1/2}}{p+1} \right)^{r}.$$
Thus we get $$\left(\lambda_f(p)^2(p^2-p)+\lambda_f(p)(p^2\sqrt{p}-\sqrt{p})-8(p^2-p)\right)\sum_{r=0}^{\infty}\frac{(r+2)(r+1)}{2}\frac{\lambda_f(p)^rp^{r/2}}{(p+1)^{r+3}},$$
which is equal to
\small
\begin{align*}
&-\frac{8(p^2-p)}{(p+1)^3}-\frac{24(p^2-p)p^{1/2}}{(p+1)^4}\lambda_f(p)-\frac{48(p^2-p)p}{(p+1)^5}\lambda_f(p)^2-\sum_{r=3}^{\infty}\frac{(r+2)(r+1)}{2}\frac{8(p^2-p)\lambda_f(p)^rp^{r/2}}{(p+1)^{r+3}}\\
&+\frac{p^{5/2}-p^{1/2}}{(p+1)^3}\lambda_f(p)+3\frac{(p^{5/2}-p^{1/2})p^{1/2}}{(p+1)^4}\lambda_f(p)^2+\sum_{r=2}^{\infty}\frac{(r+2)(r+1)}{2}\frac{(p^{5/2}-p^{1/2})\lambda_f(p)^{r+1}p^{r/2}}{(p+1)^{r+3}}\\
&+\frac{p^2-p}{(p+1)^3}\lambda_f(p)^2+\sum_{r=1}^{\infty}\frac{(r+2)(r+1)}{2}\frac{(p^2-p)\lambda_f(p)^{r+2}p^{r/2}}{(p+1)^{r+3}}.
\end{align*}\normalsize
Combining like terms yields
\begin{align*}
    &-\frac{8(p^2-p)}{(p+1)^3} + \frac{\sqrt{p}(p-1)(p^2-22p+1)}{(p+1)^4}\lambda_f(p)+\frac{4p(p-1)(p^2-10p+1)}{(p+1)^5}\lambda_f(p)^2 \\
    &+\sum_{r=3}^{\infty}\frac{(p-1)(r^2(p-1)^2-12rp-8p)p^{r/2}\lambda_f(p)^r}{(p+1)^{r+3}} .
\end{align*}
Plugging this back into equation \eqref{eq:S1 second} we get
\begin{align}\label{eq:S1 third}
\notag M_{3,2}(p)\ =&\ \frac{-8(p^2-p)}{(p+1)^3} + \frac{\sqrt{p}(p-1)(p^2-22p+1)\lambda_f(p)}{(p+1)^4}+\frac{4p(p-1)(p^2-10p+1)\lambda_f(p)^2}{(p+1)^5}\\&-\frac{4(\lambda_f(p)^2-2)}{p}-\frac{\lambda_f(p)}{\sqrt{p}}+ \sum_{r=3}^{\infty}\frac{(p-1)(r^2(p-1)^2-12rp-8p)p^{r/2}\lambda_f(p)^r}{(p+1)^{r+3}}   
\\\notag\\
=&\ \frac{32p^2+24p+8}{p(p+1)^3}-\frac{27p^3-17p^2+5p+1}{(p+1)^4}\lambda_f(p)
-\frac{64p^4-4p^3+44p^2+20p+4}{p(p+1)^5}\lambda_f(p)^2 \notag\\
&+\sum_{r=3}^{\infty}\frac{(p-1)(r^2(p-1)^2-12rp-8p)p^{r/2}\lambda_f(p)^r}{(p+1)^{r+3}}.
\end{align}\qed
\section{Proof of Lemma \ref{lemma1: both m_1,m_2 geq 3}} \label{appendix: Proof of Lemma both m1,m2 geq 3}
\noindent \textbf{Lemma \ref{lemma1: both m_1,m_2 geq 3}} \textit{Suppose $\phi_1$ and $\phi_2$ are even Schwartz function with $\hat\phi_1$ and $\hat\phi_2$ having support in $[-\sigma,\sigma].$
We then have the following estimate: }
\small{\begin{align}\label{lem:phi0.estimate.statement2}
    &\frac{1}{W_R(\cF)}\sum_{\substack{p_1,p_2}}\sum_{m_1,m_2\geq 3} C(m_1,m_2)\notag \\
    & \ = \ \frac{1}{W_R(\cF)}\sum_{p_1,p_2<R^\sigma}\sum_{m_1,m_2\geq 3}\sum_{\substack{f\in \cF\\p_1\nmid N_f\\p_2\nmid N_f }}w_R(f)\frac{\sum_{j=1}^2\left(\alpha_f(p_j)^{m_j}+\beta_f(p_j)^{m_j}\right)}{p_1^{m_1/2}p_2^{m_2/2}} \frac{\log(p_1)\log (p_2)}{\log^2(R)}\hat\phi_1(0) \hat \phi_2(0)  + O\left(\frac{1}{\log^4(R)}\right).
\end{align}}
\textit{Proof of Lemma.}
Let us write
\small
\begin{align*}
    &(A) \ :=\  \frac{1}{W_R(\cF)}\sum_{p_1,p_2}\sum_{m_1,m_2\geq 3} C(m_1,m_2),\\
   & (B)\ :=\ \frac{1}{W_R(\cF)}\sum_{p_1,p_2<R^\sigma}\sum_{m_1,m_2 \geq 3}\sum_{\substack{f\in \cF\\p_1\nmid N_f\\p_2\nmid N_f }}w_R(f)\frac{\sum_{j=1}^2\left(\alpha_f(p_j)^{m_j}+\beta_f(p_j)^{m_j}\right)}{p_1^{m_1/2}p_2^{m_2/2}} \frac{\log(p_1)\log (p_2)}{\log^2(R)}\hat\phi_1(0) \hat \phi_2(0), \\
   &(C)\  := \ \frac{1}{W_R(\cF)}\sum_{p_1,p_2<R^\sigma}\sum_{m_1,m_2 \geq 3 }\sum_{\substack{f\in \cF\\p_1\nmid N_f\\p_2\nmid N_f }}w_R(f)\frac{\sum_{j=1}^2\left(\alpha_f(p_j)^{m_j}+\beta_f(p_j)^{m_j}\right)}{p_1^{m_1/2}p_2^{m_2/2}}
       \frac{\log(p_1)\log (p_2)}{\log^2(R)} \hat\phi_1 \left(m_1 \frac{\log(p_1)}{\log(R)}\right) \hat\phi_2(0).
\end{align*}
\normalsize
Notice that the statement of the lemma is equivalent to saying that $(A)-(B) = O\left(1/\log^4(R)\right)$.
We first show that $(A)- (C) =O\left(1/\log^4(R)\right)$, then argue that $(C) - (B) = O\left(1/\log^4(R)\right).$
\footnotesize\begin{align*}\label{lemma.A-C}
&\frac{1}{W_R(\cF)}\sum_{p_1,p_2<R^\sigma}\sum_{m_1,m_2 \geq 3 }\sum_{\substack{f\in \cF\\p_1\nmid N_f\\p_2\nmid N_f }}w_R(f)\frac{\sum_{j=1}^2\left(\alpha_f(p_j)^{m_j}+\beta_f(p_j)^{m_j}\right)}{p_1^{m_1/2}p_2^{m_2/2}}
       \frac{\log(p_1)\log (p_2)\hat\phi_1 \left(m_1 \frac{\log(p_1)}{\log(R)}\right)}{\log^2(R)} \left[\hat\phi_2 \left(m_2 \frac{\log(p_2)}{\log(R)}\right) - \hat\phi_2(0)\right]
       \\
       &\ \lesssim  \  \frac{1}{\log^4(R)}\sum_{p_1<R^\sigma}\sum_{m_1\geq 3} \frac{\log(p_1)}{p_1^{m_1/2}}\hat\phi_1\left(m_1\frac{\log(p_1)} {\log(R)}\right)\\ &
       \lesssim \frac{1}{\log^4(R)}\sum_{p_1}\frac{\log(p_1)}{p(\sqrt{p}-1)}\\
       &\ \lesssim  \ \frac{1}{\log^4(R)}\sum_{p_1}\frac{\log(p_1)}{p^{3/2}}\\
       &\ \lesssim  \ \frac{1}{\log^4(R)}.
    \end{align*}\normalsize
The first approximation follows from using Taylor expansion, bounding $|\alpha_f(p)^m + \beta_f(p)^m|$ by 2, and finally observing that $\sum_{m_2\geq3 } {m_2^2}/{p_2^{m_2/2}} = O(p_2^{-3/2})$ and $\sum_{p_2}\log(p_2)^3/p_2^{3/2} < \infty.$ Therefore, we have that $(A)-(C) = O\left(\frac{1}{\log^4(R)}\right)$. 

Now we show that $(C)-(B)= O\left(\frac{1}{\log^4(R)}\right)$. By a similar approximation,   \small\begin{align*}
&\frac{1}{W_R(\cF)}\sum_{p_1,p_2<R^\sigma}\sum_{m_1,m_2 \geq 3 }\sum_{\substack{f\in \cF\\p_1\nmid N_f\\p_2\nmid N_f }}w_R(f)\frac{\sum_{j=1}^2\left(\alpha_f(p_j)^{m_j}+\beta_f(p_j)^{m_j}\right)}{p_1^{m_1/2}p_2^{m_2/2}}
       \frac{\log(p_1)\log (p_2)}{\log^2(R)}       \left(\hat\phi_1 \left(m_1 \frac{\log(p_1)}{\log(R)}\right)-\hat\phi_1(0)\right)\hat\phi_2(0) \\  \ &\lesssim \ \frac{1}{\log^4(R)} \left(\sum_{p_1<R^\sigma}\sum_{m_1\geq 3}\frac{m_1^2\log(p_1)^3}{p_1^{m_1/2}}\right)\left(\sum_{p_2<R^\sigma}\sum_{m_2\geq 3}\frac{\log(p_2)^3}{p_2^{m_2/2}}\right)\\ \ &\lesssim\  \frac{1}{\log^4(R)}\left(\sum_{p_1<R^\sigma}{p_1}^{-3/2}\right)\\ \ &\lesssim  \ 
       \frac{1}{\log^4(R)}.
    \end{align*}\normalsize
With this, the proof of Lemma \ref{lemma1: both m_1,m_2 geq 3} is complete.
\qed
 \section{Proof of Lemmas \ref{lem: 2 primes 1 fixed}, \ref{thm:q1q2infmom}, and \ref{lem:psquareinfmom}} \label{appendix: Proof of Lemmas for computing As and Bs}
In this section, we prove the Lemmas \ref{lem: 2 primes 1 fixed}, \ref{thm:q1q2infmom}, and \ref{lem:psquareinfmom}. 

\noindent\textit{Proof of Lemma \ref{lem: 2 primes 1 fixed}. }
The proofs for $A_r(p)$ and $B_{r_1,r_2}(p_1,p_2)$ are identical to the prime case so we omit it. First we look at $A'_{r,\cF}(p)$. 
There are two cases to consider: when $p \ = \ q_1$ and when $p\ =\ q_2.$ If $p\ =\ q_1$ and $r$ is even, then using the fact that $\lambda_f(q_1)^2 = \frac{1}{q_1}$, we have that $A'_{r}(q_1)\  =\ q_1^{-r/2}.$ If $r$ is odd, using Corollary 2.10 of \cite{ILS2} and substituting for $W_R(\cF)$ we found in lemma \eqref{lem: 2 primes 1 fixed}, we have: \begin{align*}
    A'_{r}(q_1) \ =\ \frac{q_1^{-\lfloor r/2\rfloor}}{W_{R}(\cF)} \sum_{f\in\cF} w_R(f) \lambda_f(q_1) \ =\ \frac{q_1^{-\lfloor r/2\rfloor}}{W_{R}(\cF)} \left(O(k^{1/6}q_1^{-1/4}\log(q_1N)\right) = O_{q_1,k}\left(\frac{\log(N)}{N}\right).
\end{align*}
Moreover, when $p\ =\ q_2,$ using triangle inequality and using the fact that $w_R(f) \geq 0$ for all $f \in \cF$, 
$$|A'_r(q_2)|\ \leq \ \frac{1}{q_2^{r/2}}\left(\frac{1}{W_R(\cF)}\sum_{f\in\cF}w_R(f)\right)\ \lesssim \ N^{-r/2}.$$
Next, we consider $B''_{r_1,r_2}(p_1,p_2)$. There are four pairs possible for $(p_1,p_2).$ They are: $(q_1,q_1), (q_1,q_2), (q_2,q_1), $ and $(q_2,q_2)$. 

In the first case, where $(p_1 , p_2) = (q_1,q_1),$ we have \begin{align*}
    B''_{r_1,r_2}(q_1,q_1) \ &= \ \frac{1}{W_R(\cF)} \sum_{f \in \cF} w_R(f) \lambda_f(q_1)^{r_1+r_2} \\&= \ \frac{1}{q_1^{\lfloor (r_1+r_2) / 2\rfloor}}\frac{1}{W_R(\cF)}\sum_{f \in \cF} w_f\lambda_f(q_1)^{((r_1+r_2)\modd 2)} \\
   &= \ \begin{cases}
    q_1^{(r_1+r_2) / 2} 
      & \text{if $r_1+r_2$ even}, \\[4pt]
    O_{k,q_1}\!\left(\frac{\log^2(N)}{N}\right) 
      & \text{if $r_1+r_2$ odd}.
\end{cases}
\end{align*}
In the last three cases, with at least one of $p_1$ or $p_2$ being $q_2,$ we have:
\begin{align*}
      B''_{r_1,r_2}(p_1,p_2)\ &=\ \frac{1}{W_R(\cF)} \sum_{f \in \cF} w_R(f) \lambda_f(p_1)^{r_1}\lambda_f(p_2)^{r_2} \\
     \ &=\ \frac{1}{p_1^{\lfloor r_1/2\rfloor}}\frac{1}{p_2^{\lfloor r_2/2\rfloor}}\left(\frac{1}{W_R(\cF)} \sum_{f \in \cF} w_R(f) \lambda_f(p_1)^{(r_1\modd 2)}\lambda_f(p_2)^{(r_1\modd 2)} \right)\\
      &\lesssim \frac{1}{q_2^{\lfloor \min(r_1,r_2)/2 \rfloor}} \left(\frac{\log(N)}{N^{1/4}}\right) \lesssim \ \frac{1}{N^{\lfloor \min(r_1,r_2)/2 \rfloor}}.
\end{align*}
Next, we look at $B'_{r_1,r_2}(p_1,p_2).$ By definition, $p_1 | N$ and $p_2 \nmid N.$ 
If $p_1 = q_1,$ we see
\begin{align*}
    B'_{r_1,r_2}(q_1,p_2)\ &= \ \frac{1}{W_R(\cF)} \sum_{f \in \cF} w_R(f) \lambda_f(q_1)^{r_1}\lambda_f(p_2)^{r_2} \\
   \ &=\ \frac{1}{q_1^{\lfloor r_1/2\rfloor}} \frac{1}{W_R(\cF)} \sum_{f \in \cF} w_R(f) \lambda_f(q_1)^{(r_1 \modd 2)}\lambda_f(p_2)^{r_2}.
\end{align*}
\noindent If $r_1 $ is even, it follows that $$B'_{r_1,r_2}(q_1,p_2) \ =\ 
        \frac{1}{q_1^{ r_1/2}} \left(A_{r_2}(p_2)\right)\ =\ \begin{cases}
         \frac{1}{q_1^{ r_1/2}} C_{r_2/2} + O_{k,q_1}\left(\frac{r_22^{r_2}p_2^{r_2/4}\log(p_2^{r_2}N)\log(N)}{N}\right)  &\text{ if $r_2$ even}   \\
         O_{k,q_1}\left(\frac{r_22^{r_2}p_2^{r_2/4}\log(p_2^{r_2}N)\log(N)}{N}\right)  &\text{ if $r_2$ odd}
        \end{cases}.$$

\noindent If $r_1$ is odd, Corollary 2.10 of \cite{ILS2} yields\begin{align*}
    B'_{r_1,r_2}(q_1,p_2)\ =\ O_{q_1,k}\left(\frac{2^{r_2}p_2^{r_2/4}\log(p_2)\log^2(N)}{N}\right).
\end{align*} 

\noindent Next, suppose we have that $p_1 = q_2.$ Then, trivially bounding by $|\lambda_f(p_2)| \leq 2$ gives:
\begin{align*}
     B'_{r_1,r_2}(q_2,p_2)\ &=\ \frac{1}{W_R(\cF)} \sum_{f \in \cF} w_R(f) \lambda_f(q_2)^{r_1}\lambda_f(p_2)^{r_2} \\
    &=\ \frac{1}{q_2^{\lfloor r_1/2\rfloor}} \frac{1}{W_R(\cF)} \sum_{f \in \cF} w_R(f) \lambda_f(q_2)^{(r_1 \modd 2)}\lambda_f(p_2)^{r_2} \\
    &=\  O\left(\frac{2^{r_2}p_2^{r_2/4}\log(p_2^{r_2}N)}{N^{\lfloor r_1/2\rfloor + 1/2}}\right),
\end{align*}
completing the proof. \qed \\

Next, we prove Lemma \ref{thm:q1q2infmom}.

\noindent\textit{Proof of Lemma \ref{thm:q1q2infmom}.}
The proofs of $A_r(p)$ and $B_{r_1,r_2}(p_1,p_2)$ is identical to the prime case so we omit it.
    We can see 
    \begin{align*}
        |A'_r(p)|&\leq\frac{1}{W_R(\cF)}\sum_{\substack{f \in \cF\\p|N}}w_{R}(f)|\lambda_{f}(p)|^r\leq \frac{1}{p^{r/2}}.
    \end{align*}
    Since $N^{-(1-\delta)/2}\lesssim N^{-\delta/2}$, this gives
    $$A'_r(p)\ =\ O\left(\frac{1}{N^{r\delta/2}} \right).$$
    We now look at $B''_{r_1,r_2}(p_1,p_2)$. We see we must have $(p_1,p_2)\in \{(q_1,q_2), (q_1,q_1),(q_1,q_2),(q_2,q_1)\}\ =\ A$, and $(q_2,q_2)$. Using the fact that $|\lambda_f(p)|\ =\ p^{-1/2}$ for $p|N$, we have 
    \begin{align}
       |B''_{r_1,r_2}(p_1,p_2)| \ &=\ \left|\frac{1}{W_R(\cF)}\sum_{\substack{f \in \cF }}w_R(f)\lambda_f(p_1)^{r_1}\lambda_f(p_2)^{r_2} \right|\notag \\
       &\leq \frac{1}{W_R(\cF) }\sum_{\substack{f \in \cF }}w_R(f)|\lambda_f(p_1)|^{r_1/2}|\lambda_f(p_2)|^{r_2/2} \notag \\
       &\leq \frac{1}{W_R(\cF)}\sum_{\substack{f \in \cF }}w_R(f) \frac{1}{p_1^{r_1/2}}\frac{1}{p_2^{r_2/2}} \notag \\
        &\lesssim\frac{1}{N^{(r_1+r_2)\delta/2}}.
    \end{align}
    For $B'_{r_1,r_2}(p_1,p_2)$, we have $p_1\ =\ q_1,q_2$. 
    We can write 
    $$B'_{r_1,r_2}(p_1,p_2)=\frac{1}{p_1^{r_1/2}W_R(\cF)}\sum_{\substack{f \in \cF }}w_R(f)\lambda_f(p_2)^{r_2}. \notag$$
    Using the triangle inequality and the fact that $|\lambda_f(p_2)|\leq 2$, we find 
    \begin{align}
        |B'_{r_1,r_2}(p_1,p_2)|\ &=\ \frac{1}{p_1^{r_1/2}W_R(\cF)}\sum_{\substack{f \in \cF }}w_R(f)|\lambda_f(p_2)|^{r_2} \leq\frac{2^{r_2}}{p^{r_2/2}}.
    \end{align}
  \qed 

Lastly, we look at the proof of Lemma \ref{lem:psquareinfmom}.
\noindent \textit{Proof of Lemma \ref{lem:psquareinfmom}.}
First, we show that $A'_{r}(p), \ B''_{r_1,r_2}(p_1, p_2)$, and $B'_{r_1,r_2}(p_1, p_2)$ for $N = p^2$ are 0. 
By definition, for $p_1,p_2 \nmid N$, we have that the sum is empty, so $B''_{r_1,r_2}(p_1, p_2) = 0$. For $p_1,~p_2 \mid N$, by \cite{BarrettEtAl2016arXiv} (2.11), we have for $p_i$ a prime,
\begin{align*}
    \lambda_f(p_i)^2 \ = \ \begin{cases}
        \frac{1}{p_i}, & p_i \mid\mid N \\
        0, & p_i^2 \mid N.
    \end{cases}
\end{align*}
Since $p_1,~p_2 \mid N$ the Hecke eigenvalue terms are zero. Thus, $B''_{r_1,r_2}(p_1, p_2) = 0$ for all values of $p_1$ and $p_2$.
Similarly for $B'_{r_1,r_2}(p_1, p_2)$, for $p_1 \nmid N$, the sum is empty, so $B'_{r_1,r_2}(p_1, p_2) = 0$. For $p_1 \mid N$, we again have that the Hecke eigenvalue term $\lambda_f(p_1)$ is zero. Thus, $B'_{r_1,r_2,}(p_1, p_2) = 0$ for all values of $p_1$ and $p_2$. The same reasoning applies to show that $A'_{r}(p)$ for all primes $p$.

Next, we compute $A_r(q)$.
For $q \mid N$, which implies $q = p$, the sum is empty, so $A_r(q) = 0$. Therefore, assume that $q \neq p$.
We then have
\begin{align}
    A_r(p)\ &=\ \frac{1}{W_R(\cF)} \sum_{f\in\cF} w_R(f) \lambda_f(p)^{r} \notag \\& = \frac{1}{W_R(\cF)} \sum_{f\in\cF} w_R(f) \sum_{n=0}^{\lfloor r/2\rfloor }b_{r,r-2n}\lambda_f(p^{r-2n}) \notag \\
    \ &=\ \frac{1}{W_R(\cF)} \sum_{n=0}^{\lfloor r/2\rfloor }b_{r,r-2n}\frac{k-1}{12} \left((p^2-1)\Delta_{k,p^2}\left(q^{r-2n}, 1\right)-p\Delta_{k,p}\left(q^{r-2n}, 1\right)\right) \notag\\
    \ &=\ \sum_{n=0}^{\lfloor r/2\rfloor }b_{r,r-2n}\delta(q^{r-2n},~1) \ + \ O_k\left(\frac{1}{N}2^rq^{r/4}\log{(2q^r)}\right)\notag\\
    \ &=\ \begin{cases}
        C_{r/2} + O_k\left(2^rq^\frac{r}{4} \frac{\log(2q^{r-2\ell})}{N}\right)& r  \ \text{even} \\
       O_k\left(2^rq^\frac{r}{4} \frac{\log(2q^{r-2\ell})}{N}\right) & r  \ \text{odd}.
    \end{cases}
\end{align}

Lastly, we compute $B_{r_1, r_2,}(p_1, p_2)$.
By definition,
\begin{align}\label{B for N=p^2}
    B_{r_1,r_2}(p_1, p_2) \ = \ \frac{1}{W_R(H_k^*(N))}\sum_{f \in H_k^*(N),~p_1\nmid N_f,~p_2 \nmid N_f} w_R(f) \lambda_f(p_1)^{r_1}\lambda_f(p_2)^{r_2}.
\end{align}
For $p_1, p_2 \mid N$, which implies $p_1$ and $ p_2$ equal $p$, the sum is empty, so $B_{r_1,~r_2,~\mathcal{H}_k^*(N)}(p_1, p_2) = 0$. Therefore, assume that $p_1, p_2 \neq p$.
We then have
\begin{align}
   B_{r_1,r_2}(p_1, p_2)\ &=\ \frac{1}{W_R(\mathcal{F})}\sum_{f\in\mathcal{F}}\sum_{k_1=0}^{\lfloor r_1/2 \rfloor}\sum_{k_2=0}^{\lfloor r_2/2 \rfloor} w_R(f)b_{r_1, r_1-2k_1}b_{r_2, r_2-2k_2}\lambda_f\left(p_1^{r_1-2k_1}\right)\lambda_f\left(p_2^{r_2-2k_2}\right) \nonumber\\
    &=\ \frac{1}{W_R(\mathcal{F})}\sum_{k_1=0}^{\lfloor r_1/2 \rfloor}\sum_{k_2=0}^{\lfloor r_2/2 \rfloor} b_{r_1, r_1-2k_1}b_{r_2, r_2-2k_2}\sum_{f\in\mathcal{F}}w_R(f)\lambda_f\left(p_1^{r_1-2k_1}\right)\lambda_f\left(p_2^{r_2-2k_2}\right) \nonumber\\
    &=\ \label{B up to Delta*} \frac{1}{W_R(\mathcal{F})}\sum_{k_1=0}^{\lfloor r_1/2 \rfloor}\sum_{k_2=0}^{\lfloor r_2/2 \rfloor} b_{r_1, r_1-2k_1}b_{r_2, r_2-2k_2}\Delta^*_{k,N}\left(p_1^{r_1-2k_1}, p_2^{r_2-2k_2}\right).
\end{align}

We now apply \cite{BarrettEtAl2016arXiv} Proposition 4.1, which holds as $p_1$ and $p_2$ are not equal $p$. Using the fact that $\mu(p^2)=0$,
\begin{align*}
    \Delta^*_{k,N}\left(p_1^{r_1-2k_1}, p_2^{r_2-2k_2}\right)\ &=\ \frac{k-1}{12}\sum_{LM=N}\mu(L)M\prod_{\substack{q^2 \mid M \\ q \ \textrm{prime}}}\left(\frac{q^2-1}{q^2} \right)\sum_{\substack{\ell \mid L^\infty\\ (\ell, M)=1}}\ell^{-1}\Delta_{k,M}(m\ell^2,n) \\
    \ &=\ \frac{k-1}{12} \left((p^2-1)\Delta_{k,p^2}\left(p_1^{r_1-2k_1}, p_2^{r_2-2k_2}\right)-p\Delta_{k,p}\left(p_1^{r_1-2k_1}, p_2^{r_2-2k_2}\right)\right).
\end{align*}
By \cite{ILS2} Proposition 2.2,
\begin{align*}
    \Delta_{k,p^2}&\left(p_1^{r_1-2k_1}, p_2^{r_2-2k_2}\right) = \delta\left(p_1^{r_1-2k_1}, p_2^{r_2-2k_2}\right) \\ &+O\left( \frac{\tau(p^2)}{p^2k^{5/6}}\frac{\tau_3\left(\left(p_1^{r_1-2k_1}, p_2^{r_2-2k_2}\right)\right)}{\sqrt{2}}\left( \frac{p_1^{r_1-2k_1}p_2^{r_2-2k_2}}{\sqrt{p_1^{r_1-2k_1}p_2^{r_2-2k_2}}+kp^2}\right)^{1/2}\log\left(2p_1^{r_1-2k_1}p_2^{r_2-2k_2}\right)\right) \\
    &=\ \delta(p_1^{r_1-2k_1},p_2^{r_2-2k_2}) \ + \ O_k\left(\frac{1}{p^2}\left( \frac{p_1^{r_1-2k_1}p_2^{r_2-2k_2}}{\sqrt{p_1^{r_1-2k_1}p_2^{r_2-2k_2}}+kp^2}\right)^{1/2}\log\left(2p_1^{r_1-2k_1}p_2^{r_2-2k_2}\right) \right)
\end{align*}
and
\begin{align*}
    \Delta_{k,p}&\left(p_1^{r_1-2k_1}, p_2^{r_2-2k_2}\right) = \delta\left(p_1^{r_1-2k_1}, p_2^{r_2-2k_2}\right)\\ &+O\left( \frac{\tau(p)}{pk^{5/6}}\frac{\tau_3\left(\left(p_1^{r_1-2k_1}, p_2^{r_2-2k_2}\right)\right)}{\sqrt{2}}\left( \frac{p_1^{r_1-2k_1}p_2^{r_2-2k_2}}{\sqrt{p_1^{r_1-2k_1}p_2^{r_2-2k_2}}+kp}\right)^{1/2}\log\left(2p_1^{r_1-2k_1}p_2^{r_2-2k_2}\right)\right) \\
    &=\ \delta\left(p_1^{r_1-2k_1}, p_2^{r_2-2k_2}\right)+O_k\left( \frac{1}{p}\left( \frac{p_1^{r_1-2k_1}p_2^{r_2-2k_2}}{\sqrt{p_1^{r_1-2k_1}p_2^{r_2-2k_2}}+kp}\right)^{1/2}\log\left(2p_1^{r_1-2k_1}p_2^{r_2-2k_2}\right)\right).
\end{align*}
After multiplying $\Delta_{k,p}(p_1, p_2)$ by $p$ and $\Delta_{k,~p^2}(p_1, p_2)$ by $(p^2-1)$, the error term associated with $\Delta_{k,~p}(p_1, p_2)$ absorbs the error terms associated with $\Delta_{k,p^2}(p_1, p_2)$.
As such, we may write
\begin{align}
  B_{r_1,r_2}(p_1, p_2)\ =\ \frac{1}{W_R(\mathcal{F})} \left(\sum_{k_1=0}^{\lfloor r_1/2 \rfloor}\sum_{k_2=0}^{\lfloor r_2/2 \rfloor} b_{r_1, r_1-2k_1}b_{r_2, r_2-2k_2}\left(\frac{k-1}{12}(p^2-p-1)\delta\left(p_1^{r_1-2k_1}, p_2^{r_2-2k_2}\right) \right. \right. \nonumber\\ \left. \ + \ O_k\left(\left(\frac{p_1^{r_1-2k_1}p_2^{r_2-2k_2}}{\sqrt{p_1^{r_1-2k_1}p_2^{r_2-2k_2}}+kp}\right)^{1/2}\log\left(2p_1^{r_1-2k_1}p_2^{r_2-2k_2}\right)\right) \right).
\end{align}
Dividing by $W_R(\mathcal{F})$, we have
\begin{align*}
    B_{r_1,r_2}(p_1, p_2)\ &=\ \sum_{k_1=0}^{\lfloor r_1/2 \rfloor}\sum_{k_2=0}^{\lfloor r_2/2 \rfloor} b_{r_1, r_1-2k_1}b_{r_2, r_2-2k_2}\delta(p_1^{r_1-2k_1}, p_2^{r_2-2k_2}) \\ &+ \ O_k\left(\frac{1}{N} \sum_{k_1=0}^{\lfloor r_1/2 \rfloor}\sum_{k_2=0}^{\lfloor r_2/2 \rfloor} b_{r_1, r_1-2k_1}b_{r_2, r_2-2k_2}\left({p_1^{r_1-2k_1}p_2^{r_2-2k_2}}\right)^{1/4}\log\left(2p_1^{r_1-2k_1}p_2^{r_2-2k_2}\right)\right).
\end{align*}
Trivial estimation and the fact that $|\sum_{\ell=0}^{\lfloor r/2\rfloor}b_{r,r-2\ell}|\lesssim2^{r}$ gives us
\begin{align*}
    B_{r_1,r_2,}(p_1, p_2)\ = \ \sum_{k_1=0}^{\lfloor r_1/2 \rfloor}\sum_{k_2=0}^{\lfloor r_2/2 \rfloor} b_{r_1, r_1-2k_1}b_{r_2, r_2-2k_2}\delta(p_1^{r_1-2k_1}, p_2^{r_2-2k_2})+ \ O_k\left(\frac{1}{N} 2^{r_1+r_2}{p_1^{r_1/4}p_2^{r_2/4}}\log\left(2p_1^{r_1}p_2^{r_2}\right)\right),
\end{align*}
yielding the lemma. \qed
%%%%%%%%%%%%%%%%%%%%%%%%%%%%%%%%%%%%%%%%%%%%%%%%
%%%%%%%%%%%%%%%%%%%%%%%%%%%%%%%%%%%%%%%%%%%%%%%%

\section{Lemmas Used for Main Terms in Theorems \ref{thm:S_B''}, \ \ref{thm:S_B'}, \ \ref{thm:S_Bf},\ and \ref{thm:S_infty}. } \label{appendix Lemmas for main term}
In this section, we state and prove various lemmas used to compute the main Terms in Theorems \ref{thm:S_B''}, \ \ref{thm:S_B'}, \ \ref{thm:S_Bf},\ and \ref{thm:S_infty}.
\begin{lemma}\label{lem:MSwork}
Let \(\hat\phi\) be a compactly supported even Schwartz test function. As in \cite{young2005lower}, define
\[
\theta(t) := \sum_{p\leq t}\log(p), 
\quad E(t) :=\theta(t)-t,
\quad S(R) := \sum_{p} \frac{2\log(p)}{p\log(R)} \hat\phi\Bigl(\frac{2\log(p)}{\log(R)}\Bigr).
\]
Then
\begin{equation}
\begin{aligned}
S(R) & \ = \ \frac{\phi(0)}{2}
+\frac{2\hat\phi(0)}{\log(R)}\Bigl(1+\int_{1}^{\infty}\frac{E(t)}{t^{2}}dt\Bigr)\\
&+\frac{4\hat\phi''(0)}{(\log(R))^{3}}
\int_{1}^{\infty}\frac{E(t)}{t^{2}}\Bigl((\log(t))^{2}-2\log(t)\Bigr)dt
+O\Bigl(\frac{1}{\log^4(R)}\Bigr).
\end{aligned}
\end{equation}
\end{lemma}
\begin{proof}
By Abel summation
\begin{equation}
\begin{aligned}
S(R)
& \ = \ \sum_{p}\frac{2\log(p)}{p\log(R)}\hat\phi\Bigl(\frac{2\log(p)}{\log(R)}\Bigr)
=\frac{2}{\log(R)}\sum_{p}\frac{\log(p)}{p}\hat\phi\Bigl(\frac{2\log(p)}{\log(R)}\Bigr)\\
&\ = \ \frac{2}{\log(R)}\lim_{x\to\infty}\Bigl[\theta(x)\frac{1}{x}\hat\phi\Bigl(\frac{2\log(x)}{\log(R)}\Bigr)
-\int_{1}^{x}\theta(t)\frac{d}{dt}\Bigl(\frac{1}{t}\hat\phi\Bigl(\frac{2\log(t)}{\log(R)}\Bigr)\Bigr)dt\Bigr]\\
& \ = \ -\frac{2}{\log(R)}\int_{1}^{\infty}\theta(t)\frac{d}{dt}\Bigl(\frac{1}{t}\hat\phi\Bigl(\frac{2\log(t)}{\log(R)}\Bigr)\Bigr)dt.
\end{aligned}
\label{l1theta}
\end{equation}
Define \(u=\frac{2\log(t)}{\log(R)}\). Then
\begin{equation}
\frac{d}{dt}\Bigl(\frac{1}{t}\hat\phi(u)\Bigr)
\ = \ -\frac{1}{t^{2}}\hat\phi(u)+\frac{2}{(\log(R))t^{2}}\hat\phi'(u).
\end{equation}
Write $\theta$ as $t + E(t)$. The first term in \eqref{l1theta} is
\begin{equation}
\frac{2}{\log(R)}\int_{1}^{\infty}\frac{1}{t}\Bigl(\hat\phi\Bigl(\frac{2\log(t)}{\log(R)}\Bigr)
-\frac{2}{\log(R)}\hat\phi'\Bigl(\frac{2\log(t)}{\log(R)}\Bigr)\Bigr)dt
\end{equation}
\begin{equation}
\begin{aligned}
& \ = \ \frac{2}{\log(R)}\int_{1}^{\infty}\frac{1}{t}\hat{\phi}\Bigl(\frac{2\log(t)}{\log(R)}\Bigr)dt-\frac{4}{(\log(R))^{2}}\int_{1}^{\infty}\frac{1}{t}\hat\phi'\Bigl(\frac{2\log(t)}{\log(R)}\Bigr)dt\\
& \ = \ \frac{2}{\log(R)}\cdot\frac{\log(R)}{2}\int_{0}^{\infty}\hat\phi(u)du-\frac{4}{(\log(R))^{2}}\cdot\frac{\log(R)}{2}\int_{0}^{\infty}\hat\phi'(u)du\\
& \ = \ \frac{2}{\log(R)}\cdot\frac{\log(R)}{2}\cdot\frac{\phi(0)}{2}-\frac{4}{(\log(R))^{2}}\cdot\frac{\log(R)}{2}\cdot\bigl(-\hat\phi(0)\bigr)\\
&  \ = \ \frac{\phi(0)}{2}+\frac{2\hat\phi(0)}{\log(R)}. \label{l1first}
\end{aligned}
\end{equation}
The second term in \eqref{l1theta} is
\begin{equation}
\frac{2}{\log(R)}\int_{1}^{\infty}\frac{E(t)}{t^{2}}\Bigl(\hat\phi\Bigl(\frac{2\log(t)}{\log(R)}\Bigr)
-\frac{2}{\log(R)}\hat\phi'\Bigl(\frac{2\log(t)}{\log(R)}\Bigr)\Bigr)dt.
\end{equation}
Because $\hat\phi$ is even, the Taylor expansions are
\begin{equation}
\hat\phi(u) \ = \ \hat\phi(0)+\frac{\hat\phi''(0)}{2}u^{2}+O(u^{4}), \quad
\hat\phi'(u) \ = \ \hat\phi''(0)u+O(u^3).
\end{equation}
Thus, the second term in \eqref{l1theta} becomes
\begin{equation}
\begin{aligned}
& \ = \ \frac{2}{\log(R)}\int_{1}^{\infty}\frac{E(t)}{t^{2}}
\Bigl[\hat\phi(0)
      +\frac{\hat\phi''(0)}{2}\Bigl(\frac{2\log(t)}{\log(R)}\Bigr)^{2}
      -\frac{2\hat\phi''(0)}{\log(R)}\Bigl(\frac{2\log(t)}{\log(R)}\Bigr)
      +O\Bigl(\Bigl(\frac{\log(t)}{\log(R)}\Bigr)^{3}\Bigr)
\Bigr]dt \\
& \ = \  \frac{2\hat\phi(0)}{\log(R)}\int_{1}^{\infty}\frac{E(t)}{t^{2}}dt
+\frac{4\hat\phi''(0)}{(\log(R))^{3}}\int_{1}^{\infty}\frac{E(t)(\log(t))^{2}}{t^{2}}dt
-\frac{8\hat\phi''(0)}{(\log(R))^{3}}\int_{1}^{\infty}\frac{E(t)\log(t)}{t^{2}}dt
+O\Bigl(\frac{1}{\log^4(R)}\Bigr). \label{l1second}
\end{aligned}
\end{equation}
Combining the terms in \eqref{l1first} and \eqref{l1second} yields the lemma.
\end{proof}

\begin{lemma}
Let $\hat\phi_1, \ \hat\phi_2$ be a compactly supported even Schwartz test function. Then
\[
\begin{aligned}
&\sum_p \frac{\log^2(p)}{p\log^2(R)} \hat\phi_1\left(\frac{\log(p)}{\log(R)}\right) \hat\phi_2\left(\frac{\log(p)}{\log(R)} \right)\\
&= \int_0^\infty u\hat\phi(u)du -\frac{\hat\phi(0)}{\log^2(R)}\int_1^{\infty}\frac{E(t)}{t^2}(1-\log(t))dt + O\left(\frac{1}{\log^4(R)}\right),
\quad \text{where} \quad \phi:=\phi_1*\phi_2.
\end{aligned}
\]
\end{lemma}
\begin{proof}
By the Convolution Theorem, we have
\begin{equation}
\hat\phi_1\left(\frac{\log(p)}{\log(R)}\right)
\hat\phi_2\left(\frac{\log(p)}{\log(R)}\right)
=\hat\phi\left(\frac{\log(p)}{\log(R)}\right),
\quad \phi:=\phi_1*\phi_2.
\end{equation}
Define
\[
f(t) := \frac{\log(t)}{t\log^2(R)}\hat\phi\left(\frac{\log(t)}{\log(R)}\right),
\quad \theta(t) := t + E(t).
\]
\[
S_1(R) := \sum_{p} \frac{\log^2(p)}{p\log^2(R)} \hat\phi\left(\frac{\log(p)}{\log(R)}\right) = \sum_p \log(p)f(p).
\]
By Abel summation, we write
\begin{equation}
S_1(R) \ = \ \lim_{x \to \infty}\theta(x)f(x) -\int_1^{\infty}\theta(t)f'(t)dt
= -\int_1^{\infty}\theta(t)f'(t)dt. \label{l2sum}
\end{equation}
Differentiate $f$ with respect to $u := \frac{\log(t)}{\log(R)}$. Then we change the variable and calculate \eqref{l2sum}
\begin{equation}
f'(t) \ = \ \frac{1}{\log^2(R)}\left(\frac{1-\log(t)}{t^{2}}\hat\phi(u) 
+ \frac{\log(t)}{t^2 \log(R)} \hat\phi'(u)\right),
\end{equation}
\begin{equation}
S_1(R) \ = \ \int_0^\infty \theta(t)\left(\frac{\log(t) - 1}{t^2\log^2(R)} \hat\phi(u) - \frac{1}{\log(R)} \frac{\log(t)}{t^2 \log^2(R)}\hat\phi'(u)\right)dt. \label{l2sumc}
\end{equation}
Write $\theta$ as $t + E(t)$, by change of variable, the first term in \eqref{l2sumc} is
\begin{equation}
\int_0^\infty u\hat\phi(u)du - \frac{1}{\log(R)} \int_0^\infty \left(\hat\phi(u) + u\hat\phi'(u)\right)du
\ = \ \int_0^\infty u\hat\phi(u)du.
\end{equation}
Since $\hat\phi$ is even. The Taylor expansions are
\begin{equation}
\hat\phi(u) \ = \ \hat\phi(0)+O(u^2), \quad
\hat\phi'(u) \ = \ \hat\phi''(0)u+O(u^3).  \label{l2taylor}
\end{equation}
Substituting \eqref{l2taylor} into $f'(t)$ yields
\begin{equation}
f'(t) \ = \ \frac{\hat\phi(0)}{t^{2}\log^2(R)}(1-\log(t))
+O\Bigl(\frac{\log^2t}{t^{2}\log^4(R)}\Bigr).
\label{l2firstt}
\end{equation}
The second term is then equal to
\begin{equation}
-\frac{\hat\phi(0)}{\log^2(R)}\int_1^{\infty}\frac{E(t)}{t^2}(1-\log(t))dt + O\left(\frac{1}{\log^4(R)}\right).
\label{l2secondt}
\end{equation}
Combining \eqref{l2firstt} and \eqref{l2secondt} yields the lemma.

\end{proof}

\begin{lemma}\label{lem:S_2(R)}
    Let $\phi_1$, $\phi_2$ be even Schwartz functions with $\widehat \phi_1$ and $\widehat \phi_2$ compactly supported. Define
$$S_2(R):=\sum_{p}\frac{\log^2(p)}{p^2\log^2(R)}\hat\phi_1\left(\frac{2\log(p)}{\log(R)}\right)\,\,\hat\phi_2\left(\frac{2\log(p)}{\log(R)}\right),$$
    then we have
\[
    S_2(R) \ = \ \frac{\hat\phi(0)}{\log^2(R)}-\frac{\hat \phi(0)+\hat\phi''(0)}{\log^2(R)}\int_1^\infty \frac{E(t)(1-2\log(t))}{t^3}dt+O\left(\frac{1}{\log^4(R)}\right).
\]
\end{lemma}
\begin{proof}
As before, we appeal to the convolution theorem, writing $\phi := \phi_1*\phi_2$ so that
\begin{equation}
\sum_{p}\frac{\log^2 (p)}{p^2\log^2(R)}\hat\phi_1\left(\frac{2\log(p)}{\log(R)}\right)\,\,\hat\phi_2\left(\frac{2\log(p)}{\log(R)}\right)\,=\,\sum_p \frac{\log^2(p)}{p^2\log^2(R)}\hat\phi\left(\frac{2\log(p)}{\log(R)}\right).
\end{equation}
Recall that we define $\theta(t):=\sum_{p\leq t} \log(p)$ and let
$$f(x) \ := \ \frac{\log(x)}{x^2}\hat\phi\left(\frac{2\log(x)}{\log(R)}\right).$$
Now, by the Abel summation formula and the fact that $\mathrm{supp}(\hat\phi)$ is bounded,
\begin{equation}
S_2(R) \ = \ \frac{1}{\log^2(R)}\lim_{x\to \infty}\left[ \theta(x) f(x)-\int_1^x \theta(t)f'(t)dt\right]=-\frac{1}{\log^2(R)}\int_1^\infty \theta(t) f'(t)dt.
\end{equation}
To ease notation, substitute $u=\frac{2\log(t)}{\log(R)}$ and differentiate to get
\begin{equation}
f'(t)=\frac{1-2\log(t)}{t^3}\hat\phi(u)+\frac{2\log(t)}{t^3\log(R)}\hat\phi'(u)
\end{equation}
\begin{equation}
\int_1^\infty t f'(t)dt \ = \ -\int_1^\infty f(t)dt.
\end{equation}
Splitting the error in $\theta(t)=t+E(t),$ we first consider the main term. Substituting in $u$ again gives
\begin{equation}
\int_1^\infty f(t)dt \ = \ \frac{\log ^2 (R)}{4}\int_0^\infty  u e^\frac{-u\log(R)}{2}\hat\phi(u)du.
\end{equation}
Now, using the evenness of $\hat\phi$, we Taylor expand and find
\begin{equation}
\hat\phi(u) \ = \ \hat\phi(0)+\frac{\hat\phi''(0)}{2}u^2+O(u^4).
\end{equation}
and note that $\int_0^\infty ue^{-\alpha u} \ = \ 1/\alpha^2$ and $\int_0^\infty u^3e^{-\alpha u} \ = \ 6/\alpha^4.$ Thus, the main term is, after bringing in the factor of $-1/\log^2(R)$,
\begin{equation}
\frac{\hat\phi(0)}{\log^2(R)}+\frac{12 \hat\phi''(0)}{\log^4(R)}+O\left(\frac{1}{\log^6(R)}\right) \ = \ \frac{\hat\phi(0)}{\log^2(R)}+O\left(\frac{1}{\log^4(R)}\right).
\end{equation}
We now compute the error term
\begin{equation}
    -\frac{1}{\log^2(R)}\int_1^\infty E(t) f'(t)dt \ = \ -\frac{1}{\log^2(R)}\int_1^\infty E(t) \left(\frac{1-2\log(t)}{t^3}\hat\phi(u)+\frac{2\log(t)}{t^3\log(R)}\hat\phi'(u)\right)dt.
\end{equation}
Since we also have from evenness that
\begin{equation}
\hat{\phi}'(u) \ = \ \hat{\phi}''(0)u+O(u^3).
\end{equation}
we expand each of $\hat \phi $ and $\hat\phi'$ to see that the above is equal to
\begin{align}
    -\frac{1}{\log^2(R)}\int_1^\infty E(t)\left(\frac{1-2\log(t)}{t^3}\hat\phi(0)+\left(\frac{2\log^2(t)}{t^3\log^2(R)}+\frac{1-2\log(t)}{t^3}\right)\hat\phi''(0)+O\left(\frac{\log^3 (t)}{\log^3(R)}\right)\right)\nonumber\\
    =-\frac{\hat \phi(0)+\hat\phi''(0)}{\log^2(R)}\int_1^\infty \frac{E(t)(1-2\log(t))}{t^3}dt+O\left(\frac{1}{\log^4(R)}\right).
\end{align}
\end{proof}


\begin{lemma}
      Let $p$ be prime. For $\alpha < 1$ and $p$ prime, we have
      \[
      \sum_{p\le x}\frac{\log(p)}{p^\alpha} 
      \ = \ \frac{x^{1-\alpha}}{1-\alpha} + O\left(\frac{x^{1-\alpha}}{\log(x)}\right).
      \]
\end{lemma}
\noindent
\begin{proof}
Define
\[
S(x):=\sum_{p\leq x}\frac{\log(p)}{p^\alpha}, \quad \theta(t) := t + E(t).
\]
By Abel summation
\begin{equation}
\begin{aligned}
S(x) 
&\ = \  x^{-\alpha}\theta(x)-2^{-\alpha}\theta(2)+\alpha\int_{2}^{x}\theta(t) t^{-\alpha-1} dt.\\
&\ = \  x^{-\alpha}\bigl(x+E(x)\bigr)-2^{-\alpha}\log (2)
+\alpha\int_{2}^{x} \bigl(t+E(t)\bigr) t^{-\alpha-1} dt\\
&\ = \  x^{1-\alpha} + \alpha\int_{2}^{x} t^{-\alpha} dt
-2 ^{-\alpha}\log (2)
+ x^{-\alpha}E(x) + \alpha\int_{2}^{x} E(t) t^{-\alpha-1} dt.
\end{aligned}
\end{equation}
After calculating the main terms, we have
\begin{equation}
S(x)\ = \ \frac{x^{1-\alpha}}{1-\alpha}
+x^{-\alpha}E(x)
+\alpha\int_{2}^{x} E(t) t^{-\alpha-1} dt
-2^{-\alpha}\log (2)-\frac{\alpha 2^{1-\alpha}}{1-\alpha}. \label{l4ppower}
\end{equation}
By the Prime Number Theorem, $E(t)=\theta(t)-t=O\bigl(t/ \log(t)\bigr)$, substituting in \eqref{l4ppower} yields
\begin{equation}
x^{-\alpha}E(x)\ = \ \Bigl(\frac{x^{1-\alpha}}{\log(x)}\Bigr),\qquad
\int_{2}^{x} E(t) t^{-\alpha-1}dt\ = \ O\Bigl(\int_{2}^{x}\frac{t^{-\alpha}}{\log(t)}dt\Bigr)
\ = \ O\left(\frac{x^{1-\alpha}}{\log(x)}\right).
\end{equation}
The constants are absorbed into the big O term, gives
\begin{equation}
S(x) \ = \  \frac{x^{1-\alpha}}{1-\alpha}+O\left(\frac{x^{1-\alpha}}{\log(x)}\right).
\end{equation}
\end{proof}
\begin{lemma}
    Let $\widehat\phi$ be an even Schwartz function, then
    \begin{equation}
    \begin{aligned}
    \sum_p\hat\phi\left(\frac{\log(p)}{\log(R)}\right)\frac{\log^2(p)}{\log^2(R)}&\left[\frac{-3p-1}{p(p+1)^2}+\sum_{i=2}^\infty C_i \frac{p^{i-1}(p-1)}{(p+1)^{2i}}\right]\\
    &\ = \ -\frac{\hat\phi(0)}{\log^2(R)}+\frac{\hat \phi(0)+\hat\phi''(0)}{\log^2(R)}\int_1^\infty \frac{E(t)(1-2\log(t))}{t^3}dt+O\left(\frac{1}{\log^4(R)}\right).
    \end{aligned} \label{l4sum}
    \end{equation}
\end{lemma}
\begin{proof}
As in \cite{Sl}, the generating function for the Catalan numbers can be written in closed-form for $|z|<1/4$:
\begin{equation}\label{catSum}
    \sum_{i=0}^\infty C_i z^i \ = \frac{1-\sqrt{1-4z}}{2z}.
\end{equation}
Letting $z=p/(p+1)^2$, we note that
\begin{equation}
    |z|\leq \frac{2}{(2+1)^2}=\frac{2}{9}<\frac{1}{4}
\end{equation}
so we always have convergence to the above formula. Factoring and subtracting off the first two terms,
\begin{equation}
\sum_{i=2}^\infty C_i\frac{p^{i-1}(p-1)}{(p+1)^{2i}} \ = \ \frac{p-1}{p}\left(\sum_{i=0}^\infty C_iz^i-C_0-C_1z\right).
\end{equation}
Since 
\begin{equation}
1-4z \ = \ \frac{(p+1)^2}{(p+1)^2}-\frac{4p}{(p+1)^2}=\left(\frac{p-1}{p+1}\right)^2,
\end{equation}
we have, using $C_0=C_1=1$,
\begin{equation}
\sum_{i=0}^\infty C_i z^i-C_0-C_1z \ = \ \frac{1-\sqrt{1-4z}}{2z}-1-z=\frac{p+1}{p}-1-z \ = \ \frac{1}{p}-z.
\end{equation}
We multiply through by the extra factors to get
\begin{equation}
\sum_{i=2}^\infty C_i \frac{p^{i-1}(p-1)}{(p+1)^{2i}} \ = \ \frac{p-1}{p}\left(\frac{1}{p}-z\right) \ = \ \frac{p-1}{p^2}-\frac{p-1}{p}z \ = \ \frac{p-1}{p^2}-\frac{p-1}{(p+1)^2}.
\end{equation}
Furthering our effort to simplify the bracket in the sum, note that
\begin{equation}
\frac{-3p-1}{p(p+1)^2}+\frac{p-1}{p^2}-\frac{p-1}{(p+1)^2}\ = \frac{-3p^2-p+(p-1)(p+1)^2-(p-1)p^2}{p^2(p+1)^2}\ = \ -\frac{1}{p^2}.
\end{equation}
hence, the sum becomes
\begin{equation}
-\sum_p\hat\phi_2\left(\frac{\log(p)}{\log(R)}\right) \frac{\log^2(p)}{p^2\log ^2 (R)}
\end{equation}
which is amenable to our techniques involving the Prime Number Theorem. Let $u=\log(t)/\log(R)$ and apply the Abel summation formula with the sequence $a_p=\log(p)$ to get the main term
\begin{equation}
    -\int_1^\infty \frac{\log(t)}{t^2\log^2(R)}\hat\phi_2\left(\frac{\log(t) }{\log(R)}\right)dt \ = \ -\int_0^\infty u e^{-u\log(R)} \hat\phi_2(u)du
\end{equation}
arising from $\theta(t)=t+E(t).$ Writing $\hat\phi(u)=\hat\phi(0)+O(u^2)$ and plugging back into the integral, the main term becomes
\begin{equation}
    -\frac{\hat\phi(0)}{\log^2(R)}+O\left(\frac{1}{\log^4(R)}\right)
\end{equation}
since $\int_0^\infty ue^{-u\alpha}du=1/\alpha^2.$
We proceed to compute the error term. We have that the error is
\begin{equation}
    \frac{1}{\log^2(R)}\int_1^\infty E(t)\dfrac{d}{dt}\left[ \frac{\log^2 t}{t^2}\phi\left(\frac{\log(t)}{\log(R)}\right)\right].
\end{equation}
Now,
\begin{equation}
    \dfrac{d}{dt}\left[ \frac{\log^2 t}{t^2}\phi\left(\frac{\log(t)}{\log(R)}\right)\right] \ = \ \frac{1-2\log(t)}{t^3}\hat\phi(u)+\frac{\log(t)}{t^3\log(R)}\hat\phi'(u)
\end{equation}
so the error term is 
\begin{equation}
    \frac{\hat \phi(0)+\hat\phi''(0)}{\log^2(R)}\int_1^\infty \frac{E(t)(1-2\log(t))}{t^3}dt+O\left(\frac{1}{\log^4(R)}\right).
\end{equation}
\end{proof}

\begin{lemma}\label{line4}
    Let $\hat\phi$ be an even Schwartz function, then
    \[
        \begin{aligned}
        \sum_p \frac{\log^2(p)}{\log^2(R)}\hat\phi_2\left(\frac{2\log(p)}{\log(R)}\right) &\frac{p^2+3p+1}{p^2(p+1)^3}\\
        & \ = \ \frac{I\hat\phi(0)}{4\log^2(R)}-\frac{\hat\phi(0)}{\log^2(R)}\int_1^\infty \frac{E(t)B(t)}{t^3(t+1)^4}dt+O\left(\frac{1}{\log^4(R)}\right),
        \end{aligned}
    \]
    where quantities $I$ and $B(t)$ are defined explicitly in the proof below.
\end{lemma}
\begin{proof}
    Let $f(x)=\log(x)\,\hat\phi_2\left(\frac{2\log(x)}{\log(R)}\right) \frac{x^2+3x+1}{x^2(x+1)^3}$ so that by our standard tricks, we have
    \begin{equation}\label{line4error}
        \sum_p \frac{\log^2(p)}{\log^2(R)}\hat\phi_2\left(\frac{2\log(p)}{\log(R)}\right) \frac{p^2+3p+1}{p^2(p+1)^3}\ = \ -\frac{1}{\log^2(R)}\int_1^\infty tf'(t)dt-\frac{1}{\log^2(R)}\int_1^\infty E(t)f'(t)dt.
    \end{equation}
    Integrating the first term by parts and letting $s=2\log(t)$, it suffices to compute
    \begin{equation}
    \int_1^\infty f(t)dt \ = \ \frac{1}{4}\int_0^\infty s \hat\phi\left(\frac{s}{\log(R)}\right) \frac{e^s+3+e^{-s}}{(e^s+1)^3}ds.
    \end{equation}
    Expanding $\hat\phi$ using evenness, this gives that the main term is
    \begin{equation}
    \frac{I\hat\phi(0)}{4\log^2(R)}+O\left(\frac{1}{\log^4(R)}\right),
    \end{equation}
where we let 
    \begin{equation}
        I \ := \ \int_0^\infty s\frac{e^s+3+e^{-s}}{(e^s+1)^3}ds
    \end{equation}
to later be integrated numerically. % Pramana: Good :)
We now compute the error term in \eqref{line4error}. Differentiating, we have
    \begin{equation}
        f'(x) \ = \ \frac{A(x)\hat\phi'\left(\frac{2\log(x)}{\log(R)}\right)+\log(R)\,\,B(x)\hat\phi\left(\frac{2\log(x)}{\log(R)}\right)}{x^3(x+1)^4\log(R)}
    \end{equation}
    where 
    \begin{align}
        &A(x)\ \coloneq \ (2x^3+8x^2+8x+2)\log(x)\\
        &B(x)\ \coloneq \ (-3x^3-12x^2-8x-2)\log(x)+x^3+4x^2+4x+1.
    \end{align}
    Using evenness to expand $\hat\phi$ and $\hat\phi'$ to compute the integral
    \begin{equation}
        -\frac{1}{\log^2(R)}\int_1^\infty E(t)f'(t)dt,
    \end{equation}
    we see that the first term is absorbed into the $O(\log^{-4}R)$ error. 
    Considering the second term, only the constant term in the expansion of $\hat\phi$ matters for us, yielding an error term of
    \begin{equation}
        -\frac{\hat\phi(0)}{\log^2(R)}\int_1^\infty \frac{E(t)B(t)}{t^3(t+1)^4}dt. 
    \end{equation} % Pramana: Good :)
\end{proof}
\begin{lemma}\label{lem:line1}
    Let $\hat\phi$ be even Schwartz function. Then
    \begin{equation}
    \begin{aligned}
    &\sum_{p_1,p_2}\frac{\log(p_1)\log(p_2)}{p_2\log^2(R)}\hat\phi\left(\frac{2\log(p_2)}{\log(R)}\right)\frac{2}{p_1(p_1+1)}\\
    &\quad =\frac{\left(-\frac{1}{2}+\log 2\right)\hat\phi(0)}{2\log(R)}\left(1+4F_1\right)+\frac{\hat\phi(0)F_2}{2\log(R)}+\frac{2\hat\phi(0)F_1F_2}{\log(R)}+O\left(\frac{1}{\log^4(R)}\right).
    \end{aligned}
    \end{equation}
\end{lemma}
\begin{proof}
    Rearranging
    \begin{equation}
    \begin{aligned}
        &\sum_{p_1,p_2}\frac{\log(p_1)\log(p_2)}{p_2\log^2(R)}\hat\phi\left(\frac{2\log(p_2)}{\log(R)}\right)\frac{2}{p_1(p_1+1)}\\
        \ = \ &\sum_{p_1} \frac{\log(p_1)}{p_1(p_1+1)\log(R)}
            \sum_{p_2}\frac{2\log(p_2)}{p_2\log(R)}\hat\phi\left(\frac{2\log(p_2)}{\log(R)}\right).
    \end{aligned}
    \end{equation}
We apply Lemma \ref{lem:MSwork} to the sum over $p_2$, getting that the above equals
\begin{equation}
    \frac{S(R)}{\log(R)}\sum_{p_1} \frac{\log(p_1)}{p_1(p_1+1)}
\end{equation}
where we can compute $S(R)$ up to $O(\log^{-4}R)$ numerically, leaving only the consideration of the sum over $p_1.$ We apply the standard method of Abel summation, setting $g(x)=\frac{1}{x(x+1)}$ so that
\begin{equation}\label{line1split}
    \sum_{p_1}\frac{\log(p_1)}{p_1(p_1+1)}\ =\ -\int_1^\infty tg'(t)dt -\int_1^\infty E(t)g'(t)dt.
\end{equation}
Integrating by parts, the first term gives a contribution of
\begin{equation}
    [-tg(t)]_1^\infty+\int_1^\infty g(t)dt \ = \ -\frac{1}{2}+\log 2.
\end{equation}
Since
\begin{equation}
    g'(x) \ = \ -\frac{2x+1}{(x^2+x)^2},
\end{equation}
the error term in \eqref{line1split}
\begin{equation}
    \int_1^\infty \frac{E(t)(2t+1)}{(t^2+t)^2}dt \eqqcolon F_1 \label{llF1}
\end{equation}
which is a constant that we can compute numerically. Hence, after substituting in our expression for $S(R)$ from Lemma \ref{lem:MSwork}, we have in total
\begin{equation}
\begin{aligned}
    &\sum_{p_1,p_2}\frac{\log(p_1)\log(p_2)}{p_2\log^2(R)}\hat\phi\left(\frac{2\log(p_2)}{\log(R)}\right)\frac{2}{p_1(p_1+1)} \ = \ \frac{S(R)}{\log(R)}\left[-\frac{1}{2}+\log(2)+F_1\right]\\
    &=\bigg[\frac{\hat\phi(0)}{2\log(R)}+\frac{2\hat\phi(0)}{\log(R)}\left(1+\int_0^\infty\frac{E(t)}{t^2}dt\right)+O\left(\frac{1}{\log^4(R)}\right)\bigg]
    \cdot \left[\left(-\frac{1}{2}+\log(2)\right)+\int_1^\infty \frac{E(t)(2t+1)}{(t^2+t)^2}dt \right]. \label{l7sum}
\end{aligned}
\end{equation}
Let 
\begin{align}
    F_2\coloneqq \int_1^\infty \frac{E(t)(2t+1)}{(t^2+t)^2}dt. \label{llF2}
\end{align}
Substituting \eqref{llF1} and \eqref{llF2} into  \eqref{l7sum} yeilds the lemma.
\end{proof}

\begin{lemma}\label{line2}
    Let $\hat\phi$ be even Schwartz function. Then
    \begin{equation}
        \begin{aligned}
        & \sum_{p_1,p_2} \frac{\log(p_1)\log(p_2)}{p_2\log^2(R)}\hat\phi\left(\frac{2\log(p_2)}{\log(R)}\right)\frac{(p_1^2+3p_1+1)}{p_1(p_1+1)^3}\\
        %& \ = \ \bigg[\frac{\hat\phi_2(0)}{4\log(R)}+\frac{F_1\hat\phi_2(0)}{\log(R)}+O\left(\frac{1}{\log^4(R)}\right)\bigg]\cdot
        %\left[\log{2} + \frac{3}{4}+F_3\right]\\
        &=\frac{\hat\phi(0)(\log 2+\frac{3}{4}+F_3)}{\log(R)}\left(\frac{1}{4}+F_1\right)+O\left(\frac{1}{\log^4(R)}\right).
        \end{aligned}
    \end{equation}
    where $F_1=1+\int_0^\infty E(t)/t^2 dt.$ 
\end{lemma}
\begin{proof}
    We begin the same as in Lemma \ref{lem:line1}, obtaining
    \begin{equation}
        \sum_{p_1,p_2} \frac{\log(p_1)\log(p_2)}{p_2\log^2(R)}\hat\phi_2\left(\frac{2\log(p_2)}{\log(R)}\right)\frac{(p_1^2+3p_1+1)}{p_1(p_1+1)^3}\ =\ \frac{S(R)}{2\log(R)}\left[\sum_{p_1} \log(p_1) \frac{p_1^2+3p_1+1}{p_1(p_1+1)^3}\right].
    \end{equation}
    We define $h(x)=\frac{x^2+3x+1}{x(x+1)^3}$ so that by the Prime Number Theorem it follows that
    \begin{equation}
        \sum_{p_1} \log(p_1) h(p) \ = \ -\int_1^\infty th'(t)dt-\int_1^\infty E(t)h'(t)dt.
    \end{equation}
    The first term is
    \begin{equation}
        -[t\cdot h(t)]_1^\infty +\int_1^\infty h(t)dt \ =\ \log 2 + \frac{3}{4}.
    \end{equation}
    % {\color{magenta} Pramana: I get \(\frac{5}{8} + \frac{3}{8} = 1\). 
    % \(th(t)\to 0\) as \(t\to \infty\) because it's \(O(x^{-2})\) and \(h(1) = \frac{5}{8}\). 
    % Source for integral: 
    % \url{https://www.wolframalpha.com/input?i=integral+1+to+infinity+%28x%5E2%2B3x%2B1%29%2F%28x%5E2%28x%2B1%29%5E3%29+}}
    Since
    \begin{equation}
        h'(x)\ = \ -\frac{2x^3+8x^2+4x+1}{x^2(x+1)^4},
    \end{equation}
    the second term is
    \begin{equation}
        F_3 \ :=\ \int_1^\infty E(t)\frac{2t^3+8t^2+4t+1}{t^2(t+1)^4}dt.
    \end{equation}
    Hence, applying Lemma \ref{lem:MSwork} to get the value of $S(R)$, we obtain
        \begin{equation}
        \begin{aligned}
        & \sum_{p_1,p_2} \frac{\log(p_1)\log(p_2)}{p_2\log^2(R)}\hat\phi\left(\frac{2\log(p_2)}{\log(R)}\right)\frac{(p_1^2+3p_1+1)}{p_1(p_1+1)^3}\\
        & \ = \ \bigg[\frac{\hat\phi(0)}{4\log(R)}+\frac{F_1\hat\phi(0)}{\log(R)}+O\left(\frac{1}{\log^4(R)}\right)\bigg]\cdot
        \left[\log{2} + \frac{3}{4}+F_3\right]\\
        &=\frac{\hat\phi(0)(\log 2+\frac{3}{4}+F_3)}{\log(R)}\left(\frac{1}{4}+F_1\right)+O\left(\frac{1}{\log^4(R)}\right),
        \end{aligned}
    \end{equation}
where $F_1=1+\int_0^\infty E(t)/t^2 dt.$ 
\end{proof}
\begin{lemma}
     Let $\hat\phi$ be even Schwartz function. Then
    \begin{equation}
        \begin{aligned}
        & \sum_{p_1,p_2} \frac{\log(p_1)\log(p_2)}{p_2\log(R)}\hat\phi\left(\frac{2\log(p_2)}{\log(R)}\right) \sum_{\ell=2}^\infty C_\ell \frac{p_1^\ell(p_1-1)}{(p_1+1)^{2\ell +1}}\\
        %& \ = \ \bigg[\frac{\hat\phi_2(0)}{4\log(R)}+\frac{F_1\hat\phi_2(0)}{\log(R)}+ O\left(\frac{1}{\log^4(R)}\right)\bigg]\cdot\left[\log 2 -\frac{1}{4}+F_4\right]\\
        & \ = \ \frac{\hat\phi(0)(\log(2)+\frac{1}{4}+F_4)}{\log(R)}\left(\frac{1}{4}+F_1\right)+O\left(\frac{1}{\log^4(R)}\right). 
    \end{aligned}
\end{equation}
\end{lemma}
\begin{proof}
    Applying the same argument as in the previous computations and the fact that
    \begin{equation}
        \sum_{\ell =2}^\infty C_\ell \frac{p_1^\ell (p_1-1)}{(p_1+1)^{2\ell+1}}=\frac{2p_1^2-p_1+1}{p_1(p+1)^3},
    \end{equation}
    it follows that the sum equals
    \begin{equation}\label{line3split}
        \frac{S(R)}{2\log(R)}\left[\sum_{p_1}\log(p_1) \frac{2p_1^2-p_1+1}{p_1(p+1)^3}\right].
    \end{equation}
By the Prime Nnumber Theorem, letting $f(x) \ = \ -\frac{2x^2-x+1}{x(x+1)^3}$ gives
\begin{equation}\label{p1sumline3}
    \sum_{p_1}\log(p_1) \frac{2p_1^2-p_1+1}{p_1(p+1)^3} \ =\ -\int_1^\infty tf'(t)dt-\int_1^\infty E(t)f'(t)dt.
\end{equation}
The first term is
\begin{equation}
    -\int_1^\infty tf'(t)dt \ = \ -[tf(t)]_1^\infty+\int_1^\infty f(t)dt-1=\log(2)+\frac{1}{4}.
\end{equation}
For the second term, note that
\begin{equation}
    f'(x) \ =\ -\frac{4x^{3} - 5x^{2} + 4x + 1}{x^{2} \left(x + 1\right)^{4}}
\end{equation}
so that
\begin{equation}
    -\int_1^\infty E(t)f'(t)dt \ = \ \int_1^\infty E(t)\frac{4x^{3} - 5x^{2} + 4x + 1}{x^{2} \left(x + 1\right)^{4}}dt=:F_4.
\end{equation}
Hence, \eqref{p1sumline3} is equal to
\begin{equation}
    \log (2)+\frac{1}{4}+F_4.
\end{equation}
Substituting into \eqref{line3split}, the sum is
    \begin{equation}
        \begin{aligned}
        & \sum_{p_1,p_2} \frac{\log(p_1)\log(p_2)}{p_2\log(R)}\hat\phi_2\left(\frac{2\log(p_2)}{\log(R)}\right) \sum_{\ell=2}^\infty C_\ell \frac{p_1^\ell(p_1-1)}{(p_1+1)^{2\ell +1}}\\
        & \ = \ \bigg[\frac{\hat\phi_2(0)}{4\log(R)}+\frac{F_1\hat\phi_2(0)}{\log(R)}+ O\left(\frac{1}{\log^4(R)}\right)\bigg]\cdot\left[\log 2 +\frac{1}{4}+F_4\right]\\
        & \ = \ \frac{\hat\phi_2(0)(\log 2+\frac{1}{4}+F_4)}{\log(R)}\left(\frac{1}{4}+F_1\right)+O\left(\frac{1}{\log^4(R)}\right). 
    \end{aligned}
\end{equation}
\end{proof}

\begin{lemma}
     Let $\hat\phi$ be even Schwartz function. Then
    \begin{equation}
    \begin{aligned}
        & \sum_{p} \frac{\log^2(p)}{\log^2(R)}\hat\phi\left(\frac{2\log(p)}{\log(R)}\right) \sum_{\ell=2}^\infty (C_{\ell+1}-C_\ell) \frac{p^\ell(p-1)}{(p+1)^{2\ell +1}}\\
        & \ = \ \frac{\hat\phi_2(0)}{4\log^2(R)}
        \ - \ \frac{\hat\phi_2(0)}{\log^2(R)}\int_1^\infty E(t)\frac{(-9t^4+10t^2+10t+3)\log(t)+3t^4+3t^3-2t^2-3t-1}{t^4(t+1)^4}dt\\
        & \ + \ O\left(\frac{1}{\log^4(R)}\right).
    \end{aligned}
    \end{equation}
\end{lemma}
\begin{proof}
    We first simplify the sum over $\ell$. Let
    \begin{equation}
        G(x) \ := \ \sum_{\ell=0}^\infty C_\ell x^\ell \ = \ \frac{1-\sqrt{1-4x}}{2x}
    \end{equation}
    and
    \begin{equation}
        S\ :=\ \sum_{\ell=2}^\infty (C_{\ell+1}-C_\ell)z^\ell
    \end{equation}
    where $z=p/(p+1)^2.$ Since $C_0=C_1=1$ and $C_2=2$, it follows that
    \begin{equation}
        S\ =\ \frac{G(z)-1-z-2z^2}{z}-G+1+z=\frac{(1-z)G-1-z^2}{z}.
    \end{equation}
    Plug in $z=p/(p+1)^2$, noting that $\sqrt{1-4z}=(p-1)/(p+1)$ to get that
    \begin{equation}
        G(z)\ =\ \frac{p+1}{p}
    \end{equation}
    and thus
    \begin{equation}
        S=\frac{3p^2+3p+1}{p^2(p+1)^2}.
    \end{equation}
    Multiplying through by what we factored out, we get
    \begin{equation}
        \sum_{\ell=2}^\infty(C_{\ell+1}-C_\ell) \frac{p^\ell(p-1)}{p(p+1)^{2\ell +1}}=\frac{(p-1)(3p^2+3p+1)}{p^3(p+1)^3}.
    \end{equation}
    The rest of the argument is analogous to \eqref{line4}. We define
    \begin{equation}
        g(x)\ :=\ \hat\phi\left(\frac{2\log(x)}{\log(R)}\right)G(x) \log(x),
    \end{equation}
    where
    \begin{equation}
        G(x)\ :=\ \frac{(x-1)(3x^2+3x+1)}{x^3(x+1)^3}.
    \end{equation}
    Hence, we can rewrite the sum we want as
    \begin{equation}
        \frac{1}{\log^2(R)}\sum_p \log(p)\, g(p)
        \ = \ -\frac{1}{\log^2(R)}\int_1^\infty tg'(t)dt -\frac{1}{\log^2(R)}\int_1^\infty E(t)g'(t)dt.
    \end{equation}
    Letting $u=2\log(p)/\log(R)$,
    \begin{equation}
        -\frac{1}{\log^2(R)}\int_1^\infty tg'(t)dt\ =\ \frac{1}{\log^2(R)}\int_1^\infty g(t)dt
        \ = \ \frac{1}{\log^2(R)}\int_1^\infty 
        \log(t)\,\hat\phi_2\left(\frac{2\log(t)}{\log(R)}\right) G(t) dt.
    \end{equation}
    This is equal to
    \begin{equation}
        \frac{\hat\phi_2(0)}{\log^2(R)}\int_1^\infty \frac{\log(t) (t-1)(3t^2+3t+1)}{t^3(t+1)^3}+O\left(\frac{1}{\log^4(R)}\right)\ =\ \frac{\hat\phi_2(0)}{4\log^2(R)}+O\left(\frac{1}{\log^4(R)}\right).
        \label{ll10f}
    \end{equation}
    We now consider the error term. We have that 
    \begin{equation}
        g'(x)\ =\ \frac{A(x)\log(x) \hat\phi_2'\left(\frac{2\log(x)}{\log(R)}\right)+(B(x)\log(x)+C(x))\log(R) \hat\phi_2\left(\frac{2\log(x)}{\log(R)}\right)}{\log(R)x^4(x+1)^4}
    \end{equation}
    where $A(x)$, $B(x)$, and $C(x)$ are polynomials. Expanding $\hat\phi$ and $\hat\phi'$, the main term doesn't contribute, so the error is
    \begin{equation}
    -\frac{\hat\phi_2(0)}{\log^2(R)}\int_1^\infty E(t)\frac{B(t)\log(t)+C(t)}{t^4(t+1)}dt
    \end{equation}
    which is explicitly
    \begin{equation}
        -\frac{\hat\phi_2(0)}{\log^2(R)}\int_1^\infty E(t)\frac{(-9t^4+10t^2+10t+3)\log(t)+3t^4+3t^3-2t^2-3t-1}{t^4(t+1)^4}dt.
        \label{ll10s}
    \end{equation}
    Combining \eqref{ll10f} and \eqref{ll10s} yields the lemma.
\end{proof}

\printbibliography

\end{document}

